% ===========================================================================
% main.tex —— 只是骨架。看這一檔就懂全書結構。
%   建置:  make full   (含 L3 實作補充)
%          make paper  (乾淨學術版,剝掉 L3)
%   或:    xelatex "\def\IMPL{1}% ===========================================================================
% main.tex —— 只是骨架。看這一檔就懂全書結構。
%   建置:  make full   (含 L3 實作補充)
%          make paper  (乾淨學術版,剝掉 L3)
%   或:    xelatex "\def\IMPL{1}% ===========================================================================
% main.tex —— 只是骨架。看這一檔就懂全書結構。
%   建置:  make full   (含 L3 實作補充)
%          make paper  (乾淨學術版,剝掉 L3)
%   或:    xelatex "\def\IMPL{1}% ===========================================================================
% main.tex —— 只是骨架。看這一檔就懂全書結構。
%   建置:  make full   (含 L3 實作補充)
%          make paper  (乾淨學術版,剝掉 L3)
%   或:    xelatex "\def\IMPL{1}\input{main}"
% ===========================================================================
\documentclass[11pt,oneside]{ctexbook}   % 中文 → ctexbook + XeLaTeX
\input{preamble}

\title{Saccade 全局數學模型}
\author{Ray Lei}

\begin{document}
% Keep page numbering format uniform across title, TOC, main text, and appendix.
\maketitle
\setcounter{page}{2}
\tableofcontents

% ---- Part I:主架構(the spine,只放地圖,不放數學)----
\part{系統總覽}
\include{chapters/00-overview}

% ---- Part II:模組細節(一章一模組,全 pipeline 等深 L1–L2)----
% 每章 = \archmap{...} 迷你地圖 + contract 合約框 + 核心數學
%        + 末尾 \suppinput 掛 L3 實作補充(paper 版會自動剝掉)
\part{模組細節}
\include{chapters/04-detection}    % pipeline 最前段
\include{chapters/05-gmc}          % ← 完整範例章
\include{chapters/06-kalman}       % track 主體 1/3
\include{chapters/07-association}  % track 主體 2/3
\include{chapters/08-auction}      % track 主體 3/3
\include{chapters/09-lifecycle}    % track 收尾 1/2
\include{chapters/10-relink}       % track 收尾 2/2

% ---- 附錄:reference 性質的大表(符號↔config)放這,不擋正文節奏 ----
\appendix
\include{chapters/A-symbols}

\end{document}
"
% ===========================================================================
\documentclass[11pt,oneside]{ctexbook}   % 中文 → ctexbook + XeLaTeX
% ===========================================================================
% preamble.tex —— 套件、深度開關、archmap 迷你地圖、合約框
% 引擎:XeLaTeX(中文必須)。建置見 README.md / Makefile。
% ===========================================================================

% --- 繁體中文字型(fandol 為簡體取向,缺繁體字形;改用 Noto CJK TC)---
\setCJKmainfont{Noto Serif CJK TC}
\setCJKsansfont{Noto Sans CJK TC}
\setCJKmonofont{Noto Sans CJK TC}
% ctexbook 預設 caption 是簡體;改回繁體(录→錄)。appendixname 同時管標題與頁首。
\ctexset{contentsname=目錄, appendixname=附錄}
% 中英混排時,允許行末額外伸縮以吸收邊際 overfull(標準作法)
\setlength{\emergencystretch}{3em}

% --- 數學 / 表格 / 引用 ---
\usepackage{amsmath, amssymb, mathtools}
\usepackage{booktabs}            % \toprule \midrule \bottomrule,禁直線
\usepackage{array}
\usepackage{tabularx}            % 自動寬度欄(職責表用)
\usepackage{longtable}           % 可跨頁的參考表(附錄符號表用)
\usepackage[hidelinks]{hyperref}
\usepackage{cleveref}            % \Cref{eq:...} 就近引用
\usepackage{xstring}             % archmap 的字串比較(跨引擎安全)
\usepackage{fancyhdr}            % 統一頁碼位置

% --- 圖 ---
\usepackage{tikz}
\usetikzlibrary{arrows.meta, positioning, fit, backgrounds}

% --- 程式碼 / 合約框 ---
\usepackage[most]{tcolorbox}
\usepackage{fvextra}             % Verbatim,ASCII pipeline 圖暫存用

% --- 頁碼 ---
% ctexbook/book 的一般頁與章首頁預設 page style 不同,會讓頁碼在頁首/頁腳間跳。
% 這裡同時覆寫一般頁型與 plain 頁型,固定頁碼在頁腳中央。
\fancypagestyle{saccadepage}{
  \fancyhf{}
  \fancyfoot[C]{\thepage}
  \renewcommand{\headrulewidth}{0pt}
  \renewcommand{\footrulewidth}{0pt}
}
\fancypagestyle{plain}{
  \fancyhf{}
  \fancyfoot[C]{\thepage}
  \renewcommand{\headrulewidth}{0pt}
  \renewcommand{\footrulewidth}{0pt}
}
\AtBeginDocument{\pagestyle{saccadepage}}

% ===========================================================================
% 深度開關:L3 實作補充是否編入
%   make full  -> \def\IMPL{1}  含 supp/*-impl.tex
%   make paper -> \def\IMPL{0}  剝掉所有 L3,得乾淨學術版
% ===========================================================================
\providecommand{\IMPL}{1}
\newif\ifimpl
\ifnum\IMPL=1 \impltrue \else \implfalse \fi

% \suppinput{file}:只有 full 版會把 L3 補充檔讀進來
\newcommand{\suppinput}[1]{\ifimpl\input{#1}\fi}

% L3 補充章節的統一外框(視覺上和學術本體區隔)
\newtcolorbox{implsupp}[1]{
  enhanced, breakable,
  colback=orange!4, colframe=orange!55!black,
  fonttitle=\bfseries\small, coltitle=orange!30!black,
  title={實作補充 (L3,僅 full 版)\quad#1},
  left=2mm, right=2mm, top=1mm, bottom=1mm}

% ===========================================================================
% 模組合約框:每個細節章開頭,濃縮 I/O / Source / Baseline
% ===========================================================================
\newtcolorbox{contract}[1]{
  enhanced, breakable,
  colback=gray!4, colframe=black!55,
  fonttitle=\bfseries, title={模組合約:#1}}

% ===========================================================================
% \archmap{focus} —— 全書共用的「你在這裡」迷你地圖
%   focus ∈ {none, detect, gmc, track, output}
%   none = 全亮(總覽章用);其餘 = 點亮焦點、淡化其它
%   track 內的子章(Kalman/association/auction/...)先一律用 track;
%   日後要更細,可在此複製出 \trackmap 畫 track 內部子階。
% ===========================================================================
% 每個 stage 的代表色(改這裡就改全書配色)——用 \colorlet 建具名顏色,
% 才能在 pgfkeys option 裡以字面名稱引用(巨集名不會被展開)。
\colorlet{Cbuf}{teal}
\colorlet{Cdet}{blue!65!cyan}
\colorlet{Cgmc}{orange!90!red}
\colorlet{Ctrk}{violet!80!blue}
\colorlet{Cout}{green!55!black}

% 全域 style:box 外型 + 三種狀態(吃顏色當 #1)+ 箭頭/標籤
\tikzset{
  ambox/.style  ={draw, rounded corners, align=center,
                  minimum height=10mm, minimum width=24mm, font=\small},
  amfull/.style ={fill=#1!25, draw=#1!85!black, line width=1pt,   text=black!88},
  amfade/.style ={fill=#1!8,  draw=#1!35,       line width=0.6pt, text=black!40},
  amfocus/.style={fill=#1!48, draw=#1!90!black, line width=1.9pt, text=black!92},
  amarr/.style  ={-{Stealth}, thick},               % 直角:不加 rounded corners
  amlbl/.style  ={font=\scriptsize, text=black!55},
}

\newcommand{\archmap}[1]{%
  % 預設 full(總覽 none 全部上色);具名 style 都是字面,pgfkeys 吃得到
  \tikzset{amSbuf/.style={amfull=Cbuf}, amSdet/.style={amfull=Cdet},
           amSgmc/.style={amfull=Cgmc}, amStrk/.style={amfull=Ctrk},
           amSout/.style={amfull=Cout}}%
  \IfStrEq{#1}{none}{}{% 指定焦點:先全 fade,再把焦點改 focus
    \tikzset{amSbuf/.style={amfade=Cbuf}, amSdet/.style={amfade=Cdet},
             amSgmc/.style={amfade=Cgmc}, amStrk/.style={amfade=Ctrk},
             amSout/.style={amfade=Cout}}%
    \IfStrEq{#1}{detect}{\tikzset{amSdet/.style={amfocus=Cdet}}}{}%
    \IfStrEq{#1}{gmc}{\tikzset{amSgmc/.style={amfocus=Cgmc}}}{}%
    \IfStrEq{#1}{track}{\tikzset{amStrk/.style={amfocus=Ctrk}}}{}%
    \IfStrEq{#1}{output}{\tikzset{amSout/.style={amfocus=Cout}}}{}%
  }%
  % 固定成欄寬比例再置中,避免視覺偏移
  \resizebox{0.92\linewidth}{!}{%
  \begin{tikzpicture}
    % 橫式佈局:左→右流,detect/GMC 為上下兩條平行支線
    \node[ambox, amSbuf] (buf) at (0,0)     {frame buffer\\\scriptsize(GPU, CHW)};
    \node[ambox, amSdet] (det) at (3.2, 1)  {detect 鏈};
    \node[ambox, amSgmc] (gmc) at (3.2,-1)  {GMC};
    \node[ambox, amStrk] (trk) at (6.6, 0)  {track\\\scriptsize tracker\_gpu.cu};
    \node[ambox, amSout] (out) at (9.8, 0)  {materialize\\\scriptsize→ MOT rows};
    % 分岔:buf 拉出 bus,再直角接上下兩支線(連線染來源色)
    \draw[amarr, Cbuf!65!black] (buf.east) -- ++(0.4,0) |- (det.west);
    \draw[amarr, Cbuf!65!black] (buf.east) -- ++(0.4,0) |- (gmc.west);
    % 匯合:兩支線直角收進 track
    \draw[amarr, Cdet!65!black] (det.east) -- ++(0.4,0) node[amlbl,above]{boxes} |- (trk.west);
    \draw[amarr, Cgmc!70!black] (gmc.east) -- ++(0.4,0) node[amlbl,below]{$W_f$} |- (trk.west);
    \draw[amarr, Ctrk!65!black] (trk.east) -- (out.west);
  \end{tikzpicture}}}

% ===========================================================================
% 章內「演算法傳遞圖」共用 style —— 每章自畫,顏色用該模組的 C*
%   step=<color>     一個演算法步驟方塊
%   flowarr=<color>  傳遞箭頭(直角時用 |- / -|)
%   flowlbl          標在邊上的「傳了什麼」(白底壓線)
% 用法見 chapters/05-gmc.tex 的範例。
% ===========================================================================
\tikzset{
  step/.style   ={draw, rounded corners=2pt, align=center, font=\footnotesize,
                  minimum height=10mm, minimum width=20mm, inner sep=4pt,
                  fill=#1!14, draw=#1!70!black, text=black!88},
  flowarr/.style={-{Stealth}, semithick, draw=#1!75!black},
  flowlbl/.style={font=\scriptsize, text=black!60, fill=white, inner sep=1.5pt},
}
% 把 GMC 章內流程一行寫完的便利指令(其他章可仿照自行展開)
\newcommand{\gmcflow}{%
  \begin{center}\resizebox{\linewidth}{!}{%
  \begin{tikzpicture}[node distance=13mm]
    \node[step=Cgmc] (n1) {灰階降採樣\\+\,Hanning};
    \node[step=Cgmc, right=of n1] (n2) {FFT\\prev / curr};
    \node[step=Cgmc, right=of n2] (n3) {cross-power\\$R=\overline{A}B/|{\cdot}|$};
    \node[step=Cgmc, right=of n3] (n4) {IFFT $\to$ peak\\PCR};
    \node[step=Cgmc, right=of n4] (n5) {confidence $\gamma$\\$\to$ warp $W_f$};
    \draw[flowarr=Cgmc] (n1) -- node[flowlbl,above]{$G_w$}        (n2);
    \draw[flowarr=Cgmc] (n2) -- node[flowlbl,above]{$A,B$}        (n3);
    \draw[flowarr=Cgmc] (n3) -- node[flowlbl,above]{$R(k)$}       (n4);
    \draw[flowarr=Cgmc] (n4) -- node[flowlbl,above]{$r$ peak}     (n5);
  \end{tikzpicture}}\end{center}}


\title{Saccade 全局數學模型}
\author{Ray Lei}

\begin{document}
% Keep page numbering format uniform across title, TOC, main text, and appendix.
\maketitle
\setcounter{page}{2}
\tableofcontents

% ---- Part I:主架構(the spine,只放地圖,不放數學)----
\part{系統總覽}
% ===========================================================================
% Part I 唯一一章:系統總覽 = 主架構(the spine)
% 原則:這裡只放「地圖」——架構圖、模組邊界、dataflow、baseline 合約。
%       所有數學一律下放到 Part II。這就是「不全展開」的關鍵。
% ===========================================================================
\chapter{系統總覽}\label{ch:overview}

整條 pipeline 分兩層:\textbf{Python orchestration}(排程、buffer 管理、
CUDA graph replay、eval)與 \textbf{C++/CUDA compute}(detector、postprocess、
GMC、tracker association)。Python 不在 hot path 做數值計算,只在 output
boundary 把結果搬回 host。

關鍵:detect 與 GMC 是從\textbf{同一個 frame buffer 分岔的兩條平行支線}——
GMC 吃 frame 灰階(+上一幀),\emph{不是} detect 的輸出;\texttt{track} 才把
兩條支線匯合。

\section{架構地圖}\label{sec:archmap}
全書各章開頭都會重畫這張圖、點亮當前模組(「你在這裡」)。總覽版全亮:

\begin{center}
  \archmap{none}
\end{center}

\section{模組職責與 I/O}\label{sec:modules}
下表概括每個模組\textbf{吃什麼、吐什麼、解決什麼問題},以及在 pipeline 中銜接誰。
逐項數學見 Part II 對應章;傳遞合約與 source 見 \Cref{tab:transfer}。

\begin{table}[h]\small\setlength{\tabcolsep}{4pt}\renewcommand{\arraystretch}{1.35}
\begin{tabularx}{\linewidth}{@{}l >{\raggedright\arraybackslash}X >{\raggedright\arraybackslash}X l@{}}
\toprule
模組 & 輸入 → 輸出 & 解決的問題 & 銜接 \\
\midrule
detect      & frame CHW tensor → raw boxes/scores/classes
            & 從單張影像偵測行人(TRT backbone + Mamba head,anchor-free decode)
            & → postprocess \\
postprocess & raw det → 過濾/NMS 後 boxes
            & 去除低分與重疊框,交給 tracker 乾淨的偵測集
            & detect → track \\
GMC         & prev+curr 灰階 → \(2\times3\) warp \(W_f\)
            & 估相機運動補償平移,讓 Kalman 的 velocity 只學物體運動
            & 與 detect 平行,補 track \\
track       & boxes/scores/\(W_f\) → track states / IDs
            & 跨幀關聯與維持 ID(分解見 \Cref{tab:track-internal})
            & 匯合 → materialize \\
output      & GPU state → MOT rows
            & 搬回 host、產生 MOT 結果(含 tracklet 內插後處理)
            & track → 終點 \\
\bottomrule
\end{tabularx}
\caption{模組職責與輸入/輸出概覽。}\label{tab:modules}
\end{table}

\texttt{track} 不是單一步驟,而由數個子機制串成,各為 Part II 的一章:

\begin{table}[h]\small\setlength{\tabcolsep}{4pt}\renewcommand{\arraystretch}{1.3}
\begin{tabularx}{\linewidth}{@{}l >{\raggedright\arraybackslash}X@{}}
\toprule
子機制 & 解決的問題 \\
\midrule
Kalman 模型      & constant-velocity 運動模型;每幀 predict 再用 detection update track state \\
Association cost & 把 IoU、OAO/velocity penalty、stability reward 匯成單一配對成本 \(c_{ij}\) \\
Auction 指派     & ByteTrack 風格分數級聯 + 單輪平行 auction,決定 track 與 detection 配對 \\
Track lifecycle  & track 生死:tentative→confirmed 的 birth/confirm,與 lost/remove \\
Bridge relink    & 找回短暫消失又重生的 ID(雙向中點外推,非 appearance ReID) \\
\bottomrule
\end{tabularx}
\caption{\texttt{track} 內部子機制(Part II 逐章展開)。}\label{tab:track-internal}
\end{table}

傳遞合約(誰留在 GPU、source 檔)另列:

\begin{table}[h]\small\setlength{\tabcolsep}{4pt}\renewcommand{\arraystretch}{1.2}
\begin{tabularx}{\linewidth}{@{}l
  >{\raggedright\arraybackslash}X
  >{\raggedright\arraybackslash}X@{}}
\toprule
模組 & 模組(Python facade / C++·CUDA core) & 傳遞方式 \\
\midrule
detect      & \texttt{mamba\_gated\_detector.py} / TRT backbone + Mamba head & GPU tensor(whole-detect graph)\\
postprocess & \texttt{PerceptionPipeline} / \texttt{pipeline.cpp} & GPU buffer,留 device \\
GMC         & \texttt{GMC} / \texttt{gmc\_kernel.cu} (cuFFT)       & GPU tensor,$2\times3$ warp \\
track       & \texttt{GPUByteTracker} / \texttt{tracker\_gpu.cu}   & device pointer + stream \\
materialize & \texttt{materialize\_\allowbreak gpu\_\allowbreak track\_\allowbreak results} & GPU→host,僅一次 copy \\
output      & \texttt{relink\_write}                               & host(fast MOT emit)\\
\bottomrule
\end{tabularx}
\caption{傳遞合約與 source。}\label{tab:transfer}
\end{table}

\section{Frame-Level Dataflow}\label{sec:dataflow}
對每個 frame \(f\),no-ReID baseline 的主線是兩條同源分支在 tracker 匯合:
\begin{center}
\small
\begin{tabular}{@{}rcl@{}}
frame tensor & \(\to\) & detect \(\to\) postprocess \(\to\) boxes \\
frame tensor & \(\to\) & GMC warp \(W_f\) \\
boxes \(+\ W_f\) & \(\to\) & tracker update \(\to\) materialize \\
materialize & \(\to\) & relink\_write / MOT rows
\end{tabular}
\end{center}

\section{Baseline 合約}\label{sec:baseline}
本文以 \texttt{mamba\_whole\_graph} preset 為基準。關鍵啟用條件:

\begin{table}[h]\small\setlength{\tabcolsep}{4pt}
\begin{tabularx}{\linewidth}{@{}l
  >{\raggedright\arraybackslash}p{4.8cm}
  >{\raggedright\arraybackslash}X@{}}
\toprule
區域 & baseline 值 & 效果 \\
\midrule
Graph     & \texttt{use\_whole\_graph: true}      & whole detect graph 單次 replay \\
GMC       & \texttt{gmc\_downscale: 4}            & GPU phase-correlation 估平移 \\
Appearance& \texttt{reid\_mode: off}              & headline 不用 appearance \\
Assoc     & \texttt{match\_thresh: 0.50}          & ByteTrack-like 多階段 gate \\
Cost      & \texttt{multiplicative\_cost: true}   & log-linear cost + stability reward \\
Relink    & \texttt{relink\_bridge\_enabled: true}& 雙向 foot bridge \\
\bottomrule
\end{tabularx}
\caption{Baseline 合約摘要。完整旋鈕表見附錄 A。}\label{tab:baseline}
\end{table}

% TODO: 把 math_model.md §3 的符號↔config 大表搬進 chapters/A-symbols.tex(附錄)


% ---- Part II:模組細節(一章一模組,全 pipeline 等深 L1–L2)----
% 每章 = \archmap{...} 迷你地圖 + contract 合約框 + 核心數學
%        + 末尾 \suppinput 掛 L3 實作補充(paper 版會自動剝掉)
\part{模組細節}
% ===========================================================================
% Detection 與 Postprocess(pipeline 最前段,故置於 Part II 第一章)
% ===========================================================================
\chapter{Detection 與 Postprocess}\label{ch:detect}

\begin{center}\archmap{detect}\end{center}

\begin{contract}{detect 鏈}
\textbf{輸入}\; frame \(640\times640\) CHW tensor\quad
\textbf{輸出}\; boxes / scores / classes(post-NMS)\\
\textbf{方法}\; YOLO26s TRT backbone+neck \(\to\) Mamba head(per-level lane)\(\to\)
anchor-free decode \(\to\) NMS\quad
\textbf{Source}\; \texttt{mamba\_gated\_detector.py}、\texttt{mamba\_head.py}、
\texttt{mamba\_gated\_detector.cpp}、\texttt{pipeline.cpp}\\
\textbf{Baseline}\; \texttt{tiling: native\_640}、\texttt{preprocess: none}、
ckpt \texttt{mamba\_gt\_v14replica\_t3\_t1};整段在 whole-detect CUDA graph 單次 replay
\end{contract}

\paragraph{演算法傳遞.} detect 鏈分兩段:detector(backbone+head,出 raw)與 native
postprocess(過濾/NMS)。邊上為各段傳遞的張量:

\begin{center}\resizebox{\linewidth}{!}{%
\begin{tikzpicture}[node distance=14mm]
  \node[step=Cdet] (n1) {frame\\$640^2$};
  \node[step=Cdet, right=of n1] (n2) {TRT backbone\\$+$ neck (FPN)};
  \node[step=Cdet, right=of n2] (n3) {Mamba head\\$\times(P_3,P_4,P_5)$};
  \node[step=Cdet, right=of n3] (n4) {decode\\anchor-free};
  \node[step=Cdet, right=of n4] (n5) {filter $+$ NMS};
  \draw[flowarr=Cdet] (n1) -- node[flowlbl,above]{image}              (n2);
  \draw[flowarr=Cdet] (n2) -- node[flowlbl,above]{$F_3,F_4,F_5$}       (n3);
  \draw[flowarr=Cdet] (n3) -- node[flowlbl,above]{$\mathrm{cls}_N,\mathrm{reg}_N$} (n4);
  \draw[flowarr=Cdet] (n4) -- node[flowlbl,above]{8400 anchors}        (n5);
\end{tikzpicture}}\end{center}

\section{整體結構與三尺度}\label{sec:det-overall}
P3/P4/P5 三條 lane 結構完全相同,只差尺度常數(輸入固定 640):

\begin{table}[h]\centering\small\renewcommand{\arraystretch}{1.2}
\begin{tabular}{@{}lcccc@{}}
\toprule
level \(N\) & stride \(s_N\) & \(H_N{=}W_N{=}640/s_N\) & channels \(C_N\) & tokens \(L_N{=}(H_N/4)^2\) \\
\midrule
P3 & 8  & 80 & 128 & 400 \\
P4 & 16 & 40 & 256 & 100 \\
P5 & 32 & 20 & 512 & 25  \\
\bottomrule
\end{tabular}
\caption{三尺度常數(\(s_N{=}2^N,\ C_N{=}2^{N+4}\))。}\label{tab:det-levels}
\end{table}

\section{單一 Lane:Mamba Head(Flow B)}\label{sec:det-mamba}
重點是 U-Net 式 skip:\(X_N\) 一邊進 down\(\to\)Mamba\(\to\)up,一邊\textbf{跳接}直接與
\(U_N\) concat,故 head 輸入是 256 ch。先 \texttt{Down4} 把序列壓到 \(L_N\le400\)
(控 scan 成本),上採樣後再補回細節:

\begin{center}\resizebox{\linewidth}{!}{%
\begin{tikzpicture}[node distance=11mm]
  \node[step=Cdet] (f)  {$F_N$\\$C_N{\times}H_N{\times}W_N$};
  \node[step=Cdet, right=of f]  (x)  {input\_proj\\$X_N$ (128)};
  \node[step=Cdet, right=of x]  (d)  {Down4\\stride-4};
  \node[step=Cdet, right=of d]  (z)  {flatten\\$Z_N$ ($L_N{\times}128$)};
  \node[step=Cdet, right=of z]  (m)  {MambaBlock\\$\times2$ ($d_{\mathrm{state}}{=}16$)};
  \node[step=Cdet, right=of m]  (u)  {Up4\\$U_N$ (128)};
  \node[step=Cdet, right=of u]  (c)  {concat\\$\to Y_N$ (256)};
  \draw[flowarr=Cdet] (f) -- (x);
  \draw[flowarr=Cdet] (x) -- (d);
  \draw[flowarr=Cdet] (d) -- node[flowlbl,above]{$\le400$ tok} (z);
  \draw[flowarr=Cdet] (z) -- (m);
  \draw[flowarr=Cdet] (m) -- (u);
  \draw[flowarr=Cdet] (u) -- (c);
  % U-Net skip:X_N 直接跳接到 concat(直角繞上方)
  \draw[flowarr=Cdet, dashed] (x.north) |- ++(0,7mm) -| node[flowlbl,above,pos=0.25]{skip $X_N$} (c.north);
\end{tikzpicture}}\end{center}

\section{Head 與 Decode(Flow C/D)}\label{sec:det-decode}
\textbf{Head}(per level,純卷積):\(Y_N(256)\) 分兩支,cls head
\(\to\mathrm{cls}_N\,(80)\)、reg head \(\to\mathrm{reg}_N\,(4)\);輸出通道
\(n_o=n_c+4\,r_{\max}=80+4=84\)(\(r_{\max}{=}1\),box 直接 4 維距離,無 DFL)。

\textbf{Decode}(anchor-free,跨三尺度合併,\(N=80^2{+}40^2{+}20^2=8400\) anchors)。
anchor 為 grid 中心 \((x{+}0.5,y{+}0.5)\),每點帶 stride:
\begin{equation}\label{eq:det-decode}
  \mathrm{reg}=(\ell t,rb)\ \xrightarrow{\text{dist2bbox}}\
  c=\mathrm{anchor}+\tfrac{rb-\ell t}{2},\quad
  wh=\ell t+rb\ \xrightarrow{\times s}\ \text{xyxy(像素)}.
\end{equation}
\begin{equation}\label{eq:det-cls}
  \text{score},\ \text{class}=\max_{80},\ \arg\max_{80}\ \sigma(\mathrm{cls}).
\end{equation}

\section{Postprocess Contract}\label{sec:det-post}
native postprocess(\Cref{ch:detect} 第二段)的核心 contract:
\begin{center}\ttfamily\small
(raw boxes/scores/classes) → score/class/geometry filter → compact/gather → NMS → output
\end{center}
baseline 旗標全關:
\begin{center}\small
\begin{tabular}{@{}ll@{\qquad}ll@{}}
\texttt{track\_person\_only} & false &
\texttt{person\_geometry\_prior} & false \\
\texttt{detection\_quality\_scaling} & false &
\texttt{geometry\_suspect\_support} & false
\end{tabular}
\end{center}
tracker 直接消費 post-NMS boxes,各 detection 進哪段 association 由 tracker
的 score gates(\(\tau_{\mathrm{track/mid/high}}\)、\texttt{new\_track\_thresh})決定。

\suppinput{supp/04-detection-impl}
    % pipeline 最前段
% ===========================================================================
% 細節章「樣板」:GMC。後續每章都照這個骨架寫,確保全 pipeline 等深。
%   1) \archmap{...}  迷你地圖,點亮本模組
%   2) contract 框    I/O / Source / Baseline,讓人不鑽公式也拿得到全貌
%   3) L1–L2 數學     學術本體(收 paper 用這層)
%   4) \suppinput     掛 L3 實作補充(paper 版自動剝掉)
% ===========================================================================
\chapter{GMC:相機運動補償}\label{ch:gmc}

\begin{center}\archmap{gmc}\end{center}

\begin{contract}{GMC}
\textbf{輸入}\; prev + curr 灰階(GPU tensor, CHW float \([0,1]\))\quad
\textbf{輸出}\; \(2\times3\) translation warp \(W_f\)\\
\textbf{方法}\; phase correlation(cross-power spectrum + Hanning window)\quad
\textbf{Source}\; \texttt{gmc\_kernel.cu}(cuFFT)\\
\textbf{Baseline}\; \texttt{gmc\_downscale: 4}、translation-only、soft confidence gate
\end{contract}

\noindent
GPU GMC 用 phase correlation 估 frame-to-frame 平移,作為 deterministic control
input 送進 \texttt{track} 的 Kalman predict,讓 velocity state 只學物體
residual motion 而非相機運動。

\paragraph{演算法傳遞.} 本章內部依序串接下列步驟,邊上標註各步之間傳遞的量
(對應 \Cref{sec:gmc-down,sec:gmc-cps,sec:gmc-warp}):

\gmcflow

\section{灰階降採樣與窗函數}\label{sec:gmc-down}
對 downscaled pixel \((x,y)\),floor-based 取樣:
\begin{equation}\label{eq:gmc-sample}
  s_x=\Big\lfloor x\tfrac{W_{\mathrm{src}}}{W_{\mathrm{dst}}}\Big\rfloor,\qquad
  s_y=\Big\lfloor y\tfrac{H_{\mathrm{src}}}{H_{\mathrm{dst}}}\Big\rfloor,\qquad
  G=0.299R+0.587G_c+0.114B.
\end{equation}
FFT 前套 Hanning window \(G_w(x,y)=G(x,y)\,w_x w_y\),其中
\(w_x=\tfrac12(1-\cos\tfrac{2\pi x}{W_{\mathrm{dst}}-1})\)(同理 \(w_y\))。

\section{Cross-Power Spectrum 與位移}\label{sec:gmc-cps}
令 \(A=\mathrm{FFT}(\text{prev})\)、\(B=\mathrm{FFT}(\text{curr})\):
\begin{equation}\label{eq:gmc-cps}
  C(k)=\overline{A(k)}B(k),\qquad
  R(k)=\frac{C(k)}{|C(k)|+10^{-6}},\qquad
  r=\mathrm{IFFT}(R).
\end{equation}
\(r\) 的 peak 給出 wrapped displacement(\(\text{peak}>\text{dim}/2\) 時減去該維長度)。

\section{Confidence 與 Warp}\label{sec:gmc-warp}
phase-correlation 可信度與軟縮放:
\begin{equation}\label{eq:gmc-pcr}
  \mathrm{PCR}=\frac{\max r}{\mathrm{RMS}(r)},\qquad
  \gamma=\begin{cases}\mathrm{clamp}(\mathrm{PCR}/\tau,0,1),&\mathrm{PCR}<\tau\\1,&\text{otherwise.}\end{cases}
\end{equation}
位移超過 downscaled 尺寸 25\% 視為不可信(退回 identity)。否則
\begin{equation}\label{eq:gmc-final}
  t_x=p_x d\,\gamma,\quad t_y=p_y d\,\gamma,\qquad
  W_f=\begin{bmatrix}1&0&t_x\\0&1&t_y\end{bmatrix}.
\end{equation}
\textbf{Sign convention}:positive \(t_x\) 把 predicted track center 往
current frame 右方推。

% ---------------------------------------------------------------------------
% L3 實作補充:kernel 名、env、四路徑 fallback 差異。paper 版自動剝掉。
% ---------------------------------------------------------------------------
\suppinput{supp/05-gmc-impl}
          % ← 完整範例章
% ===========================================================================
% track 主體 (1/3):Kalman 運動模型
% ===========================================================================
\chapter{Kalman 運動模型}\label{ch:kalman}

\begin{center}\archmap{track}\end{center}

\begin{contract}{Kalman(track 子機制)}
\textbf{輸入}\; 前一幀 track state \((x,P)\)、GMC warp \(W_f\)、配對到的 detection
measurement \(z\)\quad
\textbf{輸出}\; predict 後 \(\tilde{x}\)(供 gate/cost)、update 後 \((x,P)\)\\
\textbf{方法}\; SORT/DeepSORT 風格 constant-velocity filter,state
\((c_x,c_y,a,h,\dot{\cdot})\)\quad
\textbf{Source}\; \texttt{kalman\_gpu.cuh}\quad
\textbf{Baseline}\; \texttt{kalman\_r\_scale: 2.8}、NSA off
\end{contract}

\noindent
state 與 measurement:
\begin{equation}\label{eq:kf-state}
  x=(c_x,c_y,a,h,v_x,v_y,v_a,v_h)^{\!\top},\qquad
  z=(c_x,c_y,a,h)^{\!\top}.
\end{equation}

\paragraph{演算法傳遞.} 每幀對每個 track 依序:predict(先吃 GMC control)→ 算
innovation 協方差 → gate → 若入選則 update。邊上為各步傳遞的量:

\begin{center}\resizebox{\linewidth}{!}{%
\begin{tikzpicture}[node distance=14mm]
  \node[step=Ctrk] (n1) {GMC control\\$+$ predict\\$x'{=}Fx,\ P'{=}FPF^{\!\top}{+}Q$};
  \node[step=Ctrk, right=of n1] (n2) {innovation cov\\$S{=}MP'M^{\!\top}{+}R$};
  \node[step=Ctrk, right=of n2] (n3) {gate\\$\mathrm{IoU}\,\lor\,d^2_{\mathrm{maha}}{<}\tau$};
  \node[step=Ctrk, right=of n3] (n4) {gain\\$K{=}P'M^{\!\top}S^{-1}$};
  \node[step=Ctrk, right=of n4] (n5) {update\\$x{+}Ky,\ (I{-}KM)P'$};
  \draw[flowarr=Ctrk] (n1) -- node[flowlbl,above]{$\tilde{x},P'$} (n2);
  \draw[flowarr=Ctrk] (n2) -- node[flowlbl,above]{$S^{-1}$}       (n3);
  \draw[flowarr=Ctrk] (n3) -- node[flowlbl,above]{候選 $z_j$}     (n4);
  \draw[flowarr=Ctrk] (n4) -- node[flowlbl,above]{$y{=}z{-}Mx$}   (n5);
\end{tikzpicture}}\end{center}

\section{Prediction}\label{sec:kf-predict}
Transition 是 constant velocity,且 GMC translation 先以 control input 加到位置
(見 \Cref{ch:gmc}),再做 predict:
\begin{equation}\label{eq:kf-predict}
  F=\begin{bmatrix}I_4 & I_4\\ 0 & I_4\end{bmatrix},\qquad
  x'=Fx,\qquad P'=FPF^{\!\top}+Q(h^-).
\end{equation}
process noise 隨 predict 後框高 \(h^-\) 縮放:
\begin{equation}\label{eq:kf-Q}
  \sigma_p=\tfrac{h^-}{20},\quad \sigma_v=\tfrac{h^-}{160},\quad
  \mathrm{diag}(Q)=(\sigma_p^2,\sigma_p^2,10^{-4},\sigma_p^2,
                    \sigma_v^2,\sigma_v^2,10^{-10},\sigma_v^2).
\end{equation}

\section{Measurement Update}\label{sec:kf-update}
measurement matrix \(M=[\,I_4\ \ 0\,]\)。noise 用 predict 後 \(h^-\)(非 detection
height);baseline 下亮度/NSA 調節關閉(\(\lambda_{\mathrm{light}}{=}0,\ m_{\mathrm{NSA}}{=}1\)),
故 \(m_R=r_{\mathrm{scale}}\):
\begin{equation}\label{eq:kf-R}
  \mathrm{diag}(R)=\Big(\big(\tfrac{h^-}{20}\big)^2 m_R,\ \big(\tfrac{h^-}{20}\big)^2 m_R,\
                      10^{-2}m_R,\ \big(\tfrac{h^-}{20}\big)^2 m_R\Big).
\end{equation}
更新方程:
\begin{equation}\label{eq:kf-kalman}
\begin{aligned}
  S&=MP'M^{\!\top}+R, & K&=P'M^{\!\top}S^{-1}, & y&=z-Mx',\\
  x&\leftarrow x'+Ky, & P&\leftarrow(I-KM)P'.
\end{aligned}
\end{equation}

\section{Mahalanobis Gate}\label{sec:kf-gate}
即使 IoU 弱,只要 Mahalanobis 距離小於 \(\tau_{\mathrm{maha}}\) 仍可入選 candidate
(\(\tilde{x}_i\) 為已套 GMC control 並 predict 後的 state):
\begin{equation}\label{eq:kf-maha}
  d^2_{ij}=(z_j-M\tilde{x}_i)^{\!\top}S_i^{-1}(z_j-M\tilde{x}_i),\qquad
  \mathrm{cand}_{ij}\iff \mathrm{IoU}_{ij}>\tau_{\mathrm{iou}}\ \lor\ d^2_{ij}<\tau_{\mathrm{maha}}.
\end{equation}
IoU gate 與 cost(\Cref{ch:assoc})使用同一個 \(\tilde{x}_i\)。

\suppinput{supp/06-kalman-impl}
       % track 主體 1/3
% ===========================================================================
% track 主體 (2/3):Association Cost
% ===========================================================================
\chapter{Association Cost}\label{ch:assoc}

\begin{center}\archmap{track}\end{center}

\begin{contract}{Association cost(track 子機制)}
\textbf{輸入}\; predict 後 tracks \(\{\tilde{x}_i\}\)、detections \(\{b_j,s_j\}\)\quad
\textbf{輸出}\; 稀疏 candidate 上的成本 \(c_{ij}\)(供 auction)\\
\textbf{方法}\; IoU(+ReID)綜合質量 → 乘法式 cost,含 OAO/stability 調節\quad
\textbf{Source}\; \texttt{stage1\_cost\_fused\_kernel}(\texttt{tracker\_gpu.cu})\\
\textbf{Baseline}\; \texttt{multiplicative\_cost: true}、\texttt{sinkhorn\_lambda: 10}、
\texttt{stability\_cost\_w: 0.20}、\texttt{reid\_mode: off}
\end{contract}

\paragraph{演算法傳遞.} gate 過的 pair 先算 base quality,再匯入 penalty,經乘法式
轉成 cost,最後只把夠好的 enqueue 進稀疏 candidate list:

\begin{center}\resizebox{\linewidth}{!}{%
\begin{tikzpicture}[node distance=15mm]
  \node[step=Ctrk] (n1) {candidate gate\\$\mathrm{IoU}\,\lor\,\mathrm{Maha}$};
  \node[step=Ctrk, right=of n1] (n2) {base quality\\$A_{ij}{=}q^{\mathrm{iou}}_{ij}$};
  \node[step=Ctrk, right=of n2] (n3) {penalties\\$\Pi{=}P_{\mathrm{OAO}}{+}P_{\mathrm{vel}}{+}P_{\mathrm{occ}}{-}R_{\mathrm{stab}}$};
  \node[step=Ctrk, right=of n3] (n4) {mult.\ cost\\$c{=}1{-}A\,e^{-\Pi}$};
  \node[step=Ctrk, right=of n4] (n5) {sparse top-k\\enqueue};
  \draw[flowarr=Ctrk] (n1) -- node[flowlbl,above]{候選}        (n2);
  \draw[flowarr=Ctrk] (n2) -- node[flowlbl,above]{$A_{ij}$}     (n3);
  \draw[flowarr=Ctrk] (n3) -- node[flowlbl,above]{$\Pi_{ij}$}   (n4);
  \draw[flowarr=Ctrk] (n4) -- node[flowlbl,above]{$c_{ij}$}     (n5);
\end{tikzpicture}}\end{center}

\section{Candidate Gate 與 Base Quality}\label{sec:as-gate}
gate 同 \Cref{eq:kf-maha}。兩 gate 皆不過時 \(c_{ij}=1\)。detection score fusion:
\begin{equation}\label{eq:as-q}
  q^{\mathrm{iou}}_{ij}=\mathrm{IoU}_{ij}\bigl(1-w_{\mathrm{fuse}}(1-s_j)\bigr),
  \qquad A_{ij}=q^{\mathrm{iou}}_{ij}\ \ (\text{no ReID}).
\end{equation}
baseline \texttt{fuse\_score\_weight}=0,故 \(q^{\mathrm{iou}}_{ij}=\mathrm{IoU}_{ij}\)。
sparse list 只 enqueue 成本落在最寬鬆門檻內者:
\(c_{\max}=\max(c_{\mathrm{DDA}},\tau_{\mathrm{match}},\tau_{\mathrm{stage2}})\),
\(\mathrm{enqueue}_{ij}\iff c_{ij}\le c_{\max}\)。

\section{Multiplicative Cost}\label{sec:as-cost}
active cost form 與 penalty 匯總:
\begin{equation}\label{eq:as-cost}
  c_{ij}=\mathrm{clamp}\!\bigl(1-A_{ij}e^{-\Pi_{ij}},\,0,\,1\bigr),\qquad
  \Pi_{ij}=P_{\mathrm{OAO}}+P_{\mathrm{vel}}+P_{\mathrm{occ\_front}}-R_{\mathrm{stability}}.
\end{equation}
legacy additive path \(c_{ij}=1-A_{ij}+\sum_k P_k\) 仍存在但 baseline 不用。
baseline 只有 \(P_{\mathrm{OAO}}\) 與 \(R_{\mathrm{stability}}\) 起作用;
\(P_{\mathrm{vel}}\)(\texttt{vel\_dir\_weight})與 \(P_{\mathrm{occ\_front}}\)
(\texttt{occ\_state\_enabled})未啟用。

\section{OAO Penalty}\label{sec:as-oao}
依 predicted track-track overlap 計 occlusion 係數,並隨遮擋持續幀數 ramp:
\begin{equation}\label{eq:as-oao}
\begin{aligned}
  o^{\mathrm{base}}_i&=\max_{k\ne i}\mathrm{IoU}\!\bigl(B(x_i),B(x_k)\bigr),\qquad
  o_i=o^{\mathrm{base}}_i\min\!\Bigl(1,\tfrac{d_i}{N_{\mathrm{ramp}}}\Bigr),\\
  P_{\mathrm{OAO}}(i,j)&=\tau_{\mathrm{OAO}}\,o_i\,g_s(s_j).
\end{aligned}
\end{equation}
baseline \(\tau_{\mathrm{OAO}}=0.50\)、\(N_{\mathrm{ramp}}=25\);score 調節
\(g_s\equiv1\)(\texttt{oao\_score\_w} 未設)。

\section{Stability Reward}\label{sec:as-stab}
高度一致的 size-consistency reward(成本側,\(\lambda_{\mathrm{eff}}=\max(\lambda,1)\)):
\begin{equation}\label{eq:as-stab}
  R_{\mathrm{stability}}(i,j)=
  \frac{w_{\mathrm{stab}}/\lambda_{\mathrm{eff}}}
       {1+|h_i-h_j|/\max(h_j,10^{-3})}.
\end{equation}
因 reward 除以 \(\lambda\),進入 \(e^{-\lambda c}\)(\Cref{ch:auction})後其 boost 大致不隨
\(\lambda\) 失衡。baseline \texttt{stability\_cost\_w}=0.20。

\suppinput{supp/07-association-impl}
  % track 主體 2/3
% ===========================================================================
% track 主體 (3/3):Sparse Top-K 與 Auction Assignment
% ===========================================================================
\chapter{Sparse Top-K 與 Auction}\label{ch:auction}

\begin{center}\archmap{track}\end{center}

\begin{contract}{Auction(track 子機制)}
\textbf{輸入}\; 每個 track 的稀疏 cost candidates \(c_{ij}\)\quad
\textbf{輸出}\; 配對 \texttt{trk\_to\_det} / \texttt{det\_to\_trk}\\
\textbf{方法}\; ByteTrack 風格分數級聯 \(\times\) 單輪平行 Bertsekas auction\quad
\textbf{Source}\; \texttt{fused\_sinkhorn\_multistage\_kernel}、
\texttt{parallel\_auction\_shmem\_kernel}\\
\textbf{Baseline}\; 5-stage 級聯;\texttt{sinkhorn\_lambda} 僅作 \(e^{-\lambda c}\) value,
\emph{非}完整 Sinkhorn solve
\end{contract}

\paragraph{演算法傳遞.} cost candidates 先一次算出五個 stage 的 top-k,再依優先序
串行跑級聯;每個 stage reset 價格後跑單輪 auction 並 commit,未配掉的 track/det
流到下一 stage:

\begin{center}\resizebox{\linewidth}{!}{%
\begin{tikzpicture}[node distance=15mm]
  \node[step=Ctrk] (n1) {cost\\candidates};
  \node[step=Ctrk, right=of n1] (n2) {fused top-k\\(5 stages)};
  \node[step=Ctrk, right=of n2] (n3) {級聯\\S0$\to$S1$\to$S1b$\to$S1c$\to$S2};
  \node[step=Ctrk, right=of n3] (n4) {per-stage\\reset price $+$ auction};
  \node[step=Ctrk, right=of n4] (n5) {commit\\\texttt{trk\_to\_det}};
  \draw[flowarr=Ctrk] (n1) -- node[flowlbl,above]{$c_{ij}$} (n2);
  \draw[flowarr=Ctrk] (n2) -- node[flowlbl,above]{$p_{ij}$} (n3);
  \draw[flowarr=Ctrk] (n3) -- node[flowlbl,above]{逐 stage} (n4);
  \draw[flowarr=Ctrk] (n4) -- node[flowlbl,above]{winners}  (n5);
\end{tikzpicture}}\end{center}

\noindent
這\emph{不是} full dense Sinkhorn:每 stage 只跑一輪 bid+commit,price buffer 開頭
reset 為 0,沒有 price-raising 迭代到收斂——實質是單輪平行貪婪配對。

\section{Stage 級聯與 Value}\label{sec:au-stage}
五個 stage 依優先序貪婪,carry-over 只看前面沒配掉的 track/det:

\begin{table}[h]\centering\small\renewcommand{\arraystretch}{1.2}
\begin{tabular}{@{}llll@{}}
\toprule
Stage & Track state & Detection score & Cost cap \\
\midrule
S0 DDA       & confirmed  & \([\text{high},1.1)\)         & \(c_{\mathrm{DDA}}\) (0.12) \\
S1 high      & confirmed  & \([\text{high},1.1)\)         & \(\tau_{\mathrm{match}}\) \\
S1b mid      & confirmed  & \([\text{mid},\text{high})\)  & \(\tau_{\mathrm{match}}\) \\
S1c tentative& tentative  & \([\text{mid},1.1)\)          & \(\tau_{\mathrm{match}}\) \\
S2 low       & confirmed  & \([\text{track},\text{mid})\) & \(\tau_{\mathrm{stage2}}\) \\
\bottomrule
\end{tabular}
\caption{分數級聯(\texttt{run\_stage} 0..4)。}\label{tab:au-stage}
\end{table}

放入 top-k 的 value 與 aspect penalty:
\begin{equation}\label{eq:au-value}
  p_{ij}=e^{-\lambda c_{ij}}\,G_{\mathrm{aspect}}(b_j),\qquad
  r_j=\tfrac{\mathrm{width}_j}{\mathrm{height}_j},
\end{equation}
\begin{equation}\label{eq:au-aspect}
  G_{\mathrm{aspect}}=
  \begin{cases}
    \max(0.5,\,1-(r_j-0.8)), & r_j>0.8\\
    \max(0.5,\,1-5(0.15-r_j)), & r_j<0.15\\
    1, & \text{otherwise.}
  \end{cases}
\end{equation}

\section{Auction Bid}\label{sec:au-bid}
給定 detection 當前 price \(\rho_j\),track \(i\) 的 best/second 與 bid:
\begin{equation}\label{eq:au-bid}
  v_{ij}=p_{ij}-\rho_j,\quad
  j^\ast=\arg\max_j v_{ij},\quad
  \Delta\rho_i=v_{ij^\ast}-v^{(2)}_i+\epsilon,\quad
  \mathrm{bid}_i=\rho_{j^\ast}+\Delta\rho_i.
\end{equation}
single-round 下 \(\Delta\rho\) 只用來決定同一輪內多 track 競標同一 detection 誰勝出。
兩個 bid bias(絕對相加):
\begin{equation}\label{eq:au-bias}
  \mathrm{bid}_i \mathrel{+}=\frac{w_{\mathrm{fresh}}}{1+\mathrm{age}_i},\qquad
  \mathrm{bid}_i \mathrel{+}=\frac{w_{\mathrm{stab,bid}}}{1+|h_i-h_j|/h_j}.
\end{equation}
預設不一致,\textbf{勿一律當 off}:freshness \(w_{\mathrm{fresh}}\)
(\texttt{SACCADE\_FRESHNESS\_W})預設 0(關);stability bid bias
\(w_{\mathrm{stab,bid}}\)(\texttt{SACCADE\_STABILITY\_W})預設 \(0.1\)(\textbf{開})。

\suppinput{supp/08-auction-impl}
      % track 主體 3/3
% ===========================================================================
% track 收尾 (1/2):Track Lifecycle 與 Birth
% ===========================================================================
\chapter{Track Lifecycle 與 Birth}\label{ch:lifecycle}

\begin{center}\archmap{track}\end{center}

\begin{contract}{Lifecycle(track 子機制)}
\textbf{輸入}\; assignment 結果(matched / unmatched track / unmatched detection)\quad
\textbf{輸出}\; 更新後的 track 集合與狀態(tentative/confirmed/lost/removed)\\
\textbf{方法}\; hit-streak 累積 confirm、age 計數 lost/remove\quad
\textbf{Source}\; \texttt{tracker\_gpu.cu}\\
\textbf{Baseline}\; \texttt{new\_track\_thresh: 0.28}、\texttt{confirm\_streak: 3}、
\texttt{confirm\_score\_thresh: 0.50}、\texttt{track\_buffer: 30}、\texttt{per\_seq\_adapt: false}
\end{contract}

\paragraph{演算法傳遞.} assignment 後依配對狀態驅動 track 狀態機;下圖為主要轉移
(tentative 若連續未配會很快 drop,圖中略):

\begin{center}\resizebox{\linewidth}{!}{%
\begin{tikzpicture}[node distance=20mm]
  \node[step=Ctrk] (nw) {未配 det};
  \node[step=Ctrk, right=of nw] (te) {tentative};
  \node[step=Ctrk, right=of te] (co) {confirmed};
  \node[step=Ctrk, right=of co] (lo) {lost};
  \node[step=Ctrk, right=of lo] (rm) {removed};
  \draw[flowarr=Ctrk] (nw) -- node[flowlbl,above]{$s_j{\ge}$new\_thr} (te);
  \draw[flowarr=Ctrk] (te) -- node[flowlbl,above]{streak${\ge}3$}    (co);
  \draw[flowarr=Ctrk] (co) -- node[flowlbl,above]{unmatched}         (lo);
  \draw[flowarr=Ctrk] (lo) -- node[flowlbl,above]{age${>}$buffer}    (rm);
  \draw[flowarr=Ctrk] (lo) to[bend left=32] node[flowlbl,above]{re-match} (co);
\end{tikzpicture}}\end{center}

\section{State 轉移}\label{sec:lc-states}
assignment 後 tracker 更新 state:
\begin{itemize}\setlength\itemsep{2pt}
  \item \textbf{matched}(confirmed/tentative):以 detection \(z_j\) 做 Kalman update
        (\Cref{ch:kalman}),\texttt{hit\_streak}\(+1\),\texttt{age}/\texttt{time\_since\_update}
        歸零。
  \item \textbf{unmatched active track}:\texttt{age}\(\le\)\texttt{track\_buffer} 期間維持
        active/lost,超過 max age 後 remove。
  \item \textbf{unmatched detection}(\(s_j\ge\)\texttt{new\_track\_thresh}):生成新的
        tentative track。
\end{itemize}

\section{Birth 與 Confirm}\label{sec:lc-confirm}
新 tentative track 必須累積足夠連續 evidence 才會被視為 confirmed output:
\begin{equation}\label{eq:lc-confirm}
\begin{aligned}
  \text{confirmed}\iff{}&
  \texttt{hit\_streak}\ge\texttt{confirm\_streak}\\
  &{}\land\bar{s}\ge\texttt{confirm\_score\_thresh}.
\end{aligned}
\end{equation}
baseline:
\begin{center}\small
\begin{tabular}{@{}ll@{\qquad}ll@{}}
\texttt{confirm\_streak} & 3 &
\texttt{confirm\_score\_thresh} & 0.50 \\
\texttt{new\_track\_thresh} & 0.28 &
\texttt{track\_buffer} & 30
\end{tabular}
\end{center}
若干 birth-gate experiments 位於 \path{scripts/eval/config/lifecycle.py},但目前
preset 未啟用(\texttt{per\_seq\_adapt: false})。
普通新生 track 以 tentative 與 \texttt{hit\_streak=1} 起跑;bridge relink 復活的
slot 會保留 lost id 並直接回到 confirmed,不再走 tentative confirm gate。

\suppinput{supp/09-lifecycle-impl}
    % track 收尾 1/2
% ===========================================================================
% track 收尾 (2/2):Bridge Relink
% ===========================================================================
\chapter{Bridge Relink}\label{ch:relink}

\begin{center}\archmap{track}\end{center}

\begin{contract}{Bridge relink(track 子機制)}
\textbf{輸入}\; 新穩定的 candidate track、窗內未配的 lost confirmed track\quad
\textbf{輸出}\; candidate 採用 lost track 的 id(lost slot deactivate)\\
\textbf{方法}\; 雙向中點外推(bidirectional midpoint extrapolation),\emph{非} appearance ReID\quad
\textbf{Source}\; \texttt{tracker\_gpu.cu}(\(\ne\) \texttt{relink\_gate.cu} 的 appearance gate)\\
\textbf{Baseline}\; \texttt{relink\_bridge\_enabled: true}、\texttt{relink\_bridge\_px: 0.25}、
\texttt{margin: 0.05}、\texttt{h\_lo/h\_hi: 0.75/1.33}、\texttt{dir\_bonus: 0.8}、
anchor \texttt{adaptive}
\end{contract}

\paragraph{演算法傳遞.} 對 candidate–lost pair,先取 anchor、做 4 點速度回歸,再雙向
外推算殘差,合成 \(d_{\mathrm{bridge}}\),經 direction bonus 與多重 gate 後 commit:

\begin{center}\resizebox{\linewidth}{!}{%
\begin{tikzpicture}[node distance=12mm]
  \node[step=Ctrk] (n1) {candidate\\$+$ lost pair};
  \node[step=Ctrk, right=of n1] (n2) {anchor\\$(a_x,a_y)$};
  \node[step=Ctrk, right=of n2] (n3) {4-pt velocity\\$v_\ell,v_c$};
  \node[step=Ctrk, right=of n3] (n4) {外推殘差\\$r_{\mathrm{fwd}},r_{\mathrm{bwd}}$};
  \node[step=Ctrk, right=of n4] (n5) {$d_{\mathrm{bridge}}$\\$+$ dir bonus};
  \node[step=Ctrk, right=of n5] (n6) {gates\\$+$ commit};
  \draw[flowarr=Ctrk] (n1) -- (n2);
  \draw[flowarr=Ctrk] (n2) -- node[flowlbl,above]{$\ell,c$} (n3);
  \draw[flowarr=Ctrk] (n3) -- (n4);
  \draw[flowarr=Ctrk] (n4) -- node[flowlbl,above]{$d_h$} (n5);
  \draw[flowarr=Ctrk] (n5) -- node[flowlbl,above]{$d^{(1)}$} (n6);
\end{tikzpicture}}\end{center}

\section{History 與 Anchor}\label{sec:rl-anchor}
每個 observed slot(\texttt{age}=0)保存 center/height history ring,並對框高做 EMA:
\(\bar{h}\leftarrow0.95\bar{h}+0.05h\)。coasting lost track 保留最後 history。
baseline anchor \texttt{adaptive}:\(x\) 永遠用 center;\(y\) 候選為
\(y^{\mathrm{top}}{=}c_y{-}h/2,\ y^{\mathrm{bot}}{=}c_y{+}h/2,\ y^{\mathrm{ctr}}{=}c_y\)。
若平均相鄰 height 變化 \(\le\)\texttt{anchor\_rate}(0.03)退化為 center;否則以 top/bottom
各自四點線性回歸殘差為權重選較穩的 edge:
\begin{equation}\label{eq:rl-anchor}
  w_e=\frac{1}{\mathrm{RSS}_{\mathrm{line}}(y^e)/(\bar{h}^2+10^{-3})+0.01},\quad
  a_y=\frac{w_{\mathrm{top}}y^{\mathrm{top}}+w_{\mathrm{bot}}y^{\mathrm{bot}}}
           {w_{\mathrm{top}}+w_{\mathrm{bot}}}.
\end{equation}

\section{Candidate 條件}\label{sec:rl-cand}
candidate 採用 lost id 的前置條件:
\begin{center}\ttfamily\small
candidate hit\_streak == bridge\_at\quad candidate 有 $\ge$4 history samples\\
lost active 但本幀未配\quad lost state == confirmed\quad
bridge\_min\_lost $\le$ lost\_age $\le$ bridge\_ttl
\end{center}

\section{速度回歸與 \texorpdfstring{\(d_{\mathrm{bridge}}\)}{d\_bridge}}\label{sec:rl-bridge}
對四個等時距 scalar anchor sample \((p_0,p_1,p_2,p_3)\),4 點回歸速度
\(v=\tfrac{3p_3+p_2-p_1-3p_0}{10}\)。令 \(\ell\)=lost exit anchor、\(c\)=candidate
entry anchor、\(g\)=\texttt{lost\_age}、\(h_{\mathrm{ref}}=\max(\tfrac{\bar{h}_\ell+\bar{h}_c}{2},1)\):
\begin{equation}\label{eq:rl-resid}
  r_{\mathrm{fwd}}=\frac{\lVert(\ell+v_\ell g)-c\rVert}{h_{\mathrm{ref}}},\qquad
  r_{\mathrm{bwd}}=\frac{\lVert(c-v_c g)-\ell\rVert}{h_{\mathrm{ref}}}.
\end{equation}
以 lost 端速度大小決定外推 vs 純距離的權重:
\begin{align}\label{eq:rl-dbridge}
  d_h&=\frac{\lVert\ell-c\rVert}{h_{\mathrm{ref}}},\\
  w&=\sqrt{\mathrm{clamp}\!\Bigl(\tfrac{\lVert v_\ell\rVert/h_{\mathrm{ref}}}{0.12},0,1\Bigr)},\\
  d_{\mathrm{bridge}}&=
  w\,\frac{r_{\mathrm{fwd}}+r_{\mathrm{bwd}}}{2}+(1-w)\,d_h.
\end{align}

\section{Direction Bonus}\label{sec:rl-dir}
\texttt{dir\_bonus}>0 且方向相近(\(\cos\theta>0.5\))時,向 cross-track 誤差 \(d_{\mathrm{cross}}\) blend:
\begin{equation}\label{eq:rl-dir}
  \alpha=\min\!\bigl(w_{\mathrm{dir}}\cos^2\theta\,\eta_v\,\eta_g,\,1\bigr),\qquad
  d_{\mathrm{bridge}}\leftarrow(1-\alpha)d_{\mathrm{bridge}}+\alpha\,d_{\mathrm{cross}},
\end{equation}
其中 \(\eta_v=\mathrm{clamp}\bigl(\tfrac{\min(\lVert v_\ell\rVert,\lVert v_c\rVert)}
{\max(0.005h_{\mathrm{ref}},10^{-3})},0,1\bigr)\)、\(\eta_g=\mathrm{clamp}(g/30,0,1)\)。
baseline \texttt{dir\_bonus}=0.8。

\section{Gates 與 Commit}\label{sec:rl-gate}
接受條件:
\begin{equation}\label{eq:rl-gate}
  d_{\mathrm{bridge}}\le\tau_{\mathrm{bridge}},\quad
  \frac{\bar{h}_\ell}{\bar{h}_c}\in[h_{\mathrm{lo}},h_{\mathrm{hi}}],\quad
  d^{(2)}_{\mathrm{bridge}}-d^{(1)}_{\mathrm{bridge}}\ge m_{\mathrm{bridge}}.
\end{equation}
注意 \texttt{relink\_bridge\_px}=0.25 是歷史命名,\emph{非} 0.25 pixel——它與已除以
\(h_{\mathrm{ref}}\) 的 \(d_{\mathrm{bridge}}\) 比較,單位是 reference-height-normalized
distance。baseline \texttt{spatial\_gate}=0.0(關),不啟用 occupancy expansion。
接受後 candidate 採用 lost id、lost slot deactivate;多 candidate 競爭同一 lost id 時,
較高 detection score 贏,candidate index 作 tie-breaker。

\suppinput{supp/10-relink-impl}
       % track 收尾 2/2

% ---- 附錄:reference 性質的大表(符號↔config)放這,不擋正文節奏 ----
\appendix
% ===========================================================================
% 附錄 A:符號表 與 符號↔config 對照(math_model §3 / §3.1 / §3.2)
% reference 性質,故置於附錄。代碼錨點(cu:N 行號)見各章 L3 實作補充(full 版)。
% ===========================================================================
\chapter{符號與 Config 對照}\label{app:symbols}

baseline 值對應 \texttt{mamba\_whole\_graph}(frozen\_v2)。各機制的代碼錨點
(kernel 名、\texttt{cu:N} 行號)見對應章末的 L3 實作補充。

\section{基本符號}\label{app:basic}

\begin{longtable}{@{}>{\raggedright\arraybackslash}p{3.3cm}
  >{\raggedright\arraybackslash}p{8.2cm}@{}}
\toprule
符號 & 意義 \\* \midrule \endhead
\(f\) & 目前 frame index \\
\(I_f\) & 目前 RGB/CHW frame tensor \\
\(D_f=\{d_j\}\) & frame \(f\) postprocess 後的 detection set \\
\(b_j=(x_1,y_1,x_2,y_2)\) & detection box(original-frame pixels) \\
\(s_j,\ \mathrm{cls}_j\) & detection score / class id \\
\(T_f=\{t_i\}\) & association 前 active 或 lost tracker slots \\
\(x_i,\ P_i\) & track slot \(i\) 的 8D Kalman state / \(8\times8\) covariance \\
\(W_f\) & GMC \(2\times3\) camera warp;translation path \([1,0,t_x;0,1,t_y]\) \\
\midrule
\(x=(c_x,c_y,a,h,\dot{\cdot})\) & center、aspect、height 與對應速度 \\
\(w=a\cdot h,\ B(x)\) & state implied width / 還原的 box \\
\(z=(c_x,c_y,a,h),\ M=[I_4\ 0]\) & Kalman measurement / measurement matrix \\
\(\mathrm{IoU}(i,j),\ c_{ij}\) & \(B(x_i)\) 與 \(b_j\) 的 IoU / final association cost \\
\(p_{ij},\ A_{ij},\ \Pi_{ij}\) & auction value / base quality / penalty 總和 \\
\bottomrule
\end{longtable}

% 共用:4 欄 config 對照表的欄型(footnotesize,長名稱用 \newline 斷行)
\section{GMC}\label{app:gmc}
{\footnotesize\setlength{\tabcolsep}{3pt}
\begin{longtable}{@{}>{\raggedright\arraybackslash}p{2.8cm}
                     >{\raggedright\arraybackslash}p{2.4cm}
                     >{\raggedright\arraybackslash}p{3.8cm}
                     >{\raggedright\arraybackslash}p{1.6cm}@{}}
\toprule
符號 & 用途 & config / env & baseline \\* \midrule \endhead
\(W_f\) & 相機運動補償 warp & — & trans-only \\
PCR & phase-corr 峰值可信度 & — & — \\
\(\tau_{\mathrm{PCR}}\) & 低可信度縮小位移 & \texttt{SACCADE\_GMC\_PCR\_THRESH} & py \(5.0\) \\
\bottomrule
\end{longtable}}

\section{Kalman}\label{app:kalman}
{\footnotesize\setlength{\tabcolsep}{4pt}
\begin{longtable}{@{}>{\raggedright\arraybackslash}p{1.8cm}
                     >{\raggedright\arraybackslash}p{3.0cm}
                     >{\raggedright\arraybackslash}p{4.2cm}
                     >{\raggedright\arraybackslash}p{1.7cm}@{}}
\toprule
符號 & 用途 & config / env & baseline \\* \midrule \endhead
\(x,z,F,M\) & 8D CV state / 4D measurement & — & — \\
\(\sigma_p{=}h^-/20\) & 位置過程噪聲 & hardcoded & 1/20 \\
\(\sigma_v{=}h^-/160\) & 速度過程噪聲 & hardcoded & 1/160 \\
\(r_{\mathrm{scale}}\) & 測量噪聲縮放 & \texttt{kalman\_r\_scale} & 2.8 \\
\(m_{\mathrm{NSA}},\lambda_{\mathrm{light}}\) & NSA/亮度調節 & — & 1 / 0 \\
\(\tau_{\mathrm{maha}}\) & Mahalanobis gate & \texttt{maha\_gate} & 見實作 \\
\bottomrule
\end{longtable}}

\section{Association Cost}\label{app:assoc}
{\footnotesize\setlength{\tabcolsep}{4pt}
\begin{longtable}{@{}>{\raggedright\arraybackslash}p{1.8cm}
                     >{\raggedright\arraybackslash}p{3.0cm}
                     >{\raggedright\arraybackslash}p{4.2cm}
                     >{\raggedright\arraybackslash}p{1.7cm}@{}}
\toprule
符號 & 用途 & config / env & baseline \\* \midrule \endhead
\(A_{ij}\) & IoU(+ReID)綜合質量 & — & = IoU \\
\(w_{\mathrm{fuse}}\) & 低分檢測降權 & \texttt{fuse\_score\_weight} & 0.0 \\
\(c_{ij}\) & 乘法式 cost & \texttt{multiplicative\_cost} & true \\
\(\lambda\) & cost\(\to\)value 溫度 & \texttt{sinkhorn\_lambda} & 10 \\
\(o_i,\tau_{\mathrm{OAO}}\) & 遮擋配對抑制 & \texttt{oao\_tau} /\newline \texttt{oao\_ramp\_frames} & 0.50 / 25 \\
\(w_{\mathrm{vel}}\) & 速度反向懲罰 & \texttt{vel\_dir\_weight} & 關 \\
\(w_{\mathrm{occ}}\) & front-occluder 懲罰 & \texttt{occ\_cost\_weight} & 關 \\
\(w_{\mathrm{stab}}\) & 高度一致 reward(成本側) & \texttt{stability\_cost\_w} & 0.20 \\
\bottomrule
\end{longtable}}

\section{Auction}\label{app:auction}
{\footnotesize\setlength{\tabcolsep}{4pt}
\begin{longtable}{@{}>{\raggedright\arraybackslash}p{1.8cm}
                     >{\raggedright\arraybackslash}p{3.0cm}
                     >{\raggedright\arraybackslash}p{4.2cm}
                     >{\raggedright\arraybackslash}p{1.7cm}@{}}
\toprule
符號 & 用途 & config / env & baseline \\* \midrule \endhead
\(p_{ij}\) & auction 概率值 & — & — \\
\(G_{\mathrm{aspect}}\) & 抑制異常長寬比 & hardcoded (0.8/0.15) & — \\
\(w_{\mathrm{fresh}}\) & 新鮮度 bid bias & \texttt{SACCADE\_FRESHNESS\_W} & 0.0 \\
\(w_{\mathrm{stab,bid}}\) & 高度一致 bid bias & \texttt{SACCADE\_STABILITY\_W} & 0.1(開) \\
S0 DDA & confirmed×high 更緊 stage & \texttt{SACCADE\_ENABLE\_DDA} /\newline \texttt{\_DDA\_MAX\_COST} & on / 0.12 \\
stage thr & 分數級聯邊界 & \texttt{match} / \texttt{high} / \texttt{mid} /\newline \texttt{track} / \texttt{stage2} & .50/.45/\newline.10/.05/.50 \\
\bottomrule
\end{longtable}}

\section{Lifecycle / Bridge / Semantic relink}\label{app:rest}
{\footnotesize\setlength{\tabcolsep}{4pt}
\begin{longtable}{@{}>{\raggedright\arraybackslash}p{1.8cm}
                     >{\raggedright\arraybackslash}p{3.0cm}
                     >{\raggedright\arraybackslash}p{4.2cm}
                     >{\raggedright\arraybackslash}p{1.7cm}@{}}
\toprule
符號 & 用途 & config / env & baseline \\* \midrule \endhead
birth/\newline confirm & tentative\(\to\)confirmed & \texttt{new\_track\_thresh} /\newline \texttt{confirm\_streak} & 0.28 / 3 \\
\(\tau_{\mathrm{bridge}}\) & 雙向外推殘差門檻 & \texttt{relink\_bridge\_px} & 0.25 \\
\(\alpha\) & 方向一致偏移 & \texttt{relink\_bridge\_dir\_bonus} & 0.8 \\
\(h_{\mathrm{lo}},h_{\mathrm{hi}}\) & 高度比 gate & \texttt{relink\_bridge\_h\_lo} /\newline \texttt{\_h\_hi} & 0.75/1.33 \\
\(m_{\mathrm{bridge}}\) & best-vs-second margin & \texttt{relink\_bridge\_margin} & 0.05 \\
\(w_{\mathrm{sim/iou}}\),\newline \(w_{\mathrm{maha}}\) & semantic joint 權重 & \texttt{semantic\_w\_*\_base} & off \\
\bottomrule
\end{longtable}}

\section{方法出處 / 命名對照}\label{app:naming}
\begin{itemize}\setlength\itemsep{2pt}
  \item \textbf{GMC} = phase correlation(cross-power spectrum + Hanning window),translation-only warp。
  \item \textbf{Kalman} = SORT/DeepSORT 風格 constant-velocity filter。
  \item \textbf{Assignment} = Bertsekas auction(單輪平行貪婪)跑在 softmin-temperature top-k 上;
        分數分段 = ByteTrack 風格 cascade。\texttt{sinkhorn\_lambda} 為歷史命名——只用
        \(e^{-\lambda c}\) 當 value,\textbf{非}完整 Sinkhorn 迭代。
  \item \textbf{OAO} = occlusion-aware(track-track overlap)配對抑制 + duration ramp。
  \item \textbf{Bridge relink} = 雙向中點外推,項目自有機制,非標準 appearance ReID。
  \item \textbf{Semantic relink gate} = appearance + Mahalanobis + IoU joint gate(baseline 關)。
\end{itemize}

\paragraph{提醒.} \(w_{\mathrm{stab}}\)(\Cref{sec:as-stab},成本側 reward)與
\(w_{\mathrm{stab,bid}}\)(\Cref{sec:au-bid},auction bid bias)是\textbf{兩個不同旋鈕},
數值與作用點皆不同,雖都用高度一致性 \(|h_i-h_j|/h_j\)。


\end{document}
"
% ===========================================================================
\documentclass[11pt,oneside]{ctexbook}   % 中文 → ctexbook + XeLaTeX
% ===========================================================================
% preamble.tex —— 套件、深度開關、archmap 迷你地圖、合約框
% 引擎:XeLaTeX(中文必須)。建置見 README.md / Makefile。
% ===========================================================================

% --- 繁體中文字型(fandol 為簡體取向,缺繁體字形;改用 Noto CJK TC)---
\setCJKmainfont{Noto Serif CJK TC}
\setCJKsansfont{Noto Sans CJK TC}
\setCJKmonofont{Noto Sans CJK TC}
% ctexbook 預設 caption 是簡體;改回繁體(录→錄)。appendixname 同時管標題與頁首。
\ctexset{contentsname=目錄, appendixname=附錄}
% 中英混排時,允許行末額外伸縮以吸收邊際 overfull(標準作法)
\setlength{\emergencystretch}{3em}

% --- 數學 / 表格 / 引用 ---
\usepackage{amsmath, amssymb, mathtools}
\usepackage{booktabs}            % \toprule \midrule \bottomrule,禁直線
\usepackage{array}
\usepackage{tabularx}            % 自動寬度欄(職責表用)
\usepackage{longtable}           % 可跨頁的參考表(附錄符號表用)
\usepackage[hidelinks]{hyperref}
\usepackage{cleveref}            % \Cref{eq:...} 就近引用
\usepackage{xstring}             % archmap 的字串比較(跨引擎安全)
\usepackage{fancyhdr}            % 統一頁碼位置

% --- 圖 ---
\usepackage{tikz}
\usetikzlibrary{arrows.meta, positioning, fit, backgrounds}

% --- 程式碼 / 合約框 ---
\usepackage[most]{tcolorbox}
\usepackage{fvextra}             % Verbatim,ASCII pipeline 圖暫存用

% --- 頁碼 ---
% ctexbook/book 的一般頁與章首頁預設 page style 不同,會讓頁碼在頁首/頁腳間跳。
% 這裡同時覆寫一般頁型與 plain 頁型,固定頁碼在頁腳中央。
\fancypagestyle{saccadepage}{
  \fancyhf{}
  \fancyfoot[C]{\thepage}
  \renewcommand{\headrulewidth}{0pt}
  \renewcommand{\footrulewidth}{0pt}
}
\fancypagestyle{plain}{
  \fancyhf{}
  \fancyfoot[C]{\thepage}
  \renewcommand{\headrulewidth}{0pt}
  \renewcommand{\footrulewidth}{0pt}
}
\AtBeginDocument{\pagestyle{saccadepage}}

% ===========================================================================
% 深度開關:L3 實作補充是否編入
%   make full  -> \def\IMPL{1}  含 supp/*-impl.tex
%   make paper -> \def\IMPL{0}  剝掉所有 L3,得乾淨學術版
% ===========================================================================
\providecommand{\IMPL}{1}
\newif\ifimpl
\ifnum\IMPL=1 \impltrue \else \implfalse \fi

% \suppinput{file}:只有 full 版會把 L3 補充檔讀進來
\newcommand{\suppinput}[1]{\ifimpl\input{#1}\fi}

% L3 補充章節的統一外框(視覺上和學術本體區隔)
\newtcolorbox{implsupp}[1]{
  enhanced, breakable,
  colback=orange!4, colframe=orange!55!black,
  fonttitle=\bfseries\small, coltitle=orange!30!black,
  title={實作補充 (L3,僅 full 版)\quad#1},
  left=2mm, right=2mm, top=1mm, bottom=1mm}

% ===========================================================================
% 模組合約框:每個細節章開頭,濃縮 I/O / Source / Baseline
% ===========================================================================
\newtcolorbox{contract}[1]{
  enhanced, breakable,
  colback=gray!4, colframe=black!55,
  fonttitle=\bfseries, title={模組合約:#1}}

% ===========================================================================
% \archmap{focus} —— 全書共用的「你在這裡」迷你地圖
%   focus ∈ {none, detect, gmc, track, output}
%   none = 全亮(總覽章用);其餘 = 點亮焦點、淡化其它
%   track 內的子章(Kalman/association/auction/...)先一律用 track;
%   日後要更細,可在此複製出 \trackmap 畫 track 內部子階。
% ===========================================================================
% 每個 stage 的代表色(改這裡就改全書配色)——用 \colorlet 建具名顏色,
% 才能在 pgfkeys option 裡以字面名稱引用(巨集名不會被展開)。
\colorlet{Cbuf}{teal}
\colorlet{Cdet}{blue!65!cyan}
\colorlet{Cgmc}{orange!90!red}
\colorlet{Ctrk}{violet!80!blue}
\colorlet{Cout}{green!55!black}

% 全域 style:box 外型 + 三種狀態(吃顏色當 #1)+ 箭頭/標籤
\tikzset{
  ambox/.style  ={draw, rounded corners, align=center,
                  minimum height=10mm, minimum width=24mm, font=\small},
  amfull/.style ={fill=#1!25, draw=#1!85!black, line width=1pt,   text=black!88},
  amfade/.style ={fill=#1!8,  draw=#1!35,       line width=0.6pt, text=black!40},
  amfocus/.style={fill=#1!48, draw=#1!90!black, line width=1.9pt, text=black!92},
  amarr/.style  ={-{Stealth}, thick},               % 直角:不加 rounded corners
  amlbl/.style  ={font=\scriptsize, text=black!55},
}

\newcommand{\archmap}[1]{%
  % 預設 full(總覽 none 全部上色);具名 style 都是字面,pgfkeys 吃得到
  \tikzset{amSbuf/.style={amfull=Cbuf}, amSdet/.style={amfull=Cdet},
           amSgmc/.style={amfull=Cgmc}, amStrk/.style={amfull=Ctrk},
           amSout/.style={amfull=Cout}}%
  \IfStrEq{#1}{none}{}{% 指定焦點:先全 fade,再把焦點改 focus
    \tikzset{amSbuf/.style={amfade=Cbuf}, amSdet/.style={amfade=Cdet},
             amSgmc/.style={amfade=Cgmc}, amStrk/.style={amfade=Ctrk},
             amSout/.style={amfade=Cout}}%
    \IfStrEq{#1}{detect}{\tikzset{amSdet/.style={amfocus=Cdet}}}{}%
    \IfStrEq{#1}{gmc}{\tikzset{amSgmc/.style={amfocus=Cgmc}}}{}%
    \IfStrEq{#1}{track}{\tikzset{amStrk/.style={amfocus=Ctrk}}}{}%
    \IfStrEq{#1}{output}{\tikzset{amSout/.style={amfocus=Cout}}}{}%
  }%
  % 固定成欄寬比例再置中,避免視覺偏移
  \resizebox{0.92\linewidth}{!}{%
  \begin{tikzpicture}
    % 橫式佈局:左→右流,detect/GMC 為上下兩條平行支線
    \node[ambox, amSbuf] (buf) at (0,0)     {frame buffer\\\scriptsize(GPU, CHW)};
    \node[ambox, amSdet] (det) at (3.2, 1)  {detect 鏈};
    \node[ambox, amSgmc] (gmc) at (3.2,-1)  {GMC};
    \node[ambox, amStrk] (trk) at (6.6, 0)  {track\\\scriptsize tracker\_gpu.cu};
    \node[ambox, amSout] (out) at (9.8, 0)  {materialize\\\scriptsize→ MOT rows};
    % 分岔:buf 拉出 bus,再直角接上下兩支線(連線染來源色)
    \draw[amarr, Cbuf!65!black] (buf.east) -- ++(0.4,0) |- (det.west);
    \draw[amarr, Cbuf!65!black] (buf.east) -- ++(0.4,0) |- (gmc.west);
    % 匯合:兩支線直角收進 track
    \draw[amarr, Cdet!65!black] (det.east) -- ++(0.4,0) node[amlbl,above]{boxes} |- (trk.west);
    \draw[amarr, Cgmc!70!black] (gmc.east) -- ++(0.4,0) node[amlbl,below]{$W_f$} |- (trk.west);
    \draw[amarr, Ctrk!65!black] (trk.east) -- (out.west);
  \end{tikzpicture}}}

% ===========================================================================
% 章內「演算法傳遞圖」共用 style —— 每章自畫,顏色用該模組的 C*
%   step=<color>     一個演算法步驟方塊
%   flowarr=<color>  傳遞箭頭(直角時用 |- / -|)
%   flowlbl          標在邊上的「傳了什麼」(白底壓線)
% 用法見 chapters/05-gmc.tex 的範例。
% ===========================================================================
\tikzset{
  step/.style   ={draw, rounded corners=2pt, align=center, font=\footnotesize,
                  minimum height=10mm, minimum width=20mm, inner sep=4pt,
                  fill=#1!14, draw=#1!70!black, text=black!88},
  flowarr/.style={-{Stealth}, semithick, draw=#1!75!black},
  flowlbl/.style={font=\scriptsize, text=black!60, fill=white, inner sep=1.5pt},
}
% 把 GMC 章內流程一行寫完的便利指令(其他章可仿照自行展開)
\newcommand{\gmcflow}{%
  \begin{center}\resizebox{\linewidth}{!}{%
  \begin{tikzpicture}[node distance=13mm]
    \node[step=Cgmc] (n1) {灰階降採樣\\+\,Hanning};
    \node[step=Cgmc, right=of n1] (n2) {FFT\\prev / curr};
    \node[step=Cgmc, right=of n2] (n3) {cross-power\\$R=\overline{A}B/|{\cdot}|$};
    \node[step=Cgmc, right=of n3] (n4) {IFFT $\to$ peak\\PCR};
    \node[step=Cgmc, right=of n4] (n5) {confidence $\gamma$\\$\to$ warp $W_f$};
    \draw[flowarr=Cgmc] (n1) -- node[flowlbl,above]{$G_w$}        (n2);
    \draw[flowarr=Cgmc] (n2) -- node[flowlbl,above]{$A,B$}        (n3);
    \draw[flowarr=Cgmc] (n3) -- node[flowlbl,above]{$R(k)$}       (n4);
    \draw[flowarr=Cgmc] (n4) -- node[flowlbl,above]{$r$ peak}     (n5);
  \end{tikzpicture}}\end{center}}


\title{Saccade 全局數學模型}
\author{Ray Lei}

\begin{document}
% Keep page numbering format uniform across title, TOC, main text, and appendix.
\maketitle
\setcounter{page}{2}
\tableofcontents

% ---- Part I:主架構(the spine,只放地圖,不放數學)----
\part{系統總覽}
% ===========================================================================
% Part I 唯一一章:系統總覽 = 主架構(the spine)
% 原則:這裡只放「地圖」——架構圖、模組邊界、dataflow、baseline 合約。
%       所有數學一律下放到 Part II。這就是「不全展開」的關鍵。
% ===========================================================================
\chapter{系統總覽}\label{ch:overview}

整條 pipeline 分兩層:\textbf{Python orchestration}(排程、buffer 管理、
CUDA graph replay、eval)與 \textbf{C++/CUDA compute}(detector、postprocess、
GMC、tracker association)。Python 不在 hot path 做數值計算,只在 output
boundary 把結果搬回 host。

關鍵:detect 與 GMC 是從\textbf{同一個 frame buffer 分岔的兩條平行支線}——
GMC 吃 frame 灰階(+上一幀),\emph{不是} detect 的輸出;\texttt{track} 才把
兩條支線匯合。

\section{架構地圖}\label{sec:archmap}
全書各章開頭都會重畫這張圖、點亮當前模組(「你在這裡」)。總覽版全亮:

\begin{center}
  \archmap{none}
\end{center}

\section{模組職責與 I/O}\label{sec:modules}
下表概括每個模組\textbf{吃什麼、吐什麼、解決什麼問題},以及在 pipeline 中銜接誰。
逐項數學見 Part II 對應章;傳遞合約與 source 見 \Cref{tab:transfer}。

\begin{table}[h]\small\setlength{\tabcolsep}{4pt}\renewcommand{\arraystretch}{1.35}
\begin{tabularx}{\linewidth}{@{}l >{\raggedright\arraybackslash}X >{\raggedright\arraybackslash}X l@{}}
\toprule
模組 & 輸入 → 輸出 & 解決的問題 & 銜接 \\
\midrule
detect      & frame CHW tensor → raw boxes/scores/classes
            & 從單張影像偵測行人(TRT backbone + Mamba head,anchor-free decode)
            & → postprocess \\
postprocess & raw det → 過濾/NMS 後 boxes
            & 去除低分與重疊框,交給 tracker 乾淨的偵測集
            & detect → track \\
GMC         & prev+curr 灰階 → \(2\times3\) warp \(W_f\)
            & 估相機運動補償平移,讓 Kalman 的 velocity 只學物體運動
            & 與 detect 平行,補 track \\
track       & boxes/scores/\(W_f\) → track states / IDs
            & 跨幀關聯與維持 ID(分解見 \Cref{tab:track-internal})
            & 匯合 → materialize \\
output      & GPU state → MOT rows
            & 搬回 host、產生 MOT 結果(含 tracklet 內插後處理)
            & track → 終點 \\
\bottomrule
\end{tabularx}
\caption{模組職責與輸入/輸出概覽。}\label{tab:modules}
\end{table}

\texttt{track} 不是單一步驟,而由數個子機制串成,各為 Part II 的一章:

\begin{table}[h]\small\setlength{\tabcolsep}{4pt}\renewcommand{\arraystretch}{1.3}
\begin{tabularx}{\linewidth}{@{}l >{\raggedright\arraybackslash}X@{}}
\toprule
子機制 & 解決的問題 \\
\midrule
Kalman 模型      & constant-velocity 運動模型;每幀 predict 再用 detection update track state \\
Association cost & 把 IoU、OAO/velocity penalty、stability reward 匯成單一配對成本 \(c_{ij}\) \\
Auction 指派     & ByteTrack 風格分數級聯 + 單輪平行 auction,決定 track 與 detection 配對 \\
Track lifecycle  & track 生死:tentative→confirmed 的 birth/confirm,與 lost/remove \\
Bridge relink    & 找回短暫消失又重生的 ID(雙向中點外推,非 appearance ReID) \\
\bottomrule
\end{tabularx}
\caption{\texttt{track} 內部子機制(Part II 逐章展開)。}\label{tab:track-internal}
\end{table}

傳遞合約(誰留在 GPU、source 檔)另列:

\begin{table}[h]\small\setlength{\tabcolsep}{4pt}\renewcommand{\arraystretch}{1.2}
\begin{tabularx}{\linewidth}{@{}l
  >{\raggedright\arraybackslash}X
  >{\raggedright\arraybackslash}X@{}}
\toprule
模組 & 模組(Python facade / C++·CUDA core) & 傳遞方式 \\
\midrule
detect      & \texttt{mamba\_gated\_detector.py} / TRT backbone + Mamba head & GPU tensor(whole-detect graph)\\
postprocess & \texttt{PerceptionPipeline} / \texttt{pipeline.cpp} & GPU buffer,留 device \\
GMC         & \texttt{GMC} / \texttt{gmc\_kernel.cu} (cuFFT)       & GPU tensor,$2\times3$ warp \\
track       & \texttt{GPUByteTracker} / \texttt{tracker\_gpu.cu}   & device pointer + stream \\
materialize & \texttt{materialize\_\allowbreak gpu\_\allowbreak track\_\allowbreak results} & GPU→host,僅一次 copy \\
output      & \texttt{relink\_write}                               & host(fast MOT emit)\\
\bottomrule
\end{tabularx}
\caption{傳遞合約與 source。}\label{tab:transfer}
\end{table}

\section{Frame-Level Dataflow}\label{sec:dataflow}
對每個 frame \(f\),no-ReID baseline 的主線是兩條同源分支在 tracker 匯合:
\begin{center}
\small
\begin{tabular}{@{}rcl@{}}
frame tensor & \(\to\) & detect \(\to\) postprocess \(\to\) boxes \\
frame tensor & \(\to\) & GMC warp \(W_f\) \\
boxes \(+\ W_f\) & \(\to\) & tracker update \(\to\) materialize \\
materialize & \(\to\) & relink\_write / MOT rows
\end{tabular}
\end{center}

\section{Baseline 合約}\label{sec:baseline}
本文以 \texttt{mamba\_whole\_graph} preset 為基準。關鍵啟用條件:

\begin{table}[h]\small\setlength{\tabcolsep}{4pt}
\begin{tabularx}{\linewidth}{@{}l
  >{\raggedright\arraybackslash}p{4.8cm}
  >{\raggedright\arraybackslash}X@{}}
\toprule
區域 & baseline 值 & 效果 \\
\midrule
Graph     & \texttt{use\_whole\_graph: true}      & whole detect graph 單次 replay \\
GMC       & \texttt{gmc\_downscale: 4}            & GPU phase-correlation 估平移 \\
Appearance& \texttt{reid\_mode: off}              & headline 不用 appearance \\
Assoc     & \texttt{match\_thresh: 0.50}          & ByteTrack-like 多階段 gate \\
Cost      & \texttt{multiplicative\_cost: true}   & log-linear cost + stability reward \\
Relink    & \texttt{relink\_bridge\_enabled: true}& 雙向 foot bridge \\
\bottomrule
\end{tabularx}
\caption{Baseline 合約摘要。完整旋鈕表見附錄 A。}\label{tab:baseline}
\end{table}

% TODO: 把 math_model.md §3 的符號↔config 大表搬進 chapters/A-symbols.tex(附錄)


% ---- Part II:模組細節(一章一模組,全 pipeline 等深 L1–L2)----
% 每章 = \archmap{...} 迷你地圖 + contract 合約框 + 核心數學
%        + 末尾 \suppinput 掛 L3 實作補充(paper 版會自動剝掉)
\part{模組細節}
% ===========================================================================
% Detection 與 Postprocess(pipeline 最前段,故置於 Part II 第一章)
% ===========================================================================
\chapter{Detection 與 Postprocess}\label{ch:detect}

\begin{center}\archmap{detect}\end{center}

\begin{contract}{detect 鏈}
\textbf{輸入}\; frame \(640\times640\) CHW tensor\quad
\textbf{輸出}\; boxes / scores / classes(post-NMS)\\
\textbf{方法}\; YOLO26s TRT backbone+neck \(\to\) Mamba head(per-level lane)\(\to\)
anchor-free decode \(\to\) NMS\quad
\textbf{Source}\; \texttt{mamba\_gated\_detector.py}、\texttt{mamba\_head.py}、
\texttt{mamba\_gated\_detector.cpp}、\texttt{pipeline.cpp}\\
\textbf{Baseline}\; \texttt{tiling: native\_640}、\texttt{preprocess: none}、
ckpt \texttt{mamba\_gt\_v14replica\_t3\_t1};整段在 whole-detect CUDA graph 單次 replay
\end{contract}

\paragraph{演算法傳遞.} detect 鏈分兩段:detector(backbone+head,出 raw)與 native
postprocess(過濾/NMS)。邊上為各段傳遞的張量:

\begin{center}\resizebox{\linewidth}{!}{%
\begin{tikzpicture}[node distance=14mm]
  \node[step=Cdet] (n1) {frame\\$640^2$};
  \node[step=Cdet, right=of n1] (n2) {TRT backbone\\$+$ neck (FPN)};
  \node[step=Cdet, right=of n2] (n3) {Mamba head\\$\times(P_3,P_4,P_5)$};
  \node[step=Cdet, right=of n3] (n4) {decode\\anchor-free};
  \node[step=Cdet, right=of n4] (n5) {filter $+$ NMS};
  \draw[flowarr=Cdet] (n1) -- node[flowlbl,above]{image}              (n2);
  \draw[flowarr=Cdet] (n2) -- node[flowlbl,above]{$F_3,F_4,F_5$}       (n3);
  \draw[flowarr=Cdet] (n3) -- node[flowlbl,above]{$\mathrm{cls}_N,\mathrm{reg}_N$} (n4);
  \draw[flowarr=Cdet] (n4) -- node[flowlbl,above]{8400 anchors}        (n5);
\end{tikzpicture}}\end{center}

\section{整體結構與三尺度}\label{sec:det-overall}
P3/P4/P5 三條 lane 結構完全相同,只差尺度常數(輸入固定 640):

\begin{table}[h]\centering\small\renewcommand{\arraystretch}{1.2}
\begin{tabular}{@{}lcccc@{}}
\toprule
level \(N\) & stride \(s_N\) & \(H_N{=}W_N{=}640/s_N\) & channels \(C_N\) & tokens \(L_N{=}(H_N/4)^2\) \\
\midrule
P3 & 8  & 80 & 128 & 400 \\
P4 & 16 & 40 & 256 & 100 \\
P5 & 32 & 20 & 512 & 25  \\
\bottomrule
\end{tabular}
\caption{三尺度常數(\(s_N{=}2^N,\ C_N{=}2^{N+4}\))。}\label{tab:det-levels}
\end{table}

\section{單一 Lane:Mamba Head(Flow B)}\label{sec:det-mamba}
重點是 U-Net 式 skip:\(X_N\) 一邊進 down\(\to\)Mamba\(\to\)up,一邊\textbf{跳接}直接與
\(U_N\) concat,故 head 輸入是 256 ch。先 \texttt{Down4} 把序列壓到 \(L_N\le400\)
(控 scan 成本),上採樣後再補回細節:

\begin{center}\resizebox{\linewidth}{!}{%
\begin{tikzpicture}[node distance=11mm]
  \node[step=Cdet] (f)  {$F_N$\\$C_N{\times}H_N{\times}W_N$};
  \node[step=Cdet, right=of f]  (x)  {input\_proj\\$X_N$ (128)};
  \node[step=Cdet, right=of x]  (d)  {Down4\\stride-4};
  \node[step=Cdet, right=of d]  (z)  {flatten\\$Z_N$ ($L_N{\times}128$)};
  \node[step=Cdet, right=of z]  (m)  {MambaBlock\\$\times2$ ($d_{\mathrm{state}}{=}16$)};
  \node[step=Cdet, right=of m]  (u)  {Up4\\$U_N$ (128)};
  \node[step=Cdet, right=of u]  (c)  {concat\\$\to Y_N$ (256)};
  \draw[flowarr=Cdet] (f) -- (x);
  \draw[flowarr=Cdet] (x) -- (d);
  \draw[flowarr=Cdet] (d) -- node[flowlbl,above]{$\le400$ tok} (z);
  \draw[flowarr=Cdet] (z) -- (m);
  \draw[flowarr=Cdet] (m) -- (u);
  \draw[flowarr=Cdet] (u) -- (c);
  % U-Net skip:X_N 直接跳接到 concat(直角繞上方)
  \draw[flowarr=Cdet, dashed] (x.north) |- ++(0,7mm) -| node[flowlbl,above,pos=0.25]{skip $X_N$} (c.north);
\end{tikzpicture}}\end{center}

\section{Head 與 Decode(Flow C/D)}\label{sec:det-decode}
\textbf{Head}(per level,純卷積):\(Y_N(256)\) 分兩支,cls head
\(\to\mathrm{cls}_N\,(80)\)、reg head \(\to\mathrm{reg}_N\,(4)\);輸出通道
\(n_o=n_c+4\,r_{\max}=80+4=84\)(\(r_{\max}{=}1\),box 直接 4 維距離,無 DFL)。

\textbf{Decode}(anchor-free,跨三尺度合併,\(N=80^2{+}40^2{+}20^2=8400\) anchors)。
anchor 為 grid 中心 \((x{+}0.5,y{+}0.5)\),每點帶 stride:
\begin{equation}\label{eq:det-decode}
  \mathrm{reg}=(\ell t,rb)\ \xrightarrow{\text{dist2bbox}}\
  c=\mathrm{anchor}+\tfrac{rb-\ell t}{2},\quad
  wh=\ell t+rb\ \xrightarrow{\times s}\ \text{xyxy(像素)}.
\end{equation}
\begin{equation}\label{eq:det-cls}
  \text{score},\ \text{class}=\max_{80},\ \arg\max_{80}\ \sigma(\mathrm{cls}).
\end{equation}

\section{Postprocess Contract}\label{sec:det-post}
native postprocess(\Cref{ch:detect} 第二段)的核心 contract:
\begin{center}\ttfamily\small
(raw boxes/scores/classes) → score/class/geometry filter → compact/gather → NMS → output
\end{center}
baseline 旗標全關:
\begin{center}\small
\begin{tabular}{@{}ll@{\qquad}ll@{}}
\texttt{track\_person\_only} & false &
\texttt{person\_geometry\_prior} & false \\
\texttt{detection\_quality\_scaling} & false &
\texttt{geometry\_suspect\_support} & false
\end{tabular}
\end{center}
tracker 直接消費 post-NMS boxes,各 detection 進哪段 association 由 tracker
的 score gates(\(\tau_{\mathrm{track/mid/high}}\)、\texttt{new\_track\_thresh})決定。

\suppinput{supp/04-detection-impl}
    % pipeline 最前段
% ===========================================================================
% 細節章「樣板」:GMC。後續每章都照這個骨架寫,確保全 pipeline 等深。
%   1) \archmap{...}  迷你地圖,點亮本模組
%   2) contract 框    I/O / Source / Baseline,讓人不鑽公式也拿得到全貌
%   3) L1–L2 數學     學術本體(收 paper 用這層)
%   4) \suppinput     掛 L3 實作補充(paper 版自動剝掉)
% ===========================================================================
\chapter{GMC:相機運動補償}\label{ch:gmc}

\begin{center}\archmap{gmc}\end{center}

\begin{contract}{GMC}
\textbf{輸入}\; prev + curr 灰階(GPU tensor, CHW float \([0,1]\))\quad
\textbf{輸出}\; \(2\times3\) translation warp \(W_f\)\\
\textbf{方法}\; phase correlation(cross-power spectrum + Hanning window)\quad
\textbf{Source}\; \texttt{gmc\_kernel.cu}(cuFFT)\\
\textbf{Baseline}\; \texttt{gmc\_downscale: 4}、translation-only、soft confidence gate
\end{contract}

\noindent
GPU GMC 用 phase correlation 估 frame-to-frame 平移,作為 deterministic control
input 送進 \texttt{track} 的 Kalman predict,讓 velocity state 只學物體
residual motion 而非相機運動。

\paragraph{演算法傳遞.} 本章內部依序串接下列步驟,邊上標註各步之間傳遞的量
(對應 \Cref{sec:gmc-down,sec:gmc-cps,sec:gmc-warp}):

\gmcflow

\section{灰階降採樣與窗函數}\label{sec:gmc-down}
對 downscaled pixel \((x,y)\),floor-based 取樣:
\begin{equation}\label{eq:gmc-sample}
  s_x=\Big\lfloor x\tfrac{W_{\mathrm{src}}}{W_{\mathrm{dst}}}\Big\rfloor,\qquad
  s_y=\Big\lfloor y\tfrac{H_{\mathrm{src}}}{H_{\mathrm{dst}}}\Big\rfloor,\qquad
  G=0.299R+0.587G_c+0.114B.
\end{equation}
FFT 前套 Hanning window \(G_w(x,y)=G(x,y)\,w_x w_y\),其中
\(w_x=\tfrac12(1-\cos\tfrac{2\pi x}{W_{\mathrm{dst}}-1})\)(同理 \(w_y\))。

\section{Cross-Power Spectrum 與位移}\label{sec:gmc-cps}
令 \(A=\mathrm{FFT}(\text{prev})\)、\(B=\mathrm{FFT}(\text{curr})\):
\begin{equation}\label{eq:gmc-cps}
  C(k)=\overline{A(k)}B(k),\qquad
  R(k)=\frac{C(k)}{|C(k)|+10^{-6}},\qquad
  r=\mathrm{IFFT}(R).
\end{equation}
\(r\) 的 peak 給出 wrapped displacement(\(\text{peak}>\text{dim}/2\) 時減去該維長度)。

\section{Confidence 與 Warp}\label{sec:gmc-warp}
phase-correlation 可信度與軟縮放:
\begin{equation}\label{eq:gmc-pcr}
  \mathrm{PCR}=\frac{\max r}{\mathrm{RMS}(r)},\qquad
  \gamma=\begin{cases}\mathrm{clamp}(\mathrm{PCR}/\tau,0,1),&\mathrm{PCR}<\tau\\1,&\text{otherwise.}\end{cases}
\end{equation}
位移超過 downscaled 尺寸 25\% 視為不可信(退回 identity)。否則
\begin{equation}\label{eq:gmc-final}
  t_x=p_x d\,\gamma,\quad t_y=p_y d\,\gamma,\qquad
  W_f=\begin{bmatrix}1&0&t_x\\0&1&t_y\end{bmatrix}.
\end{equation}
\textbf{Sign convention}:positive \(t_x\) 把 predicted track center 往
current frame 右方推。

% ---------------------------------------------------------------------------
% L3 實作補充:kernel 名、env、四路徑 fallback 差異。paper 版自動剝掉。
% ---------------------------------------------------------------------------
\suppinput{supp/05-gmc-impl}
          % ← 完整範例章
% ===========================================================================
% track 主體 (1/3):Kalman 運動模型
% ===========================================================================
\chapter{Kalman 運動模型}\label{ch:kalman}

\begin{center}\archmap{track}\end{center}

\begin{contract}{Kalman(track 子機制)}
\textbf{輸入}\; 前一幀 track state \((x,P)\)、GMC warp \(W_f\)、配對到的 detection
measurement \(z\)\quad
\textbf{輸出}\; predict 後 \(\tilde{x}\)(供 gate/cost)、update 後 \((x,P)\)\\
\textbf{方法}\; SORT/DeepSORT 風格 constant-velocity filter,state
\((c_x,c_y,a,h,\dot{\cdot})\)\quad
\textbf{Source}\; \texttt{kalman\_gpu.cuh}\quad
\textbf{Baseline}\; \texttt{kalman\_r\_scale: 2.8}、NSA off
\end{contract}

\noindent
state 與 measurement:
\begin{equation}\label{eq:kf-state}
  x=(c_x,c_y,a,h,v_x,v_y,v_a,v_h)^{\!\top},\qquad
  z=(c_x,c_y,a,h)^{\!\top}.
\end{equation}

\paragraph{演算法傳遞.} 每幀對每個 track 依序:predict(先吃 GMC control)→ 算
innovation 協方差 → gate → 若入選則 update。邊上為各步傳遞的量:

\begin{center}\resizebox{\linewidth}{!}{%
\begin{tikzpicture}[node distance=14mm]
  \node[step=Ctrk] (n1) {GMC control\\$+$ predict\\$x'{=}Fx,\ P'{=}FPF^{\!\top}{+}Q$};
  \node[step=Ctrk, right=of n1] (n2) {innovation cov\\$S{=}MP'M^{\!\top}{+}R$};
  \node[step=Ctrk, right=of n2] (n3) {gate\\$\mathrm{IoU}\,\lor\,d^2_{\mathrm{maha}}{<}\tau$};
  \node[step=Ctrk, right=of n3] (n4) {gain\\$K{=}P'M^{\!\top}S^{-1}$};
  \node[step=Ctrk, right=of n4] (n5) {update\\$x{+}Ky,\ (I{-}KM)P'$};
  \draw[flowarr=Ctrk] (n1) -- node[flowlbl,above]{$\tilde{x},P'$} (n2);
  \draw[flowarr=Ctrk] (n2) -- node[flowlbl,above]{$S^{-1}$}       (n3);
  \draw[flowarr=Ctrk] (n3) -- node[flowlbl,above]{候選 $z_j$}     (n4);
  \draw[flowarr=Ctrk] (n4) -- node[flowlbl,above]{$y{=}z{-}Mx$}   (n5);
\end{tikzpicture}}\end{center}

\section{Prediction}\label{sec:kf-predict}
Transition 是 constant velocity,且 GMC translation 先以 control input 加到位置
(見 \Cref{ch:gmc}),再做 predict:
\begin{equation}\label{eq:kf-predict}
  F=\begin{bmatrix}I_4 & I_4\\ 0 & I_4\end{bmatrix},\qquad
  x'=Fx,\qquad P'=FPF^{\!\top}+Q(h^-).
\end{equation}
process noise 隨 predict 後框高 \(h^-\) 縮放:
\begin{equation}\label{eq:kf-Q}
  \sigma_p=\tfrac{h^-}{20},\quad \sigma_v=\tfrac{h^-}{160},\quad
  \mathrm{diag}(Q)=(\sigma_p^2,\sigma_p^2,10^{-4},\sigma_p^2,
                    \sigma_v^2,\sigma_v^2,10^{-10},\sigma_v^2).
\end{equation}

\section{Measurement Update}\label{sec:kf-update}
measurement matrix \(M=[\,I_4\ \ 0\,]\)。noise 用 predict 後 \(h^-\)(非 detection
height);baseline 下亮度/NSA 調節關閉(\(\lambda_{\mathrm{light}}{=}0,\ m_{\mathrm{NSA}}{=}1\)),
故 \(m_R=r_{\mathrm{scale}}\):
\begin{equation}\label{eq:kf-R}
  \mathrm{diag}(R)=\Big(\big(\tfrac{h^-}{20}\big)^2 m_R,\ \big(\tfrac{h^-}{20}\big)^2 m_R,\
                      10^{-2}m_R,\ \big(\tfrac{h^-}{20}\big)^2 m_R\Big).
\end{equation}
更新方程:
\begin{equation}\label{eq:kf-kalman}
\begin{aligned}
  S&=MP'M^{\!\top}+R, & K&=P'M^{\!\top}S^{-1}, & y&=z-Mx',\\
  x&\leftarrow x'+Ky, & P&\leftarrow(I-KM)P'.
\end{aligned}
\end{equation}

\section{Mahalanobis Gate}\label{sec:kf-gate}
即使 IoU 弱,只要 Mahalanobis 距離小於 \(\tau_{\mathrm{maha}}\) 仍可入選 candidate
(\(\tilde{x}_i\) 為已套 GMC control 並 predict 後的 state):
\begin{equation}\label{eq:kf-maha}
  d^2_{ij}=(z_j-M\tilde{x}_i)^{\!\top}S_i^{-1}(z_j-M\tilde{x}_i),\qquad
  \mathrm{cand}_{ij}\iff \mathrm{IoU}_{ij}>\tau_{\mathrm{iou}}\ \lor\ d^2_{ij}<\tau_{\mathrm{maha}}.
\end{equation}
IoU gate 與 cost(\Cref{ch:assoc})使用同一個 \(\tilde{x}_i\)。

\suppinput{supp/06-kalman-impl}
       % track 主體 1/3
% ===========================================================================
% track 主體 (2/3):Association Cost
% ===========================================================================
\chapter{Association Cost}\label{ch:assoc}

\begin{center}\archmap{track}\end{center}

\begin{contract}{Association cost(track 子機制)}
\textbf{輸入}\; predict 後 tracks \(\{\tilde{x}_i\}\)、detections \(\{b_j,s_j\}\)\quad
\textbf{輸出}\; 稀疏 candidate 上的成本 \(c_{ij}\)(供 auction)\\
\textbf{方法}\; IoU(+ReID)綜合質量 → 乘法式 cost,含 OAO/stability 調節\quad
\textbf{Source}\; \texttt{stage1\_cost\_fused\_kernel}(\texttt{tracker\_gpu.cu})\\
\textbf{Baseline}\; \texttt{multiplicative\_cost: true}、\texttt{sinkhorn\_lambda: 10}、
\texttt{stability\_cost\_w: 0.20}、\texttt{reid\_mode: off}
\end{contract}

\paragraph{演算法傳遞.} gate 過的 pair 先算 base quality,再匯入 penalty,經乘法式
轉成 cost,最後只把夠好的 enqueue 進稀疏 candidate list:

\begin{center}\resizebox{\linewidth}{!}{%
\begin{tikzpicture}[node distance=15mm]
  \node[step=Ctrk] (n1) {candidate gate\\$\mathrm{IoU}\,\lor\,\mathrm{Maha}$};
  \node[step=Ctrk, right=of n1] (n2) {base quality\\$A_{ij}{=}q^{\mathrm{iou}}_{ij}$};
  \node[step=Ctrk, right=of n2] (n3) {penalties\\$\Pi{=}P_{\mathrm{OAO}}{+}P_{\mathrm{vel}}{+}P_{\mathrm{occ}}{-}R_{\mathrm{stab}}$};
  \node[step=Ctrk, right=of n3] (n4) {mult.\ cost\\$c{=}1{-}A\,e^{-\Pi}$};
  \node[step=Ctrk, right=of n4] (n5) {sparse top-k\\enqueue};
  \draw[flowarr=Ctrk] (n1) -- node[flowlbl,above]{候選}        (n2);
  \draw[flowarr=Ctrk] (n2) -- node[flowlbl,above]{$A_{ij}$}     (n3);
  \draw[flowarr=Ctrk] (n3) -- node[flowlbl,above]{$\Pi_{ij}$}   (n4);
  \draw[flowarr=Ctrk] (n4) -- node[flowlbl,above]{$c_{ij}$}     (n5);
\end{tikzpicture}}\end{center}

\section{Candidate Gate 與 Base Quality}\label{sec:as-gate}
gate 同 \Cref{eq:kf-maha}。兩 gate 皆不過時 \(c_{ij}=1\)。detection score fusion:
\begin{equation}\label{eq:as-q}
  q^{\mathrm{iou}}_{ij}=\mathrm{IoU}_{ij}\bigl(1-w_{\mathrm{fuse}}(1-s_j)\bigr),
  \qquad A_{ij}=q^{\mathrm{iou}}_{ij}\ \ (\text{no ReID}).
\end{equation}
baseline \texttt{fuse\_score\_weight}=0,故 \(q^{\mathrm{iou}}_{ij}=\mathrm{IoU}_{ij}\)。
sparse list 只 enqueue 成本落在最寬鬆門檻內者:
\(c_{\max}=\max(c_{\mathrm{DDA}},\tau_{\mathrm{match}},\tau_{\mathrm{stage2}})\),
\(\mathrm{enqueue}_{ij}\iff c_{ij}\le c_{\max}\)。

\section{Multiplicative Cost}\label{sec:as-cost}
active cost form 與 penalty 匯總:
\begin{equation}\label{eq:as-cost}
  c_{ij}=\mathrm{clamp}\!\bigl(1-A_{ij}e^{-\Pi_{ij}},\,0,\,1\bigr),\qquad
  \Pi_{ij}=P_{\mathrm{OAO}}+P_{\mathrm{vel}}+P_{\mathrm{occ\_front}}-R_{\mathrm{stability}}.
\end{equation}
legacy additive path \(c_{ij}=1-A_{ij}+\sum_k P_k\) 仍存在但 baseline 不用。
baseline 只有 \(P_{\mathrm{OAO}}\) 與 \(R_{\mathrm{stability}}\) 起作用;
\(P_{\mathrm{vel}}\)(\texttt{vel\_dir\_weight})與 \(P_{\mathrm{occ\_front}}\)
(\texttt{occ\_state\_enabled})未啟用。

\section{OAO Penalty}\label{sec:as-oao}
依 predicted track-track overlap 計 occlusion 係數,並隨遮擋持續幀數 ramp:
\begin{equation}\label{eq:as-oao}
\begin{aligned}
  o^{\mathrm{base}}_i&=\max_{k\ne i}\mathrm{IoU}\!\bigl(B(x_i),B(x_k)\bigr),\qquad
  o_i=o^{\mathrm{base}}_i\min\!\Bigl(1,\tfrac{d_i}{N_{\mathrm{ramp}}}\Bigr),\\
  P_{\mathrm{OAO}}(i,j)&=\tau_{\mathrm{OAO}}\,o_i\,g_s(s_j).
\end{aligned}
\end{equation}
baseline \(\tau_{\mathrm{OAO}}=0.50\)、\(N_{\mathrm{ramp}}=25\);score 調節
\(g_s\equiv1\)(\texttt{oao\_score\_w} 未設)。

\section{Stability Reward}\label{sec:as-stab}
高度一致的 size-consistency reward(成本側,\(\lambda_{\mathrm{eff}}=\max(\lambda,1)\)):
\begin{equation}\label{eq:as-stab}
  R_{\mathrm{stability}}(i,j)=
  \frac{w_{\mathrm{stab}}/\lambda_{\mathrm{eff}}}
       {1+|h_i-h_j|/\max(h_j,10^{-3})}.
\end{equation}
因 reward 除以 \(\lambda\),進入 \(e^{-\lambda c}\)(\Cref{ch:auction})後其 boost 大致不隨
\(\lambda\) 失衡。baseline \texttt{stability\_cost\_w}=0.20。

\suppinput{supp/07-association-impl}
  % track 主體 2/3
% ===========================================================================
% track 主體 (3/3):Sparse Top-K 與 Auction Assignment
% ===========================================================================
\chapter{Sparse Top-K 與 Auction}\label{ch:auction}

\begin{center}\archmap{track}\end{center}

\begin{contract}{Auction(track 子機制)}
\textbf{輸入}\; 每個 track 的稀疏 cost candidates \(c_{ij}\)\quad
\textbf{輸出}\; 配對 \texttt{trk\_to\_det} / \texttt{det\_to\_trk}\\
\textbf{方法}\; ByteTrack 風格分數級聯 \(\times\) 單輪平行 Bertsekas auction\quad
\textbf{Source}\; \texttt{fused\_sinkhorn\_multistage\_kernel}、
\texttt{parallel\_auction\_shmem\_kernel}\\
\textbf{Baseline}\; 5-stage 級聯;\texttt{sinkhorn\_lambda} 僅作 \(e^{-\lambda c}\) value,
\emph{非}完整 Sinkhorn solve
\end{contract}

\paragraph{演算法傳遞.} cost candidates 先一次算出五個 stage 的 top-k,再依優先序
串行跑級聯;每個 stage reset 價格後跑單輪 auction 並 commit,未配掉的 track/det
流到下一 stage:

\begin{center}\resizebox{\linewidth}{!}{%
\begin{tikzpicture}[node distance=15mm]
  \node[step=Ctrk] (n1) {cost\\candidates};
  \node[step=Ctrk, right=of n1] (n2) {fused top-k\\(5 stages)};
  \node[step=Ctrk, right=of n2] (n3) {級聯\\S0$\to$S1$\to$S1b$\to$S1c$\to$S2};
  \node[step=Ctrk, right=of n3] (n4) {per-stage\\reset price $+$ auction};
  \node[step=Ctrk, right=of n4] (n5) {commit\\\texttt{trk\_to\_det}};
  \draw[flowarr=Ctrk] (n1) -- node[flowlbl,above]{$c_{ij}$} (n2);
  \draw[flowarr=Ctrk] (n2) -- node[flowlbl,above]{$p_{ij}$} (n3);
  \draw[flowarr=Ctrk] (n3) -- node[flowlbl,above]{逐 stage} (n4);
  \draw[flowarr=Ctrk] (n4) -- node[flowlbl,above]{winners}  (n5);
\end{tikzpicture}}\end{center}

\noindent
這\emph{不是} full dense Sinkhorn:每 stage 只跑一輪 bid+commit,price buffer 開頭
reset 為 0,沒有 price-raising 迭代到收斂——實質是單輪平行貪婪配對。

\section{Stage 級聯與 Value}\label{sec:au-stage}
五個 stage 依優先序貪婪,carry-over 只看前面沒配掉的 track/det:

\begin{table}[h]\centering\small\renewcommand{\arraystretch}{1.2}
\begin{tabular}{@{}llll@{}}
\toprule
Stage & Track state & Detection score & Cost cap \\
\midrule
S0 DDA       & confirmed  & \([\text{high},1.1)\)         & \(c_{\mathrm{DDA}}\) (0.12) \\
S1 high      & confirmed  & \([\text{high},1.1)\)         & \(\tau_{\mathrm{match}}\) \\
S1b mid      & confirmed  & \([\text{mid},\text{high})\)  & \(\tau_{\mathrm{match}}\) \\
S1c tentative& tentative  & \([\text{mid},1.1)\)          & \(\tau_{\mathrm{match}}\) \\
S2 low       & confirmed  & \([\text{track},\text{mid})\) & \(\tau_{\mathrm{stage2}}\) \\
\bottomrule
\end{tabular}
\caption{分數級聯(\texttt{run\_stage} 0..4)。}\label{tab:au-stage}
\end{table}

放入 top-k 的 value 與 aspect penalty:
\begin{equation}\label{eq:au-value}
  p_{ij}=e^{-\lambda c_{ij}}\,G_{\mathrm{aspect}}(b_j),\qquad
  r_j=\tfrac{\mathrm{width}_j}{\mathrm{height}_j},
\end{equation}
\begin{equation}\label{eq:au-aspect}
  G_{\mathrm{aspect}}=
  \begin{cases}
    \max(0.5,\,1-(r_j-0.8)), & r_j>0.8\\
    \max(0.5,\,1-5(0.15-r_j)), & r_j<0.15\\
    1, & \text{otherwise.}
  \end{cases}
\end{equation}

\section{Auction Bid}\label{sec:au-bid}
給定 detection 當前 price \(\rho_j\),track \(i\) 的 best/second 與 bid:
\begin{equation}\label{eq:au-bid}
  v_{ij}=p_{ij}-\rho_j,\quad
  j^\ast=\arg\max_j v_{ij},\quad
  \Delta\rho_i=v_{ij^\ast}-v^{(2)}_i+\epsilon,\quad
  \mathrm{bid}_i=\rho_{j^\ast}+\Delta\rho_i.
\end{equation}
single-round 下 \(\Delta\rho\) 只用來決定同一輪內多 track 競標同一 detection 誰勝出。
兩個 bid bias(絕對相加):
\begin{equation}\label{eq:au-bias}
  \mathrm{bid}_i \mathrel{+}=\frac{w_{\mathrm{fresh}}}{1+\mathrm{age}_i},\qquad
  \mathrm{bid}_i \mathrel{+}=\frac{w_{\mathrm{stab,bid}}}{1+|h_i-h_j|/h_j}.
\end{equation}
預設不一致,\textbf{勿一律當 off}:freshness \(w_{\mathrm{fresh}}\)
(\texttt{SACCADE\_FRESHNESS\_W})預設 0(關);stability bid bias
\(w_{\mathrm{stab,bid}}\)(\texttt{SACCADE\_STABILITY\_W})預設 \(0.1\)(\textbf{開})。

\suppinput{supp/08-auction-impl}
      % track 主體 3/3
% ===========================================================================
% track 收尾 (1/2):Track Lifecycle 與 Birth
% ===========================================================================
\chapter{Track Lifecycle 與 Birth}\label{ch:lifecycle}

\begin{center}\archmap{track}\end{center}

\begin{contract}{Lifecycle(track 子機制)}
\textbf{輸入}\; assignment 結果(matched / unmatched track / unmatched detection)\quad
\textbf{輸出}\; 更新後的 track 集合與狀態(tentative/confirmed/lost/removed)\\
\textbf{方法}\; hit-streak 累積 confirm、age 計數 lost/remove\quad
\textbf{Source}\; \texttt{tracker\_gpu.cu}\\
\textbf{Baseline}\; \texttt{new\_track\_thresh: 0.28}、\texttt{confirm\_streak: 3}、
\texttt{confirm\_score\_thresh: 0.50}、\texttt{track\_buffer: 30}、\texttt{per\_seq\_adapt: false}
\end{contract}

\paragraph{演算法傳遞.} assignment 後依配對狀態驅動 track 狀態機;下圖為主要轉移
(tentative 若連續未配會很快 drop,圖中略):

\begin{center}\resizebox{\linewidth}{!}{%
\begin{tikzpicture}[node distance=20mm]
  \node[step=Ctrk] (nw) {未配 det};
  \node[step=Ctrk, right=of nw] (te) {tentative};
  \node[step=Ctrk, right=of te] (co) {confirmed};
  \node[step=Ctrk, right=of co] (lo) {lost};
  \node[step=Ctrk, right=of lo] (rm) {removed};
  \draw[flowarr=Ctrk] (nw) -- node[flowlbl,above]{$s_j{\ge}$new\_thr} (te);
  \draw[flowarr=Ctrk] (te) -- node[flowlbl,above]{streak${\ge}3$}    (co);
  \draw[flowarr=Ctrk] (co) -- node[flowlbl,above]{unmatched}         (lo);
  \draw[flowarr=Ctrk] (lo) -- node[flowlbl,above]{age${>}$buffer}    (rm);
  \draw[flowarr=Ctrk] (lo) to[bend left=32] node[flowlbl,above]{re-match} (co);
\end{tikzpicture}}\end{center}

\section{State 轉移}\label{sec:lc-states}
assignment 後 tracker 更新 state:
\begin{itemize}\setlength\itemsep{2pt}
  \item \textbf{matched}(confirmed/tentative):以 detection \(z_j\) 做 Kalman update
        (\Cref{ch:kalman}),\texttt{hit\_streak}\(+1\),\texttt{age}/\texttt{time\_since\_update}
        歸零。
  \item \textbf{unmatched active track}:\texttt{age}\(\le\)\texttt{track\_buffer} 期間維持
        active/lost,超過 max age 後 remove。
  \item \textbf{unmatched detection}(\(s_j\ge\)\texttt{new\_track\_thresh}):生成新的
        tentative track。
\end{itemize}

\section{Birth 與 Confirm}\label{sec:lc-confirm}
新 tentative track 必須累積足夠連續 evidence 才會被視為 confirmed output:
\begin{equation}\label{eq:lc-confirm}
\begin{aligned}
  \text{confirmed}\iff{}&
  \texttt{hit\_streak}\ge\texttt{confirm\_streak}\\
  &{}\land\bar{s}\ge\texttt{confirm\_score\_thresh}.
\end{aligned}
\end{equation}
baseline:
\begin{center}\small
\begin{tabular}{@{}ll@{\qquad}ll@{}}
\texttt{confirm\_streak} & 3 &
\texttt{confirm\_score\_thresh} & 0.50 \\
\texttt{new\_track\_thresh} & 0.28 &
\texttt{track\_buffer} & 30
\end{tabular}
\end{center}
若干 birth-gate experiments 位於 \path{scripts/eval/config/lifecycle.py},但目前
preset 未啟用(\texttt{per\_seq\_adapt: false})。
普通新生 track 以 tentative 與 \texttt{hit\_streak=1} 起跑;bridge relink 復活的
slot 會保留 lost id 並直接回到 confirmed,不再走 tentative confirm gate。

\suppinput{supp/09-lifecycle-impl}
    % track 收尾 1/2
% ===========================================================================
% track 收尾 (2/2):Bridge Relink
% ===========================================================================
\chapter{Bridge Relink}\label{ch:relink}

\begin{center}\archmap{track}\end{center}

\begin{contract}{Bridge relink(track 子機制)}
\textbf{輸入}\; 新穩定的 candidate track、窗內未配的 lost confirmed track\quad
\textbf{輸出}\; candidate 採用 lost track 的 id(lost slot deactivate)\\
\textbf{方法}\; 雙向中點外推(bidirectional midpoint extrapolation),\emph{非} appearance ReID\quad
\textbf{Source}\; \texttt{tracker\_gpu.cu}(\(\ne\) \texttt{relink\_gate.cu} 的 appearance gate)\\
\textbf{Baseline}\; \texttt{relink\_bridge\_enabled: true}、\texttt{relink\_bridge\_px: 0.25}、
\texttt{margin: 0.05}、\texttt{h\_lo/h\_hi: 0.75/1.33}、\texttt{dir\_bonus: 0.8}、
anchor \texttt{adaptive}
\end{contract}

\paragraph{演算法傳遞.} 對 candidate–lost pair,先取 anchor、做 4 點速度回歸,再雙向
外推算殘差,合成 \(d_{\mathrm{bridge}}\),經 direction bonus 與多重 gate 後 commit:

\begin{center}\resizebox{\linewidth}{!}{%
\begin{tikzpicture}[node distance=12mm]
  \node[step=Ctrk] (n1) {candidate\\$+$ lost pair};
  \node[step=Ctrk, right=of n1] (n2) {anchor\\$(a_x,a_y)$};
  \node[step=Ctrk, right=of n2] (n3) {4-pt velocity\\$v_\ell,v_c$};
  \node[step=Ctrk, right=of n3] (n4) {外推殘差\\$r_{\mathrm{fwd}},r_{\mathrm{bwd}}$};
  \node[step=Ctrk, right=of n4] (n5) {$d_{\mathrm{bridge}}$\\$+$ dir bonus};
  \node[step=Ctrk, right=of n5] (n6) {gates\\$+$ commit};
  \draw[flowarr=Ctrk] (n1) -- (n2);
  \draw[flowarr=Ctrk] (n2) -- node[flowlbl,above]{$\ell,c$} (n3);
  \draw[flowarr=Ctrk] (n3) -- (n4);
  \draw[flowarr=Ctrk] (n4) -- node[flowlbl,above]{$d_h$} (n5);
  \draw[flowarr=Ctrk] (n5) -- node[flowlbl,above]{$d^{(1)}$} (n6);
\end{tikzpicture}}\end{center}

\section{History 與 Anchor}\label{sec:rl-anchor}
每個 observed slot(\texttt{age}=0)保存 center/height history ring,並對框高做 EMA:
\(\bar{h}\leftarrow0.95\bar{h}+0.05h\)。coasting lost track 保留最後 history。
baseline anchor \texttt{adaptive}:\(x\) 永遠用 center;\(y\) 候選為
\(y^{\mathrm{top}}{=}c_y{-}h/2,\ y^{\mathrm{bot}}{=}c_y{+}h/2,\ y^{\mathrm{ctr}}{=}c_y\)。
若平均相鄰 height 變化 \(\le\)\texttt{anchor\_rate}(0.03)退化為 center;否則以 top/bottom
各自四點線性回歸殘差為權重選較穩的 edge:
\begin{equation}\label{eq:rl-anchor}
  w_e=\frac{1}{\mathrm{RSS}_{\mathrm{line}}(y^e)/(\bar{h}^2+10^{-3})+0.01},\quad
  a_y=\frac{w_{\mathrm{top}}y^{\mathrm{top}}+w_{\mathrm{bot}}y^{\mathrm{bot}}}
           {w_{\mathrm{top}}+w_{\mathrm{bot}}}.
\end{equation}

\section{Candidate 條件}\label{sec:rl-cand}
candidate 採用 lost id 的前置條件:
\begin{center}\ttfamily\small
candidate hit\_streak == bridge\_at\quad candidate 有 $\ge$4 history samples\\
lost active 但本幀未配\quad lost state == confirmed\quad
bridge\_min\_lost $\le$ lost\_age $\le$ bridge\_ttl
\end{center}

\section{速度回歸與 \texorpdfstring{\(d_{\mathrm{bridge}}\)}{d\_bridge}}\label{sec:rl-bridge}
對四個等時距 scalar anchor sample \((p_0,p_1,p_2,p_3)\),4 點回歸速度
\(v=\tfrac{3p_3+p_2-p_1-3p_0}{10}\)。令 \(\ell\)=lost exit anchor、\(c\)=candidate
entry anchor、\(g\)=\texttt{lost\_age}、\(h_{\mathrm{ref}}=\max(\tfrac{\bar{h}_\ell+\bar{h}_c}{2},1)\):
\begin{equation}\label{eq:rl-resid}
  r_{\mathrm{fwd}}=\frac{\lVert(\ell+v_\ell g)-c\rVert}{h_{\mathrm{ref}}},\qquad
  r_{\mathrm{bwd}}=\frac{\lVert(c-v_c g)-\ell\rVert}{h_{\mathrm{ref}}}.
\end{equation}
以 lost 端速度大小決定外推 vs 純距離的權重:
\begin{align}\label{eq:rl-dbridge}
  d_h&=\frac{\lVert\ell-c\rVert}{h_{\mathrm{ref}}},\\
  w&=\sqrt{\mathrm{clamp}\!\Bigl(\tfrac{\lVert v_\ell\rVert/h_{\mathrm{ref}}}{0.12},0,1\Bigr)},\\
  d_{\mathrm{bridge}}&=
  w\,\frac{r_{\mathrm{fwd}}+r_{\mathrm{bwd}}}{2}+(1-w)\,d_h.
\end{align}

\section{Direction Bonus}\label{sec:rl-dir}
\texttt{dir\_bonus}>0 且方向相近(\(\cos\theta>0.5\))時,向 cross-track 誤差 \(d_{\mathrm{cross}}\) blend:
\begin{equation}\label{eq:rl-dir}
  \alpha=\min\!\bigl(w_{\mathrm{dir}}\cos^2\theta\,\eta_v\,\eta_g,\,1\bigr),\qquad
  d_{\mathrm{bridge}}\leftarrow(1-\alpha)d_{\mathrm{bridge}}+\alpha\,d_{\mathrm{cross}},
\end{equation}
其中 \(\eta_v=\mathrm{clamp}\bigl(\tfrac{\min(\lVert v_\ell\rVert,\lVert v_c\rVert)}
{\max(0.005h_{\mathrm{ref}},10^{-3})},0,1\bigr)\)、\(\eta_g=\mathrm{clamp}(g/30,0,1)\)。
baseline \texttt{dir\_bonus}=0.8。

\section{Gates 與 Commit}\label{sec:rl-gate}
接受條件:
\begin{equation}\label{eq:rl-gate}
  d_{\mathrm{bridge}}\le\tau_{\mathrm{bridge}},\quad
  \frac{\bar{h}_\ell}{\bar{h}_c}\in[h_{\mathrm{lo}},h_{\mathrm{hi}}],\quad
  d^{(2)}_{\mathrm{bridge}}-d^{(1)}_{\mathrm{bridge}}\ge m_{\mathrm{bridge}}.
\end{equation}
注意 \texttt{relink\_bridge\_px}=0.25 是歷史命名,\emph{非} 0.25 pixel——它與已除以
\(h_{\mathrm{ref}}\) 的 \(d_{\mathrm{bridge}}\) 比較,單位是 reference-height-normalized
distance。baseline \texttt{spatial\_gate}=0.0(關),不啟用 occupancy expansion。
接受後 candidate 採用 lost id、lost slot deactivate;多 candidate 競爭同一 lost id 時,
較高 detection score 贏,candidate index 作 tie-breaker。

\suppinput{supp/10-relink-impl}
       % track 收尾 2/2

% ---- 附錄:reference 性質的大表(符號↔config)放這,不擋正文節奏 ----
\appendix
% ===========================================================================
% 附錄 A:符號表 與 符號↔config 對照(math_model §3 / §3.1 / §3.2)
% reference 性質,故置於附錄。代碼錨點(cu:N 行號)見各章 L3 實作補充(full 版)。
% ===========================================================================
\chapter{符號與 Config 對照}\label{app:symbols}

baseline 值對應 \texttt{mamba\_whole\_graph}(frozen\_v2)。各機制的代碼錨點
(kernel 名、\texttt{cu:N} 行號)見對應章末的 L3 實作補充。

\section{基本符號}\label{app:basic}

\begin{longtable}{@{}>{\raggedright\arraybackslash}p{3.3cm}
  >{\raggedright\arraybackslash}p{8.2cm}@{}}
\toprule
符號 & 意義 \\* \midrule \endhead
\(f\) & 目前 frame index \\
\(I_f\) & 目前 RGB/CHW frame tensor \\
\(D_f=\{d_j\}\) & frame \(f\) postprocess 後的 detection set \\
\(b_j=(x_1,y_1,x_2,y_2)\) & detection box(original-frame pixels) \\
\(s_j,\ \mathrm{cls}_j\) & detection score / class id \\
\(T_f=\{t_i\}\) & association 前 active 或 lost tracker slots \\
\(x_i,\ P_i\) & track slot \(i\) 的 8D Kalman state / \(8\times8\) covariance \\
\(W_f\) & GMC \(2\times3\) camera warp;translation path \([1,0,t_x;0,1,t_y]\) \\
\midrule
\(x=(c_x,c_y,a,h,\dot{\cdot})\) & center、aspect、height 與對應速度 \\
\(w=a\cdot h,\ B(x)\) & state implied width / 還原的 box \\
\(z=(c_x,c_y,a,h),\ M=[I_4\ 0]\) & Kalman measurement / measurement matrix \\
\(\mathrm{IoU}(i,j),\ c_{ij}\) & \(B(x_i)\) 與 \(b_j\) 的 IoU / final association cost \\
\(p_{ij},\ A_{ij},\ \Pi_{ij}\) & auction value / base quality / penalty 總和 \\
\bottomrule
\end{longtable}

% 共用:4 欄 config 對照表的欄型(footnotesize,長名稱用 \newline 斷行)
\section{GMC}\label{app:gmc}
{\footnotesize\setlength{\tabcolsep}{3pt}
\begin{longtable}{@{}>{\raggedright\arraybackslash}p{2.8cm}
                     >{\raggedright\arraybackslash}p{2.4cm}
                     >{\raggedright\arraybackslash}p{3.8cm}
                     >{\raggedright\arraybackslash}p{1.6cm}@{}}
\toprule
符號 & 用途 & config / env & baseline \\* \midrule \endhead
\(W_f\) & 相機運動補償 warp & — & trans-only \\
PCR & phase-corr 峰值可信度 & — & — \\
\(\tau_{\mathrm{PCR}}\) & 低可信度縮小位移 & \texttt{SACCADE\_GMC\_PCR\_THRESH} & py \(5.0\) \\
\bottomrule
\end{longtable}}

\section{Kalman}\label{app:kalman}
{\footnotesize\setlength{\tabcolsep}{4pt}
\begin{longtable}{@{}>{\raggedright\arraybackslash}p{1.8cm}
                     >{\raggedright\arraybackslash}p{3.0cm}
                     >{\raggedright\arraybackslash}p{4.2cm}
                     >{\raggedright\arraybackslash}p{1.7cm}@{}}
\toprule
符號 & 用途 & config / env & baseline \\* \midrule \endhead
\(x,z,F,M\) & 8D CV state / 4D measurement & — & — \\
\(\sigma_p{=}h^-/20\) & 位置過程噪聲 & hardcoded & 1/20 \\
\(\sigma_v{=}h^-/160\) & 速度過程噪聲 & hardcoded & 1/160 \\
\(r_{\mathrm{scale}}\) & 測量噪聲縮放 & \texttt{kalman\_r\_scale} & 2.8 \\
\(m_{\mathrm{NSA}},\lambda_{\mathrm{light}}\) & NSA/亮度調節 & — & 1 / 0 \\
\(\tau_{\mathrm{maha}}\) & Mahalanobis gate & \texttt{maha\_gate} & 見實作 \\
\bottomrule
\end{longtable}}

\section{Association Cost}\label{app:assoc}
{\footnotesize\setlength{\tabcolsep}{4pt}
\begin{longtable}{@{}>{\raggedright\arraybackslash}p{1.8cm}
                     >{\raggedright\arraybackslash}p{3.0cm}
                     >{\raggedright\arraybackslash}p{4.2cm}
                     >{\raggedright\arraybackslash}p{1.7cm}@{}}
\toprule
符號 & 用途 & config / env & baseline \\* \midrule \endhead
\(A_{ij}\) & IoU(+ReID)綜合質量 & — & = IoU \\
\(w_{\mathrm{fuse}}\) & 低分檢測降權 & \texttt{fuse\_score\_weight} & 0.0 \\
\(c_{ij}\) & 乘法式 cost & \texttt{multiplicative\_cost} & true \\
\(\lambda\) & cost\(\to\)value 溫度 & \texttt{sinkhorn\_lambda} & 10 \\
\(o_i,\tau_{\mathrm{OAO}}\) & 遮擋配對抑制 & \texttt{oao\_tau} /\newline \texttt{oao\_ramp\_frames} & 0.50 / 25 \\
\(w_{\mathrm{vel}}\) & 速度反向懲罰 & \texttt{vel\_dir\_weight} & 關 \\
\(w_{\mathrm{occ}}\) & front-occluder 懲罰 & \texttt{occ\_cost\_weight} & 關 \\
\(w_{\mathrm{stab}}\) & 高度一致 reward(成本側) & \texttt{stability\_cost\_w} & 0.20 \\
\bottomrule
\end{longtable}}

\section{Auction}\label{app:auction}
{\footnotesize\setlength{\tabcolsep}{4pt}
\begin{longtable}{@{}>{\raggedright\arraybackslash}p{1.8cm}
                     >{\raggedright\arraybackslash}p{3.0cm}
                     >{\raggedright\arraybackslash}p{4.2cm}
                     >{\raggedright\arraybackslash}p{1.7cm}@{}}
\toprule
符號 & 用途 & config / env & baseline \\* \midrule \endhead
\(p_{ij}\) & auction 概率值 & — & — \\
\(G_{\mathrm{aspect}}\) & 抑制異常長寬比 & hardcoded (0.8/0.15) & — \\
\(w_{\mathrm{fresh}}\) & 新鮮度 bid bias & \texttt{SACCADE\_FRESHNESS\_W} & 0.0 \\
\(w_{\mathrm{stab,bid}}\) & 高度一致 bid bias & \texttt{SACCADE\_STABILITY\_W} & 0.1(開) \\
S0 DDA & confirmed×high 更緊 stage & \texttt{SACCADE\_ENABLE\_DDA} /\newline \texttt{\_DDA\_MAX\_COST} & on / 0.12 \\
stage thr & 分數級聯邊界 & \texttt{match} / \texttt{high} / \texttt{mid} /\newline \texttt{track} / \texttt{stage2} & .50/.45/\newline.10/.05/.50 \\
\bottomrule
\end{longtable}}

\section{Lifecycle / Bridge / Semantic relink}\label{app:rest}
{\footnotesize\setlength{\tabcolsep}{4pt}
\begin{longtable}{@{}>{\raggedright\arraybackslash}p{1.8cm}
                     >{\raggedright\arraybackslash}p{3.0cm}
                     >{\raggedright\arraybackslash}p{4.2cm}
                     >{\raggedright\arraybackslash}p{1.7cm}@{}}
\toprule
符號 & 用途 & config / env & baseline \\* \midrule \endhead
birth/\newline confirm & tentative\(\to\)confirmed & \texttt{new\_track\_thresh} /\newline \texttt{confirm\_streak} & 0.28 / 3 \\
\(\tau_{\mathrm{bridge}}\) & 雙向外推殘差門檻 & \texttt{relink\_bridge\_px} & 0.25 \\
\(\alpha\) & 方向一致偏移 & \texttt{relink\_bridge\_dir\_bonus} & 0.8 \\
\(h_{\mathrm{lo}},h_{\mathrm{hi}}\) & 高度比 gate & \texttt{relink\_bridge\_h\_lo} /\newline \texttt{\_h\_hi} & 0.75/1.33 \\
\(m_{\mathrm{bridge}}\) & best-vs-second margin & \texttt{relink\_bridge\_margin} & 0.05 \\
\(w_{\mathrm{sim/iou}}\),\newline \(w_{\mathrm{maha}}\) & semantic joint 權重 & \texttt{semantic\_w\_*\_base} & off \\
\bottomrule
\end{longtable}}

\section{方法出處 / 命名對照}\label{app:naming}
\begin{itemize}\setlength\itemsep{2pt}
  \item \textbf{GMC} = phase correlation(cross-power spectrum + Hanning window),translation-only warp。
  \item \textbf{Kalman} = SORT/DeepSORT 風格 constant-velocity filter。
  \item \textbf{Assignment} = Bertsekas auction(單輪平行貪婪)跑在 softmin-temperature top-k 上;
        分數分段 = ByteTrack 風格 cascade。\texttt{sinkhorn\_lambda} 為歷史命名——只用
        \(e^{-\lambda c}\) 當 value,\textbf{非}完整 Sinkhorn 迭代。
  \item \textbf{OAO} = occlusion-aware(track-track overlap)配對抑制 + duration ramp。
  \item \textbf{Bridge relink} = 雙向中點外推,項目自有機制,非標準 appearance ReID。
  \item \textbf{Semantic relink gate} = appearance + Mahalanobis + IoU joint gate(baseline 關)。
\end{itemize}

\paragraph{提醒.} \(w_{\mathrm{stab}}\)(\Cref{sec:as-stab},成本側 reward)與
\(w_{\mathrm{stab,bid}}\)(\Cref{sec:au-bid},auction bid bias)是\textbf{兩個不同旋鈕},
數值與作用點皆不同,雖都用高度一致性 \(|h_i-h_j|/h_j\)。


\end{document}
"
% ===========================================================================
\documentclass[11pt,oneside]{ctexbook}   % 中文 → ctexbook + XeLaTeX
% ===========================================================================
% preamble.tex —— 套件、深度開關、archmap 迷你地圖、合約框
% 引擎:XeLaTeX(中文必須)。建置見 README.md / Makefile。
% ===========================================================================

% --- 繁體中文字型(fandol 為簡體取向,缺繁體字形;改用 Noto CJK TC)---
\setCJKmainfont{Noto Serif CJK TC}
\setCJKsansfont{Noto Sans CJK TC}
\setCJKmonofont{Noto Sans CJK TC}
% ctexbook 預設 caption 是簡體;改回繁體(录→錄)。appendixname 同時管標題與頁首。
\ctexset{contentsname=目錄, appendixname=附錄}
% 中英混排時,允許行末額外伸縮以吸收邊際 overfull(標準作法)
\setlength{\emergencystretch}{3em}

% --- 數學 / 表格 / 引用 ---
\usepackage{amsmath, amssymb, mathtools}
\usepackage{booktabs}            % \toprule \midrule \bottomrule,禁直線
\usepackage{array}
\usepackage{tabularx}            % 自動寬度欄(職責表用)
\usepackage{longtable}           % 可跨頁的參考表(附錄符號表用)
\usepackage[hidelinks]{hyperref}
\usepackage{cleveref}            % \Cref{eq:...} 就近引用
\usepackage{xstring}             % archmap 的字串比較(跨引擎安全)
\usepackage{fancyhdr}            % 統一頁碼位置

% --- 圖 ---
\usepackage{tikz}
\usetikzlibrary{arrows.meta, positioning, fit, backgrounds}

% --- 程式碼 / 合約框 ---
\usepackage[most]{tcolorbox}
\usepackage{fvextra}             % Verbatim,ASCII pipeline 圖暫存用

% --- 頁碼 ---
% ctexbook/book 的一般頁與章首頁預設 page style 不同,會讓頁碼在頁首/頁腳間跳。
% 這裡同時覆寫一般頁型與 plain 頁型,固定頁碼在頁腳中央。
\fancypagestyle{saccadepage}{
  \fancyhf{}
  \fancyfoot[C]{\thepage}
  \renewcommand{\headrulewidth}{0pt}
  \renewcommand{\footrulewidth}{0pt}
}
\fancypagestyle{plain}{
  \fancyhf{}
  \fancyfoot[C]{\thepage}
  \renewcommand{\headrulewidth}{0pt}
  \renewcommand{\footrulewidth}{0pt}
}
\AtBeginDocument{\pagestyle{saccadepage}}

% ===========================================================================
% 深度開關:L3 實作補充是否編入
%   make full  -> \def\IMPL{1}  含 supp/*-impl.tex
%   make paper -> \def\IMPL{0}  剝掉所有 L3,得乾淨學術版
% ===========================================================================
\providecommand{\IMPL}{1}
\newif\ifimpl
\ifnum\IMPL=1 \impltrue \else \implfalse \fi

% \suppinput{file}:只有 full 版會把 L3 補充檔讀進來
\newcommand{\suppinput}[1]{\ifimpl\input{#1}\fi}

% L3 補充章節的統一外框(視覺上和學術本體區隔)
\newtcolorbox{implsupp}[1]{
  enhanced, breakable,
  colback=orange!4, colframe=orange!55!black,
  fonttitle=\bfseries\small, coltitle=orange!30!black,
  title={實作補充 (L3,僅 full 版)\quad#1},
  left=2mm, right=2mm, top=1mm, bottom=1mm}

% ===========================================================================
% 模組合約框:每個細節章開頭,濃縮 I/O / Source / Baseline
% ===========================================================================
\newtcolorbox{contract}[1]{
  enhanced, breakable,
  colback=gray!4, colframe=black!55,
  fonttitle=\bfseries, title={模組合約:#1}}

% ===========================================================================
% \archmap{focus} —— 全書共用的「你在這裡」迷你地圖
%   focus ∈ {none, detect, gmc, track, output}
%   none = 全亮(總覽章用);其餘 = 點亮焦點、淡化其它
%   track 內的子章(Kalman/association/auction/...)先一律用 track;
%   日後要更細,可在此複製出 \trackmap 畫 track 內部子階。
% ===========================================================================
% 每個 stage 的代表色(改這裡就改全書配色)——用 \colorlet 建具名顏色,
% 才能在 pgfkeys option 裡以字面名稱引用(巨集名不會被展開)。
\colorlet{Cbuf}{teal}
\colorlet{Cdet}{blue!65!cyan}
\colorlet{Cgmc}{orange!90!red}
\colorlet{Ctrk}{violet!80!blue}
\colorlet{Cout}{green!55!black}

% 全域 style:box 外型 + 三種狀態(吃顏色當 #1)+ 箭頭/標籤
\tikzset{
  ambox/.style  ={draw, rounded corners, align=center,
                  minimum height=10mm, minimum width=24mm, font=\small},
  amfull/.style ={fill=#1!25, draw=#1!85!black, line width=1pt,   text=black!88},
  amfade/.style ={fill=#1!8,  draw=#1!35,       line width=0.6pt, text=black!40},
  amfocus/.style={fill=#1!48, draw=#1!90!black, line width=1.9pt, text=black!92},
  amarr/.style  ={-{Stealth}, thick},               % 直角:不加 rounded corners
  amlbl/.style  ={font=\scriptsize, text=black!55},
}

\newcommand{\archmap}[1]{%
  % 預設 full(總覽 none 全部上色);具名 style 都是字面,pgfkeys 吃得到
  \tikzset{amSbuf/.style={amfull=Cbuf}, amSdet/.style={amfull=Cdet},
           amSgmc/.style={amfull=Cgmc}, amStrk/.style={amfull=Ctrk},
           amSout/.style={amfull=Cout}}%
  \IfStrEq{#1}{none}{}{% 指定焦點:先全 fade,再把焦點改 focus
    \tikzset{amSbuf/.style={amfade=Cbuf}, amSdet/.style={amfade=Cdet},
             amSgmc/.style={amfade=Cgmc}, amStrk/.style={amfade=Ctrk},
             amSout/.style={amfade=Cout}}%
    \IfStrEq{#1}{detect}{\tikzset{amSdet/.style={amfocus=Cdet}}}{}%
    \IfStrEq{#1}{gmc}{\tikzset{amSgmc/.style={amfocus=Cgmc}}}{}%
    \IfStrEq{#1}{track}{\tikzset{amStrk/.style={amfocus=Ctrk}}}{}%
    \IfStrEq{#1}{output}{\tikzset{amSout/.style={amfocus=Cout}}}{}%
  }%
  % 固定成欄寬比例再置中,避免視覺偏移
  \resizebox{0.92\linewidth}{!}{%
  \begin{tikzpicture}
    % 橫式佈局:左→右流,detect/GMC 為上下兩條平行支線
    \node[ambox, amSbuf] (buf) at (0,0)     {frame buffer\\\scriptsize(GPU, CHW)};
    \node[ambox, amSdet] (det) at (3.2, 1)  {detect 鏈};
    \node[ambox, amSgmc] (gmc) at (3.2,-1)  {GMC};
    \node[ambox, amStrk] (trk) at (6.6, 0)  {track\\\scriptsize tracker\_gpu.cu};
    \node[ambox, amSout] (out) at (9.8, 0)  {materialize\\\scriptsize→ MOT rows};
    % 分岔:buf 拉出 bus,再直角接上下兩支線(連線染來源色)
    \draw[amarr, Cbuf!65!black] (buf.east) -- ++(0.4,0) |- (det.west);
    \draw[amarr, Cbuf!65!black] (buf.east) -- ++(0.4,0) |- (gmc.west);
    % 匯合:兩支線直角收進 track
    \draw[amarr, Cdet!65!black] (det.east) -- ++(0.4,0) node[amlbl,above]{boxes} |- (trk.west);
    \draw[amarr, Cgmc!70!black] (gmc.east) -- ++(0.4,0) node[amlbl,below]{$W_f$} |- (trk.west);
    \draw[amarr, Ctrk!65!black] (trk.east) -- (out.west);
  \end{tikzpicture}}}

% ===========================================================================
% 章內「演算法傳遞圖」共用 style —— 每章自畫,顏色用該模組的 C*
%   step=<color>     一個演算法步驟方塊
%   flowarr=<color>  傳遞箭頭(直角時用 |- / -|)
%   flowlbl          標在邊上的「傳了什麼」(白底壓線)
% 用法見 chapters/05-gmc.tex 的範例。
% ===========================================================================
\tikzset{
  step/.style   ={draw, rounded corners=2pt, align=center, font=\footnotesize,
                  minimum height=10mm, minimum width=20mm, inner sep=4pt,
                  fill=#1!14, draw=#1!70!black, text=black!88},
  flowarr/.style={-{Stealth}, semithick, draw=#1!75!black},
  flowlbl/.style={font=\scriptsize, text=black!60, fill=white, inner sep=1.5pt},
}
% 把 GMC 章內流程一行寫完的便利指令(其他章可仿照自行展開)
\newcommand{\gmcflow}{%
  \begin{center}\resizebox{\linewidth}{!}{%
  \begin{tikzpicture}[node distance=13mm]
    \node[step=Cgmc] (n1) {灰階降採樣\\+\,Hanning};
    \node[step=Cgmc, right=of n1] (n2) {FFT\\prev / curr};
    \node[step=Cgmc, right=of n2] (n3) {cross-power\\$R=\overline{A}B/|{\cdot}|$};
    \node[step=Cgmc, right=of n3] (n4) {IFFT $\to$ peak\\PCR};
    \node[step=Cgmc, right=of n4] (n5) {confidence $\gamma$\\$\to$ warp $W_f$};
    \draw[flowarr=Cgmc] (n1) -- node[flowlbl,above]{$G_w$}        (n2);
    \draw[flowarr=Cgmc] (n2) -- node[flowlbl,above]{$A,B$}        (n3);
    \draw[flowarr=Cgmc] (n3) -- node[flowlbl,above]{$R(k)$}       (n4);
    \draw[flowarr=Cgmc] (n4) -- node[flowlbl,above]{$r$ peak}     (n5);
  \end{tikzpicture}}\end{center}}


\title{Saccade 全局數學模型}
\author{Ray Lei}

\begin{document}
% Keep page numbering format uniform across title, TOC, main text, and appendix.
\maketitle
\setcounter{page}{2}
\tableofcontents

% ---- Part I:主架構(the spine,只放地圖,不放數學)----
\part{系統總覽}
% ===========================================================================
% Part I 唯一一章:系統總覽 = 主架構(the spine)
% 原則:這裡只放「地圖」——架構圖、模組邊界、dataflow、baseline 合約。
%       所有數學一律下放到 Part II。這就是「不全展開」的關鍵。
% ===========================================================================
\chapter{系統總覽}\label{ch:overview}

整條 pipeline 分兩層:\textbf{Python orchestration}(排程、buffer 管理、
CUDA graph replay、eval)與 \textbf{C++/CUDA compute}(detector、postprocess、
GMC、tracker association)。Python 不在 hot path 做數值計算,只在 output
boundary 把結果搬回 host。

關鍵:detect 與 GMC 是從\textbf{同一個 frame buffer 分岔的兩條平行支線}——
GMC 吃 frame 灰階(+上一幀),\emph{不是} detect 的輸出;\texttt{track} 才把
兩條支線匯合。

\section{架構地圖}\label{sec:archmap}
全書各章開頭都會重畫這張圖、點亮當前模組(「你在這裡」)。總覽版全亮:

\begin{center}
  \archmap{none}
\end{center}

\section{模組職責與 I/O}\label{sec:modules}
下表概括每個模組\textbf{吃什麼、吐什麼、解決什麼問題},以及在 pipeline 中銜接誰。
逐項數學見 Part II 對應章;傳遞合約與 source 見 \Cref{tab:transfer}。

\begin{table}[h]\small\setlength{\tabcolsep}{4pt}\renewcommand{\arraystretch}{1.35}
\begin{tabularx}{\linewidth}{@{}l >{\raggedright\arraybackslash}X >{\raggedright\arraybackslash}X l@{}}
\toprule
模組 & 輸入 → 輸出 & 解決的問題 & 銜接 \\
\midrule
detect      & frame CHW tensor → raw boxes/scores/classes
            & 從單張影像偵測行人(TRT backbone + Mamba head,anchor-free decode)
            & → postprocess \\
postprocess & raw det → 過濾/NMS 後 boxes
            & 去除低分與重疊框,交給 tracker 乾淨的偵測集
            & detect → track \\
GMC         & prev+curr 灰階 → \(2\times3\) warp \(W_f\)
            & 估相機運動補償平移,讓 Kalman 的 velocity 只學物體運動
            & 與 detect 平行,補 track \\
track       & boxes/scores/\(W_f\) → track states / IDs
            & 跨幀關聯與維持 ID(分解見 \Cref{tab:track-internal})
            & 匯合 → materialize \\
output      & GPU state → MOT rows
            & 搬回 host、產生 MOT 結果(含 tracklet 內插後處理)
            & track → 終點 \\
\bottomrule
\end{tabularx}
\caption{模組職責與輸入/輸出概覽。}\label{tab:modules}
\end{table}

\texttt{track} 不是單一步驟,而由數個子機制串成,各為 Part II 的一章:

\begin{table}[h]\small\setlength{\tabcolsep}{4pt}\renewcommand{\arraystretch}{1.3}
\begin{tabularx}{\linewidth}{@{}l >{\raggedright\arraybackslash}X@{}}
\toprule
子機制 & 解決的問題 \\
\midrule
Kalman 模型      & constant-velocity 運動模型;每幀 predict 再用 detection update track state \\
Association cost & 把 IoU、OAO/velocity penalty、stability reward 匯成單一配對成本 \(c_{ij}\) \\
Auction 指派     & ByteTrack 風格分數級聯 + 單輪平行 auction,決定 track 與 detection 配對 \\
Track lifecycle  & track 生死:tentative→confirmed 的 birth/confirm,與 lost/remove \\
Bridge relink    & 找回短暫消失又重生的 ID(雙向中點外推,非 appearance ReID) \\
\bottomrule
\end{tabularx}
\caption{\texttt{track} 內部子機制(Part II 逐章展開)。}\label{tab:track-internal}
\end{table}

傳遞合約(誰留在 GPU、source 檔)另列:

\begin{table}[h]\small\setlength{\tabcolsep}{4pt}\renewcommand{\arraystretch}{1.2}
\begin{tabularx}{\linewidth}{@{}l
  >{\raggedright\arraybackslash}X
  >{\raggedright\arraybackslash}X@{}}
\toprule
模組 & 模組(Python facade / C++·CUDA core) & 傳遞方式 \\
\midrule
detect      & \texttt{mamba\_gated\_detector.py} / TRT backbone + Mamba head & GPU tensor(whole-detect graph)\\
postprocess & \texttt{PerceptionPipeline} / \texttt{pipeline.cpp} & GPU buffer,留 device \\
GMC         & \texttt{GMC} / \texttt{gmc\_kernel.cu} (cuFFT)       & GPU tensor,$2\times3$ warp \\
track       & \texttt{GPUByteTracker} / \texttt{tracker\_gpu.cu}   & device pointer + stream \\
materialize & \texttt{materialize\_\allowbreak gpu\_\allowbreak track\_\allowbreak results} & GPU→host,僅一次 copy \\
output      & \texttt{relink\_write}                               & host(fast MOT emit)\\
\bottomrule
\end{tabularx}
\caption{傳遞合約與 source。}\label{tab:transfer}
\end{table}

\section{Frame-Level Dataflow}\label{sec:dataflow}
對每個 frame \(f\),no-ReID baseline 的主線是兩條同源分支在 tracker 匯合:
\begin{center}
\small
\begin{tabular}{@{}rcl@{}}
frame tensor & \(\to\) & detect \(\to\) postprocess \(\to\) boxes \\
frame tensor & \(\to\) & GMC warp \(W_f\) \\
boxes \(+\ W_f\) & \(\to\) & tracker update \(\to\) materialize \\
materialize & \(\to\) & relink\_write / MOT rows
\end{tabular}
\end{center}

\section{Baseline 合約}\label{sec:baseline}
本文以 \texttt{mamba\_whole\_graph} preset 為基準。關鍵啟用條件:

\begin{table}[h]\small\setlength{\tabcolsep}{4pt}
\begin{tabularx}{\linewidth}{@{}l
  >{\raggedright\arraybackslash}p{4.8cm}
  >{\raggedright\arraybackslash}X@{}}
\toprule
區域 & baseline 值 & 效果 \\
\midrule
Graph     & \texttt{use\_whole\_graph: true}      & whole detect graph 單次 replay \\
GMC       & \texttt{gmc\_downscale: 4}            & GPU phase-correlation 估平移 \\
Appearance& \texttt{reid\_mode: off}              & headline 不用 appearance \\
Assoc     & \texttt{match\_thresh: 0.50}          & ByteTrack-like 多階段 gate \\
Cost      & \texttt{multiplicative\_cost: true}   & log-linear cost + stability reward \\
Relink    & \texttt{relink\_bridge\_enabled: true}& 雙向 foot bridge \\
\bottomrule
\end{tabularx}
\caption{Baseline 合約摘要。完整旋鈕表見附錄 A。}\label{tab:baseline}
\end{table}

% TODO: 把 math_model.md §3 的符號↔config 大表搬進 chapters/A-symbols.tex(附錄)


% ---- Part II:模組細節(一章一模組,全 pipeline 等深 L1–L2)----
% 每章 = \archmap{...} 迷你地圖 + contract 合約框 + 核心數學
%        + 末尾 \suppinput 掛 L3 實作補充(paper 版會自動剝掉)
\part{模組細節}
% ===========================================================================
% Detection 與 Postprocess(pipeline 最前段,故置於 Part II 第一章)
% ===========================================================================
\chapter{Detection 與 Postprocess}\label{ch:detect}

\begin{center}\archmap{detect}\end{center}

\begin{contract}{detect 鏈}
\textbf{輸入}\; frame \(640\times640\) CHW tensor\quad
\textbf{輸出}\; boxes / scores / classes(post-NMS)\\
\textbf{方法}\; YOLO26s TRT backbone+neck \(\to\) Mamba head(per-level lane)\(\to\)
anchor-free decode \(\to\) NMS\quad
\textbf{Source}\; \texttt{mamba\_gated\_detector.py}、\texttt{mamba\_head.py}、
\texttt{mamba\_gated\_detector.cpp}、\texttt{pipeline.cpp}\\
\textbf{Baseline}\; \texttt{tiling: native\_640}、\texttt{preprocess: none}、
ckpt \texttt{mamba\_gt\_v14replica\_t3\_t1};整段在 whole-detect CUDA graph 單次 replay
\end{contract}

\paragraph{演算法傳遞.} detect 鏈分兩段:detector(backbone+head,出 raw)與 native
postprocess(過濾/NMS)。邊上為各段傳遞的張量:

\begin{center}\resizebox{\linewidth}{!}{%
\begin{tikzpicture}[node distance=14mm]
  \node[step=Cdet] (n1) {frame\\$640^2$};
  \node[step=Cdet, right=of n1] (n2) {TRT backbone\\$+$ neck (FPN)};
  \node[step=Cdet, right=of n2] (n3) {Mamba head\\$\times(P_3,P_4,P_5)$};
  \node[step=Cdet, right=of n3] (n4) {decode\\anchor-free};
  \node[step=Cdet, right=of n4] (n5) {filter $+$ NMS};
  \draw[flowarr=Cdet] (n1) -- node[flowlbl,above]{image}              (n2);
  \draw[flowarr=Cdet] (n2) -- node[flowlbl,above]{$F_3,F_4,F_5$}       (n3);
  \draw[flowarr=Cdet] (n3) -- node[flowlbl,above]{$\mathrm{cls}_N,\mathrm{reg}_N$} (n4);
  \draw[flowarr=Cdet] (n4) -- node[flowlbl,above]{8400 anchors}        (n5);
\end{tikzpicture}}\end{center}

\section{整體結構與三尺度}\label{sec:det-overall}
P3/P4/P5 三條 lane 結構完全相同,只差尺度常數(輸入固定 640):

\begin{table}[h]\centering\small\renewcommand{\arraystretch}{1.2}
\begin{tabular}{@{}lcccc@{}}
\toprule
level \(N\) & stride \(s_N\) & \(H_N{=}W_N{=}640/s_N\) & channels \(C_N\) & tokens \(L_N{=}(H_N/4)^2\) \\
\midrule
P3 & 8  & 80 & 128 & 400 \\
P4 & 16 & 40 & 256 & 100 \\
P5 & 32 & 20 & 512 & 25  \\
\bottomrule
\end{tabular}
\caption{三尺度常數(\(s_N{=}2^N,\ C_N{=}2^{N+4}\))。}\label{tab:det-levels}
\end{table}

\section{單一 Lane:Mamba Head(Flow B)}\label{sec:det-mamba}
重點是 U-Net 式 skip:\(X_N\) 一邊進 down\(\to\)Mamba\(\to\)up,一邊\textbf{跳接}直接與
\(U_N\) concat,故 head 輸入是 256 ch。先 \texttt{Down4} 把序列壓到 \(L_N\le400\)
(控 scan 成本),上採樣後再補回細節:

\begin{center}\resizebox{\linewidth}{!}{%
\begin{tikzpicture}[node distance=11mm]
  \node[step=Cdet] (f)  {$F_N$\\$C_N{\times}H_N{\times}W_N$};
  \node[step=Cdet, right=of f]  (x)  {input\_proj\\$X_N$ (128)};
  \node[step=Cdet, right=of x]  (d)  {Down4\\stride-4};
  \node[step=Cdet, right=of d]  (z)  {flatten\\$Z_N$ ($L_N{\times}128$)};
  \node[step=Cdet, right=of z]  (m)  {MambaBlock\\$\times2$ ($d_{\mathrm{state}}{=}16$)};
  \node[step=Cdet, right=of m]  (u)  {Up4\\$U_N$ (128)};
  \node[step=Cdet, right=of u]  (c)  {concat\\$\to Y_N$ (256)};
  \draw[flowarr=Cdet] (f) -- (x);
  \draw[flowarr=Cdet] (x) -- (d);
  \draw[flowarr=Cdet] (d) -- node[flowlbl,above]{$\le400$ tok} (z);
  \draw[flowarr=Cdet] (z) -- (m);
  \draw[flowarr=Cdet] (m) -- (u);
  \draw[flowarr=Cdet] (u) -- (c);
  % U-Net skip:X_N 直接跳接到 concat(直角繞上方)
  \draw[flowarr=Cdet, dashed] (x.north) |- ++(0,7mm) -| node[flowlbl,above,pos=0.25]{skip $X_N$} (c.north);
\end{tikzpicture}}\end{center}

\section{Head 與 Decode(Flow C/D)}\label{sec:det-decode}
\textbf{Head}(per level,純卷積):\(Y_N(256)\) 分兩支,cls head
\(\to\mathrm{cls}_N\,(80)\)、reg head \(\to\mathrm{reg}_N\,(4)\);輸出通道
\(n_o=n_c+4\,r_{\max}=80+4=84\)(\(r_{\max}{=}1\),box 直接 4 維距離,無 DFL)。

\textbf{Decode}(anchor-free,跨三尺度合併,\(N=80^2{+}40^2{+}20^2=8400\) anchors)。
anchor 為 grid 中心 \((x{+}0.5,y{+}0.5)\),每點帶 stride:
\begin{equation}\label{eq:det-decode}
  \mathrm{reg}=(\ell t,rb)\ \xrightarrow{\text{dist2bbox}}\
  c=\mathrm{anchor}+\tfrac{rb-\ell t}{2},\quad
  wh=\ell t+rb\ \xrightarrow{\times s}\ \text{xyxy(像素)}.
\end{equation}
\begin{equation}\label{eq:det-cls}
  \text{score},\ \text{class}=\max_{80},\ \arg\max_{80}\ \sigma(\mathrm{cls}).
\end{equation}

\section{Postprocess Contract}\label{sec:det-post}
native postprocess(\Cref{ch:detect} 第二段)的核心 contract:
\begin{center}\ttfamily\small
(raw boxes/scores/classes) → score/class/geometry filter → compact/gather → NMS → output
\end{center}
baseline 旗標全關:
\begin{center}\small
\begin{tabular}{@{}ll@{\qquad}ll@{}}
\texttt{track\_person\_only} & false &
\texttt{person\_geometry\_prior} & false \\
\texttt{detection\_quality\_scaling} & false &
\texttt{geometry\_suspect\_support} & false
\end{tabular}
\end{center}
tracker 直接消費 post-NMS boxes,各 detection 進哪段 association 由 tracker
的 score gates(\(\tau_{\mathrm{track/mid/high}}\)、\texttt{new\_track\_thresh})決定。

\suppinput{supp/04-detection-impl}
    % pipeline 最前段
% ===========================================================================
% 細節章「樣板」:GMC。後續每章都照這個骨架寫,確保全 pipeline 等深。
%   1) \archmap{...}  迷你地圖,點亮本模組
%   2) contract 框    I/O / Source / Baseline,讓人不鑽公式也拿得到全貌
%   3) L1–L2 數學     學術本體(收 paper 用這層)
%   4) \suppinput     掛 L3 實作補充(paper 版自動剝掉)
% ===========================================================================
\chapter{GMC:相機運動補償}\label{ch:gmc}

\begin{center}\archmap{gmc}\end{center}

\begin{contract}{GMC}
\textbf{輸入}\; prev + curr 灰階(GPU tensor, CHW float \([0,1]\))\quad
\textbf{輸出}\; \(2\times3\) translation warp \(W_f\)\\
\textbf{方法}\; phase correlation(cross-power spectrum + Hanning window)\quad
\textbf{Source}\; \texttt{gmc\_kernel.cu}(cuFFT)\\
\textbf{Baseline}\; \texttt{gmc\_downscale: 4}、translation-only、soft confidence gate
\end{contract}

\noindent
GPU GMC 用 phase correlation 估 frame-to-frame 平移,作為 deterministic control
input 送進 \texttt{track} 的 Kalman predict,讓 velocity state 只學物體
residual motion 而非相機運動。

\paragraph{演算法傳遞.} 本章內部依序串接下列步驟,邊上標註各步之間傳遞的量
(對應 \Cref{sec:gmc-down,sec:gmc-cps,sec:gmc-warp}):

\gmcflow

\section{灰階降採樣與窗函數}\label{sec:gmc-down}
對 downscaled pixel \((x,y)\),floor-based 取樣:
\begin{equation}\label{eq:gmc-sample}
  s_x=\Big\lfloor x\tfrac{W_{\mathrm{src}}}{W_{\mathrm{dst}}}\Big\rfloor,\qquad
  s_y=\Big\lfloor y\tfrac{H_{\mathrm{src}}}{H_{\mathrm{dst}}}\Big\rfloor,\qquad
  G=0.299R+0.587G_c+0.114B.
\end{equation}
FFT 前套 Hanning window \(G_w(x,y)=G(x,y)\,w_x w_y\),其中
\(w_x=\tfrac12(1-\cos\tfrac{2\pi x}{W_{\mathrm{dst}}-1})\)(同理 \(w_y\))。

\section{Cross-Power Spectrum 與位移}\label{sec:gmc-cps}
令 \(A=\mathrm{FFT}(\text{prev})\)、\(B=\mathrm{FFT}(\text{curr})\):
\begin{equation}\label{eq:gmc-cps}
  C(k)=\overline{A(k)}B(k),\qquad
  R(k)=\frac{C(k)}{|C(k)|+10^{-6}},\qquad
  r=\mathrm{IFFT}(R).
\end{equation}
\(r\) 的 peak 給出 wrapped displacement(\(\text{peak}>\text{dim}/2\) 時減去該維長度)。

\section{Confidence 與 Warp}\label{sec:gmc-warp}
phase-correlation 可信度與軟縮放:
\begin{equation}\label{eq:gmc-pcr}
  \mathrm{PCR}=\frac{\max r}{\mathrm{RMS}(r)},\qquad
  \gamma=\begin{cases}\mathrm{clamp}(\mathrm{PCR}/\tau,0,1),&\mathrm{PCR}<\tau\\1,&\text{otherwise.}\end{cases}
\end{equation}
位移超過 downscaled 尺寸 25\% 視為不可信(退回 identity)。否則
\begin{equation}\label{eq:gmc-final}
  t_x=p_x d\,\gamma,\quad t_y=p_y d\,\gamma,\qquad
  W_f=\begin{bmatrix}1&0&t_x\\0&1&t_y\end{bmatrix}.
\end{equation}
\textbf{Sign convention}:positive \(t_x\) 把 predicted track center 往
current frame 右方推。

% ---------------------------------------------------------------------------
% L3 實作補充:kernel 名、env、四路徑 fallback 差異。paper 版自動剝掉。
% ---------------------------------------------------------------------------
\suppinput{supp/05-gmc-impl}
          % ← 完整範例章
% ===========================================================================
% track 主體 (1/3):Kalman 運動模型
% ===========================================================================
\chapter{Kalman 運動模型}\label{ch:kalman}

\begin{center}\archmap{track}\end{center}

\begin{contract}{Kalman(track 子機制)}
\textbf{輸入}\; 前一幀 track state \((x,P)\)、GMC warp \(W_f\)、配對到的 detection
measurement \(z\)\quad
\textbf{輸出}\; predict 後 \(\tilde{x}\)(供 gate/cost)、update 後 \((x,P)\)\\
\textbf{方法}\; SORT/DeepSORT 風格 constant-velocity filter,state
\((c_x,c_y,a,h,\dot{\cdot})\)\quad
\textbf{Source}\; \texttt{kalman\_gpu.cuh}\quad
\textbf{Baseline}\; \texttt{kalman\_r\_scale: 2.8}、NSA off
\end{contract}

\noindent
state 與 measurement:
\begin{equation}\label{eq:kf-state}
  x=(c_x,c_y,a,h,v_x,v_y,v_a,v_h)^{\!\top},\qquad
  z=(c_x,c_y,a,h)^{\!\top}.
\end{equation}

\paragraph{演算法傳遞.} 每幀對每個 track 依序:predict(先吃 GMC control)→ 算
innovation 協方差 → gate → 若入選則 update。邊上為各步傳遞的量:

\begin{center}\resizebox{\linewidth}{!}{%
\begin{tikzpicture}[node distance=14mm]
  \node[step=Ctrk] (n1) {GMC control\\$+$ predict\\$x'{=}Fx,\ P'{=}FPF^{\!\top}{+}Q$};
  \node[step=Ctrk, right=of n1] (n2) {innovation cov\\$S{=}MP'M^{\!\top}{+}R$};
  \node[step=Ctrk, right=of n2] (n3) {gate\\$\mathrm{IoU}\,\lor\,d^2_{\mathrm{maha}}{<}\tau$};
  \node[step=Ctrk, right=of n3] (n4) {gain\\$K{=}P'M^{\!\top}S^{-1}$};
  \node[step=Ctrk, right=of n4] (n5) {update\\$x{+}Ky,\ (I{-}KM)P'$};
  \draw[flowarr=Ctrk] (n1) -- node[flowlbl,above]{$\tilde{x},P'$} (n2);
  \draw[flowarr=Ctrk] (n2) -- node[flowlbl,above]{$S^{-1}$}       (n3);
  \draw[flowarr=Ctrk] (n3) -- node[flowlbl,above]{候選 $z_j$}     (n4);
  \draw[flowarr=Ctrk] (n4) -- node[flowlbl,above]{$y{=}z{-}Mx$}   (n5);
\end{tikzpicture}}\end{center}

\section{Prediction}\label{sec:kf-predict}
Transition 是 constant velocity,且 GMC translation 先以 control input 加到位置
(見 \Cref{ch:gmc}),再做 predict:
\begin{equation}\label{eq:kf-predict}
  F=\begin{bmatrix}I_4 & I_4\\ 0 & I_4\end{bmatrix},\qquad
  x'=Fx,\qquad P'=FPF^{\!\top}+Q(h^-).
\end{equation}
process noise 隨 predict 後框高 \(h^-\) 縮放:
\begin{equation}\label{eq:kf-Q}
  \sigma_p=\tfrac{h^-}{20},\quad \sigma_v=\tfrac{h^-}{160},\quad
  \mathrm{diag}(Q)=(\sigma_p^2,\sigma_p^2,10^{-4},\sigma_p^2,
                    \sigma_v^2,\sigma_v^2,10^{-10},\sigma_v^2).
\end{equation}

\section{Measurement Update}\label{sec:kf-update}
measurement matrix \(M=[\,I_4\ \ 0\,]\)。noise 用 predict 後 \(h^-\)(非 detection
height);baseline 下亮度/NSA 調節關閉(\(\lambda_{\mathrm{light}}{=}0,\ m_{\mathrm{NSA}}{=}1\)),
故 \(m_R=r_{\mathrm{scale}}\):
\begin{equation}\label{eq:kf-R}
  \mathrm{diag}(R)=\Big(\big(\tfrac{h^-}{20}\big)^2 m_R,\ \big(\tfrac{h^-}{20}\big)^2 m_R,\
                      10^{-2}m_R,\ \big(\tfrac{h^-}{20}\big)^2 m_R\Big).
\end{equation}
更新方程:
\begin{equation}\label{eq:kf-kalman}
\begin{aligned}
  S&=MP'M^{\!\top}+R, & K&=P'M^{\!\top}S^{-1}, & y&=z-Mx',\\
  x&\leftarrow x'+Ky, & P&\leftarrow(I-KM)P'.
\end{aligned}
\end{equation}

\section{Mahalanobis Gate}\label{sec:kf-gate}
即使 IoU 弱,只要 Mahalanobis 距離小於 \(\tau_{\mathrm{maha}}\) 仍可入選 candidate
(\(\tilde{x}_i\) 為已套 GMC control 並 predict 後的 state):
\begin{equation}\label{eq:kf-maha}
  d^2_{ij}=(z_j-M\tilde{x}_i)^{\!\top}S_i^{-1}(z_j-M\tilde{x}_i),\qquad
  \mathrm{cand}_{ij}\iff \mathrm{IoU}_{ij}>\tau_{\mathrm{iou}}\ \lor\ d^2_{ij}<\tau_{\mathrm{maha}}.
\end{equation}
IoU gate 與 cost(\Cref{ch:assoc})使用同一個 \(\tilde{x}_i\)。

\suppinput{supp/06-kalman-impl}
       % track 主體 1/3
% ===========================================================================
% track 主體 (2/3):Association Cost
% ===========================================================================
\chapter{Association Cost}\label{ch:assoc}

\begin{center}\archmap{track}\end{center}

\begin{contract}{Association cost(track 子機制)}
\textbf{輸入}\; predict 後 tracks \(\{\tilde{x}_i\}\)、detections \(\{b_j,s_j\}\)\quad
\textbf{輸出}\; 稀疏 candidate 上的成本 \(c_{ij}\)(供 auction)\\
\textbf{方法}\; IoU(+ReID)綜合質量 → 乘法式 cost,含 OAO/stability 調節\quad
\textbf{Source}\; \texttt{stage1\_cost\_fused\_kernel}(\texttt{tracker\_gpu.cu})\\
\textbf{Baseline}\; \texttt{multiplicative\_cost: true}、\texttt{sinkhorn\_lambda: 10}、
\texttt{stability\_cost\_w: 0.20}、\texttt{reid\_mode: off}
\end{contract}

\paragraph{演算法傳遞.} gate 過的 pair 先算 base quality,再匯入 penalty,經乘法式
轉成 cost,最後只把夠好的 enqueue 進稀疏 candidate list:

\begin{center}\resizebox{\linewidth}{!}{%
\begin{tikzpicture}[node distance=15mm]
  \node[step=Ctrk] (n1) {candidate gate\\$\mathrm{IoU}\,\lor\,\mathrm{Maha}$};
  \node[step=Ctrk, right=of n1] (n2) {base quality\\$A_{ij}{=}q^{\mathrm{iou}}_{ij}$};
  \node[step=Ctrk, right=of n2] (n3) {penalties\\$\Pi{=}P_{\mathrm{OAO}}{+}P_{\mathrm{vel}}{+}P_{\mathrm{occ}}{-}R_{\mathrm{stab}}$};
  \node[step=Ctrk, right=of n3] (n4) {mult.\ cost\\$c{=}1{-}A\,e^{-\Pi}$};
  \node[step=Ctrk, right=of n4] (n5) {sparse top-k\\enqueue};
  \draw[flowarr=Ctrk] (n1) -- node[flowlbl,above]{候選}        (n2);
  \draw[flowarr=Ctrk] (n2) -- node[flowlbl,above]{$A_{ij}$}     (n3);
  \draw[flowarr=Ctrk] (n3) -- node[flowlbl,above]{$\Pi_{ij}$}   (n4);
  \draw[flowarr=Ctrk] (n4) -- node[flowlbl,above]{$c_{ij}$}     (n5);
\end{tikzpicture}}\end{center}

\section{Candidate Gate 與 Base Quality}\label{sec:as-gate}
gate 同 \Cref{eq:kf-maha}。兩 gate 皆不過時 \(c_{ij}=1\)。detection score fusion:
\begin{equation}\label{eq:as-q}
  q^{\mathrm{iou}}_{ij}=\mathrm{IoU}_{ij}\bigl(1-w_{\mathrm{fuse}}(1-s_j)\bigr),
  \qquad A_{ij}=q^{\mathrm{iou}}_{ij}\ \ (\text{no ReID}).
\end{equation}
baseline \texttt{fuse\_score\_weight}=0,故 \(q^{\mathrm{iou}}_{ij}=\mathrm{IoU}_{ij}\)。
sparse list 只 enqueue 成本落在最寬鬆門檻內者:
\(c_{\max}=\max(c_{\mathrm{DDA}},\tau_{\mathrm{match}},\tau_{\mathrm{stage2}})\),
\(\mathrm{enqueue}_{ij}\iff c_{ij}\le c_{\max}\)。

\section{Multiplicative Cost}\label{sec:as-cost}
active cost form 與 penalty 匯總:
\begin{equation}\label{eq:as-cost}
  c_{ij}=\mathrm{clamp}\!\bigl(1-A_{ij}e^{-\Pi_{ij}},\,0,\,1\bigr),\qquad
  \Pi_{ij}=P_{\mathrm{OAO}}+P_{\mathrm{vel}}+P_{\mathrm{occ\_front}}-R_{\mathrm{stability}}.
\end{equation}
legacy additive path \(c_{ij}=1-A_{ij}+\sum_k P_k\) 仍存在但 baseline 不用。
baseline 只有 \(P_{\mathrm{OAO}}\) 與 \(R_{\mathrm{stability}}\) 起作用;
\(P_{\mathrm{vel}}\)(\texttt{vel\_dir\_weight})與 \(P_{\mathrm{occ\_front}}\)
(\texttt{occ\_state\_enabled})未啟用。

\section{OAO Penalty}\label{sec:as-oao}
依 predicted track-track overlap 計 occlusion 係數,並隨遮擋持續幀數 ramp:
\begin{equation}\label{eq:as-oao}
\begin{aligned}
  o^{\mathrm{base}}_i&=\max_{k\ne i}\mathrm{IoU}\!\bigl(B(x_i),B(x_k)\bigr),\qquad
  o_i=o^{\mathrm{base}}_i\min\!\Bigl(1,\tfrac{d_i}{N_{\mathrm{ramp}}}\Bigr),\\
  P_{\mathrm{OAO}}(i,j)&=\tau_{\mathrm{OAO}}\,o_i\,g_s(s_j).
\end{aligned}
\end{equation}
baseline \(\tau_{\mathrm{OAO}}=0.50\)、\(N_{\mathrm{ramp}}=25\);score 調節
\(g_s\equiv1\)(\texttt{oao\_score\_w} 未設)。

\section{Stability Reward}\label{sec:as-stab}
高度一致的 size-consistency reward(成本側,\(\lambda_{\mathrm{eff}}=\max(\lambda,1)\)):
\begin{equation}\label{eq:as-stab}
  R_{\mathrm{stability}}(i,j)=
  \frac{w_{\mathrm{stab}}/\lambda_{\mathrm{eff}}}
       {1+|h_i-h_j|/\max(h_j,10^{-3})}.
\end{equation}
因 reward 除以 \(\lambda\),進入 \(e^{-\lambda c}\)(\Cref{ch:auction})後其 boost 大致不隨
\(\lambda\) 失衡。baseline \texttt{stability\_cost\_w}=0.20。

\suppinput{supp/07-association-impl}
  % track 主體 2/3
% ===========================================================================
% track 主體 (3/3):Sparse Top-K 與 Auction Assignment
% ===========================================================================
\chapter{Sparse Top-K 與 Auction}\label{ch:auction}

\begin{center}\archmap{track}\end{center}

\begin{contract}{Auction(track 子機制)}
\textbf{輸入}\; 每個 track 的稀疏 cost candidates \(c_{ij}\)\quad
\textbf{輸出}\; 配對 \texttt{trk\_to\_det} / \texttt{det\_to\_trk}\\
\textbf{方法}\; ByteTrack 風格分數級聯 \(\times\) 單輪平行 Bertsekas auction\quad
\textbf{Source}\; \texttt{fused\_sinkhorn\_multistage\_kernel}、
\texttt{parallel\_auction\_shmem\_kernel}\\
\textbf{Baseline}\; 5-stage 級聯;\texttt{sinkhorn\_lambda} 僅作 \(e^{-\lambda c}\) value,
\emph{非}完整 Sinkhorn solve
\end{contract}

\paragraph{演算法傳遞.} cost candidates 先一次算出五個 stage 的 top-k,再依優先序
串行跑級聯;每個 stage reset 價格後跑單輪 auction 並 commit,未配掉的 track/det
流到下一 stage:

\begin{center}\resizebox{\linewidth}{!}{%
\begin{tikzpicture}[node distance=15mm]
  \node[step=Ctrk] (n1) {cost\\candidates};
  \node[step=Ctrk, right=of n1] (n2) {fused top-k\\(5 stages)};
  \node[step=Ctrk, right=of n2] (n3) {級聯\\S0$\to$S1$\to$S1b$\to$S1c$\to$S2};
  \node[step=Ctrk, right=of n3] (n4) {per-stage\\reset price $+$ auction};
  \node[step=Ctrk, right=of n4] (n5) {commit\\\texttt{trk\_to\_det}};
  \draw[flowarr=Ctrk] (n1) -- node[flowlbl,above]{$c_{ij}$} (n2);
  \draw[flowarr=Ctrk] (n2) -- node[flowlbl,above]{$p_{ij}$} (n3);
  \draw[flowarr=Ctrk] (n3) -- node[flowlbl,above]{逐 stage} (n4);
  \draw[flowarr=Ctrk] (n4) -- node[flowlbl,above]{winners}  (n5);
\end{tikzpicture}}\end{center}

\noindent
這\emph{不是} full dense Sinkhorn:每 stage 只跑一輪 bid+commit,price buffer 開頭
reset 為 0,沒有 price-raising 迭代到收斂——實質是單輪平行貪婪配對。

\section{Stage 級聯與 Value}\label{sec:au-stage}
五個 stage 依優先序貪婪,carry-over 只看前面沒配掉的 track/det:

\begin{table}[h]\centering\small\renewcommand{\arraystretch}{1.2}
\begin{tabular}{@{}llll@{}}
\toprule
Stage & Track state & Detection score & Cost cap \\
\midrule
S0 DDA       & confirmed  & \([\text{high},1.1)\)         & \(c_{\mathrm{DDA}}\) (0.12) \\
S1 high      & confirmed  & \([\text{high},1.1)\)         & \(\tau_{\mathrm{match}}\) \\
S1b mid      & confirmed  & \([\text{mid},\text{high})\)  & \(\tau_{\mathrm{match}}\) \\
S1c tentative& tentative  & \([\text{mid},1.1)\)          & \(\tau_{\mathrm{match}}\) \\
S2 low       & confirmed  & \([\text{track},\text{mid})\) & \(\tau_{\mathrm{stage2}}\) \\
\bottomrule
\end{tabular}
\caption{分數級聯(\texttt{run\_stage} 0..4)。}\label{tab:au-stage}
\end{table}

放入 top-k 的 value 與 aspect penalty:
\begin{equation}\label{eq:au-value}
  p_{ij}=e^{-\lambda c_{ij}}\,G_{\mathrm{aspect}}(b_j),\qquad
  r_j=\tfrac{\mathrm{width}_j}{\mathrm{height}_j},
\end{equation}
\begin{equation}\label{eq:au-aspect}
  G_{\mathrm{aspect}}=
  \begin{cases}
    \max(0.5,\,1-(r_j-0.8)), & r_j>0.8\\
    \max(0.5,\,1-5(0.15-r_j)), & r_j<0.15\\
    1, & \text{otherwise.}
  \end{cases}
\end{equation}

\section{Auction Bid}\label{sec:au-bid}
給定 detection 當前 price \(\rho_j\),track \(i\) 的 best/second 與 bid:
\begin{equation}\label{eq:au-bid}
  v_{ij}=p_{ij}-\rho_j,\quad
  j^\ast=\arg\max_j v_{ij},\quad
  \Delta\rho_i=v_{ij^\ast}-v^{(2)}_i+\epsilon,\quad
  \mathrm{bid}_i=\rho_{j^\ast}+\Delta\rho_i.
\end{equation}
single-round 下 \(\Delta\rho\) 只用來決定同一輪內多 track 競標同一 detection 誰勝出。
兩個 bid bias(絕對相加):
\begin{equation}\label{eq:au-bias}
  \mathrm{bid}_i \mathrel{+}=\frac{w_{\mathrm{fresh}}}{1+\mathrm{age}_i},\qquad
  \mathrm{bid}_i \mathrel{+}=\frac{w_{\mathrm{stab,bid}}}{1+|h_i-h_j|/h_j}.
\end{equation}
預設不一致,\textbf{勿一律當 off}:freshness \(w_{\mathrm{fresh}}\)
(\texttt{SACCADE\_FRESHNESS\_W})預設 0(關);stability bid bias
\(w_{\mathrm{stab,bid}}\)(\texttt{SACCADE\_STABILITY\_W})預設 \(0.1\)(\textbf{開})。

\suppinput{supp/08-auction-impl}
      % track 主體 3/3
% ===========================================================================
% track 收尾 (1/2):Track Lifecycle 與 Birth
% ===========================================================================
\chapter{Track Lifecycle 與 Birth}\label{ch:lifecycle}

\begin{center}\archmap{track}\end{center}

\begin{contract}{Lifecycle(track 子機制)}
\textbf{輸入}\; assignment 結果(matched / unmatched track / unmatched detection)\quad
\textbf{輸出}\; 更新後的 track 集合與狀態(tentative/confirmed/lost/removed)\\
\textbf{方法}\; hit-streak 累積 confirm、age 計數 lost/remove\quad
\textbf{Source}\; \texttt{tracker\_gpu.cu}\\
\textbf{Baseline}\; \texttt{new\_track\_thresh: 0.28}、\texttt{confirm\_streak: 3}、
\texttt{confirm\_score\_thresh: 0.50}、\texttt{track\_buffer: 30}、\texttt{per\_seq\_adapt: false}
\end{contract}

\paragraph{演算法傳遞.} assignment 後依配對狀態驅動 track 狀態機;下圖為主要轉移
(tentative 若連續未配會很快 drop,圖中略):

\begin{center}\resizebox{\linewidth}{!}{%
\begin{tikzpicture}[node distance=20mm]
  \node[step=Ctrk] (nw) {未配 det};
  \node[step=Ctrk, right=of nw] (te) {tentative};
  \node[step=Ctrk, right=of te] (co) {confirmed};
  \node[step=Ctrk, right=of co] (lo) {lost};
  \node[step=Ctrk, right=of lo] (rm) {removed};
  \draw[flowarr=Ctrk] (nw) -- node[flowlbl,above]{$s_j{\ge}$new\_thr} (te);
  \draw[flowarr=Ctrk] (te) -- node[flowlbl,above]{streak${\ge}3$}    (co);
  \draw[flowarr=Ctrk] (co) -- node[flowlbl,above]{unmatched}         (lo);
  \draw[flowarr=Ctrk] (lo) -- node[flowlbl,above]{age${>}$buffer}    (rm);
  \draw[flowarr=Ctrk] (lo) to[bend left=32] node[flowlbl,above]{re-match} (co);
\end{tikzpicture}}\end{center}

\section{State 轉移}\label{sec:lc-states}
assignment 後 tracker 更新 state:
\begin{itemize}\setlength\itemsep{2pt}
  \item \textbf{matched}(confirmed/tentative):以 detection \(z_j\) 做 Kalman update
        (\Cref{ch:kalman}),\texttt{hit\_streak}\(+1\),\texttt{age}/\texttt{time\_since\_update}
        歸零。
  \item \textbf{unmatched active track}:\texttt{age}\(\le\)\texttt{track\_buffer} 期間維持
        active/lost,超過 max age 後 remove。
  \item \textbf{unmatched detection}(\(s_j\ge\)\texttt{new\_track\_thresh}):生成新的
        tentative track。
\end{itemize}

\section{Birth 與 Confirm}\label{sec:lc-confirm}
新 tentative track 必須累積足夠連續 evidence 才會被視為 confirmed output:
\begin{equation}\label{eq:lc-confirm}
\begin{aligned}
  \text{confirmed}\iff{}&
  \texttt{hit\_streak}\ge\texttt{confirm\_streak}\\
  &{}\land\bar{s}\ge\texttt{confirm\_score\_thresh}.
\end{aligned}
\end{equation}
baseline:
\begin{center}\small
\begin{tabular}{@{}ll@{\qquad}ll@{}}
\texttt{confirm\_streak} & 3 &
\texttt{confirm\_score\_thresh} & 0.50 \\
\texttt{new\_track\_thresh} & 0.28 &
\texttt{track\_buffer} & 30
\end{tabular}
\end{center}
若干 birth-gate experiments 位於 \path{scripts/eval/config/lifecycle.py},但目前
preset 未啟用(\texttt{per\_seq\_adapt: false})。
普通新生 track 以 tentative 與 \texttt{hit\_streak=1} 起跑;bridge relink 復活的
slot 會保留 lost id 並直接回到 confirmed,不再走 tentative confirm gate。

\suppinput{supp/09-lifecycle-impl}
    % track 收尾 1/2
% ===========================================================================
% track 收尾 (2/2):Bridge Relink
% ===========================================================================
\chapter{Bridge Relink}\label{ch:relink}

\begin{center}\archmap{track}\end{center}

\begin{contract}{Bridge relink(track 子機制)}
\textbf{輸入}\; 新穩定的 candidate track、窗內未配的 lost confirmed track\quad
\textbf{輸出}\; candidate 採用 lost track 的 id(lost slot deactivate)\\
\textbf{方法}\; 雙向中點外推(bidirectional midpoint extrapolation),\emph{非} appearance ReID\quad
\textbf{Source}\; \texttt{tracker\_gpu.cu}(\(\ne\) \texttt{relink\_gate.cu} 的 appearance gate)\\
\textbf{Baseline}\; \texttt{relink\_bridge\_enabled: true}、\texttt{relink\_bridge\_px: 0.25}、
\texttt{margin: 0.05}、\texttt{h\_lo/h\_hi: 0.75/1.33}、\texttt{dir\_bonus: 0.8}、
anchor \texttt{adaptive}
\end{contract}

\paragraph{演算法傳遞.} 對 candidate–lost pair,先取 anchor、做 4 點速度回歸,再雙向
外推算殘差,合成 \(d_{\mathrm{bridge}}\),經 direction bonus 與多重 gate 後 commit:

\begin{center}\resizebox{\linewidth}{!}{%
\begin{tikzpicture}[node distance=12mm]
  \node[step=Ctrk] (n1) {candidate\\$+$ lost pair};
  \node[step=Ctrk, right=of n1] (n2) {anchor\\$(a_x,a_y)$};
  \node[step=Ctrk, right=of n2] (n3) {4-pt velocity\\$v_\ell,v_c$};
  \node[step=Ctrk, right=of n3] (n4) {外推殘差\\$r_{\mathrm{fwd}},r_{\mathrm{bwd}}$};
  \node[step=Ctrk, right=of n4] (n5) {$d_{\mathrm{bridge}}$\\$+$ dir bonus};
  \node[step=Ctrk, right=of n5] (n6) {gates\\$+$ commit};
  \draw[flowarr=Ctrk] (n1) -- (n2);
  \draw[flowarr=Ctrk] (n2) -- node[flowlbl,above]{$\ell,c$} (n3);
  \draw[flowarr=Ctrk] (n3) -- (n4);
  \draw[flowarr=Ctrk] (n4) -- node[flowlbl,above]{$d_h$} (n5);
  \draw[flowarr=Ctrk] (n5) -- node[flowlbl,above]{$d^{(1)}$} (n6);
\end{tikzpicture}}\end{center}

\section{History 與 Anchor}\label{sec:rl-anchor}
每個 observed slot(\texttt{age}=0)保存 center/height history ring,並對框高做 EMA:
\(\bar{h}\leftarrow0.95\bar{h}+0.05h\)。coasting lost track 保留最後 history。
baseline anchor \texttt{adaptive}:\(x\) 永遠用 center;\(y\) 候選為
\(y^{\mathrm{top}}{=}c_y{-}h/2,\ y^{\mathrm{bot}}{=}c_y{+}h/2,\ y^{\mathrm{ctr}}{=}c_y\)。
若平均相鄰 height 變化 \(\le\)\texttt{anchor\_rate}(0.03)退化為 center;否則以 top/bottom
各自四點線性回歸殘差為權重選較穩的 edge:
\begin{equation}\label{eq:rl-anchor}
  w_e=\frac{1}{\mathrm{RSS}_{\mathrm{line}}(y^e)/(\bar{h}^2+10^{-3})+0.01},\quad
  a_y=\frac{w_{\mathrm{top}}y^{\mathrm{top}}+w_{\mathrm{bot}}y^{\mathrm{bot}}}
           {w_{\mathrm{top}}+w_{\mathrm{bot}}}.
\end{equation}

\section{Candidate 條件}\label{sec:rl-cand}
candidate 採用 lost id 的前置條件:
\begin{center}\ttfamily\small
candidate hit\_streak == bridge\_at\quad candidate 有 $\ge$4 history samples\\
lost active 但本幀未配\quad lost state == confirmed\quad
bridge\_min\_lost $\le$ lost\_age $\le$ bridge\_ttl
\end{center}

\section{速度回歸與 \texorpdfstring{\(d_{\mathrm{bridge}}\)}{d\_bridge}}\label{sec:rl-bridge}
對四個等時距 scalar anchor sample \((p_0,p_1,p_2,p_3)\),4 點回歸速度
\(v=\tfrac{3p_3+p_2-p_1-3p_0}{10}\)。令 \(\ell\)=lost exit anchor、\(c\)=candidate
entry anchor、\(g\)=\texttt{lost\_age}、\(h_{\mathrm{ref}}=\max(\tfrac{\bar{h}_\ell+\bar{h}_c}{2},1)\):
\begin{equation}\label{eq:rl-resid}
  r_{\mathrm{fwd}}=\frac{\lVert(\ell+v_\ell g)-c\rVert}{h_{\mathrm{ref}}},\qquad
  r_{\mathrm{bwd}}=\frac{\lVert(c-v_c g)-\ell\rVert}{h_{\mathrm{ref}}}.
\end{equation}
以 lost 端速度大小決定外推 vs 純距離的權重:
\begin{align}\label{eq:rl-dbridge}
  d_h&=\frac{\lVert\ell-c\rVert}{h_{\mathrm{ref}}},\\
  w&=\sqrt{\mathrm{clamp}\!\Bigl(\tfrac{\lVert v_\ell\rVert/h_{\mathrm{ref}}}{0.12},0,1\Bigr)},\\
  d_{\mathrm{bridge}}&=
  w\,\frac{r_{\mathrm{fwd}}+r_{\mathrm{bwd}}}{2}+(1-w)\,d_h.
\end{align}

\section{Direction Bonus}\label{sec:rl-dir}
\texttt{dir\_bonus}>0 且方向相近(\(\cos\theta>0.5\))時,向 cross-track 誤差 \(d_{\mathrm{cross}}\) blend:
\begin{equation}\label{eq:rl-dir}
  \alpha=\min\!\bigl(w_{\mathrm{dir}}\cos^2\theta\,\eta_v\,\eta_g,\,1\bigr),\qquad
  d_{\mathrm{bridge}}\leftarrow(1-\alpha)d_{\mathrm{bridge}}+\alpha\,d_{\mathrm{cross}},
\end{equation}
其中 \(\eta_v=\mathrm{clamp}\bigl(\tfrac{\min(\lVert v_\ell\rVert,\lVert v_c\rVert)}
{\max(0.005h_{\mathrm{ref}},10^{-3})},0,1\bigr)\)、\(\eta_g=\mathrm{clamp}(g/30,0,1)\)。
baseline \texttt{dir\_bonus}=0.8。

\section{Gates 與 Commit}\label{sec:rl-gate}
接受條件:
\begin{equation}\label{eq:rl-gate}
  d_{\mathrm{bridge}}\le\tau_{\mathrm{bridge}},\quad
  \frac{\bar{h}_\ell}{\bar{h}_c}\in[h_{\mathrm{lo}},h_{\mathrm{hi}}],\quad
  d^{(2)}_{\mathrm{bridge}}-d^{(1)}_{\mathrm{bridge}}\ge m_{\mathrm{bridge}}.
\end{equation}
注意 \texttt{relink\_bridge\_px}=0.25 是歷史命名,\emph{非} 0.25 pixel——它與已除以
\(h_{\mathrm{ref}}\) 的 \(d_{\mathrm{bridge}}\) 比較,單位是 reference-height-normalized
distance。baseline \texttt{spatial\_gate}=0.0(關),不啟用 occupancy expansion。
接受後 candidate 採用 lost id、lost slot deactivate;多 candidate 競爭同一 lost id 時,
較高 detection score 贏,candidate index 作 tie-breaker。

\suppinput{supp/10-relink-impl}
       % track 收尾 2/2

% ---- 附錄:reference 性質的大表(符號↔config)放這,不擋正文節奏 ----
\appendix
% ===========================================================================
% 附錄 A:符號表 與 符號↔config 對照(math_model §3 / §3.1 / §3.2)
% reference 性質,故置於附錄。代碼錨點(cu:N 行號)見各章 L3 實作補充(full 版)。
% ===========================================================================
\chapter{符號與 Config 對照}\label{app:symbols}

baseline 值對應 \texttt{mamba\_whole\_graph}(frozen\_v2)。各機制的代碼錨點
(kernel 名、\texttt{cu:N} 行號)見對應章末的 L3 實作補充。

\section{基本符號}\label{app:basic}

\begin{longtable}{@{}>{\raggedright\arraybackslash}p{3.3cm}
  >{\raggedright\arraybackslash}p{8.2cm}@{}}
\toprule
符號 & 意義 \\* \midrule \endhead
\(f\) & 目前 frame index \\
\(I_f\) & 目前 RGB/CHW frame tensor \\
\(D_f=\{d_j\}\) & frame \(f\) postprocess 後的 detection set \\
\(b_j=(x_1,y_1,x_2,y_2)\) & detection box(original-frame pixels) \\
\(s_j,\ \mathrm{cls}_j\) & detection score / class id \\
\(T_f=\{t_i\}\) & association 前 active 或 lost tracker slots \\
\(x_i,\ P_i\) & track slot \(i\) 的 8D Kalman state / \(8\times8\) covariance \\
\(W_f\) & GMC \(2\times3\) camera warp;translation path \([1,0,t_x;0,1,t_y]\) \\
\midrule
\(x=(c_x,c_y,a,h,\dot{\cdot})\) & center、aspect、height 與對應速度 \\
\(w=a\cdot h,\ B(x)\) & state implied width / 還原的 box \\
\(z=(c_x,c_y,a,h),\ M=[I_4\ 0]\) & Kalman measurement / measurement matrix \\
\(\mathrm{IoU}(i,j),\ c_{ij}\) & \(B(x_i)\) 與 \(b_j\) 的 IoU / final association cost \\
\(p_{ij},\ A_{ij},\ \Pi_{ij}\) & auction value / base quality / penalty 總和 \\
\bottomrule
\end{longtable}

% 共用:4 欄 config 對照表的欄型(footnotesize,長名稱用 \newline 斷行)
\section{GMC}\label{app:gmc}
{\footnotesize\setlength{\tabcolsep}{3pt}
\begin{longtable}{@{}>{\raggedright\arraybackslash}p{2.8cm}
                     >{\raggedright\arraybackslash}p{2.4cm}
                     >{\raggedright\arraybackslash}p{3.8cm}
                     >{\raggedright\arraybackslash}p{1.6cm}@{}}
\toprule
符號 & 用途 & config / env & baseline \\* \midrule \endhead
\(W_f\) & 相機運動補償 warp & — & trans-only \\
PCR & phase-corr 峰值可信度 & — & — \\
\(\tau_{\mathrm{PCR}}\) & 低可信度縮小位移 & \texttt{SACCADE\_GMC\_PCR\_THRESH} & py \(5.0\) \\
\bottomrule
\end{longtable}}

\section{Kalman}\label{app:kalman}
{\footnotesize\setlength{\tabcolsep}{4pt}
\begin{longtable}{@{}>{\raggedright\arraybackslash}p{1.8cm}
                     >{\raggedright\arraybackslash}p{3.0cm}
                     >{\raggedright\arraybackslash}p{4.2cm}
                     >{\raggedright\arraybackslash}p{1.7cm}@{}}
\toprule
符號 & 用途 & config / env & baseline \\* \midrule \endhead
\(x,z,F,M\) & 8D CV state / 4D measurement & — & — \\
\(\sigma_p{=}h^-/20\) & 位置過程噪聲 & hardcoded & 1/20 \\
\(\sigma_v{=}h^-/160\) & 速度過程噪聲 & hardcoded & 1/160 \\
\(r_{\mathrm{scale}}\) & 測量噪聲縮放 & \texttt{kalman\_r\_scale} & 2.8 \\
\(m_{\mathrm{NSA}},\lambda_{\mathrm{light}}\) & NSA/亮度調節 & — & 1 / 0 \\
\(\tau_{\mathrm{maha}}\) & Mahalanobis gate & \texttt{maha\_gate} & 見實作 \\
\bottomrule
\end{longtable}}

\section{Association Cost}\label{app:assoc}
{\footnotesize\setlength{\tabcolsep}{4pt}
\begin{longtable}{@{}>{\raggedright\arraybackslash}p{1.8cm}
                     >{\raggedright\arraybackslash}p{3.0cm}
                     >{\raggedright\arraybackslash}p{4.2cm}
                     >{\raggedright\arraybackslash}p{1.7cm}@{}}
\toprule
符號 & 用途 & config / env & baseline \\* \midrule \endhead
\(A_{ij}\) & IoU(+ReID)綜合質量 & — & = IoU \\
\(w_{\mathrm{fuse}}\) & 低分檢測降權 & \texttt{fuse\_score\_weight} & 0.0 \\
\(c_{ij}\) & 乘法式 cost & \texttt{multiplicative\_cost} & true \\
\(\lambda\) & cost\(\to\)value 溫度 & \texttt{sinkhorn\_lambda} & 10 \\
\(o_i,\tau_{\mathrm{OAO}}\) & 遮擋配對抑制 & \texttt{oao\_tau} /\newline \texttt{oao\_ramp\_frames} & 0.50 / 25 \\
\(w_{\mathrm{vel}}\) & 速度反向懲罰 & \texttt{vel\_dir\_weight} & 關 \\
\(w_{\mathrm{occ}}\) & front-occluder 懲罰 & \texttt{occ\_cost\_weight} & 關 \\
\(w_{\mathrm{stab}}\) & 高度一致 reward(成本側) & \texttt{stability\_cost\_w} & 0.20 \\
\bottomrule
\end{longtable}}

\section{Auction}\label{app:auction}
{\footnotesize\setlength{\tabcolsep}{4pt}
\begin{longtable}{@{}>{\raggedright\arraybackslash}p{1.8cm}
                     >{\raggedright\arraybackslash}p{3.0cm}
                     >{\raggedright\arraybackslash}p{4.2cm}
                     >{\raggedright\arraybackslash}p{1.7cm}@{}}
\toprule
符號 & 用途 & config / env & baseline \\* \midrule \endhead
\(p_{ij}\) & auction 概率值 & — & — \\
\(G_{\mathrm{aspect}}\) & 抑制異常長寬比 & hardcoded (0.8/0.15) & — \\
\(w_{\mathrm{fresh}}\) & 新鮮度 bid bias & \texttt{SACCADE\_FRESHNESS\_W} & 0.0 \\
\(w_{\mathrm{stab,bid}}\) & 高度一致 bid bias & \texttt{SACCADE\_STABILITY\_W} & 0.1(開) \\
S0 DDA & confirmed×high 更緊 stage & \texttt{SACCADE\_ENABLE\_DDA} /\newline \texttt{\_DDA\_MAX\_COST} & on / 0.12 \\
stage thr & 分數級聯邊界 & \texttt{match} / \texttt{high} / \texttt{mid} /\newline \texttt{track} / \texttt{stage2} & .50/.45/\newline.10/.05/.50 \\
\bottomrule
\end{longtable}}

\section{Lifecycle / Bridge / Semantic relink}\label{app:rest}
{\footnotesize\setlength{\tabcolsep}{4pt}
\begin{longtable}{@{}>{\raggedright\arraybackslash}p{1.8cm}
                     >{\raggedright\arraybackslash}p{3.0cm}
                     >{\raggedright\arraybackslash}p{4.2cm}
                     >{\raggedright\arraybackslash}p{1.7cm}@{}}
\toprule
符號 & 用途 & config / env & baseline \\* \midrule \endhead
birth/\newline confirm & tentative\(\to\)confirmed & \texttt{new\_track\_thresh} /\newline \texttt{confirm\_streak} & 0.28 / 3 \\
\(\tau_{\mathrm{bridge}}\) & 雙向外推殘差門檻 & \texttt{relink\_bridge\_px} & 0.25 \\
\(\alpha\) & 方向一致偏移 & \texttt{relink\_bridge\_dir\_bonus} & 0.8 \\
\(h_{\mathrm{lo}},h_{\mathrm{hi}}\) & 高度比 gate & \texttt{relink\_bridge\_h\_lo} /\newline \texttt{\_h\_hi} & 0.75/1.33 \\
\(m_{\mathrm{bridge}}\) & best-vs-second margin & \texttt{relink\_bridge\_margin} & 0.05 \\
\(w_{\mathrm{sim/iou}}\),\newline \(w_{\mathrm{maha}}\) & semantic joint 權重 & \texttt{semantic\_w\_*\_base} & off \\
\bottomrule
\end{longtable}}

\section{方法出處 / 命名對照}\label{app:naming}
\begin{itemize}\setlength\itemsep{2pt}
  \item \textbf{GMC} = phase correlation(cross-power spectrum + Hanning window),translation-only warp。
  \item \textbf{Kalman} = SORT/DeepSORT 風格 constant-velocity filter。
  \item \textbf{Assignment} = Bertsekas auction(單輪平行貪婪)跑在 softmin-temperature top-k 上;
        分數分段 = ByteTrack 風格 cascade。\texttt{sinkhorn\_lambda} 為歷史命名——只用
        \(e^{-\lambda c}\) 當 value,\textbf{非}完整 Sinkhorn 迭代。
  \item \textbf{OAO} = occlusion-aware(track-track overlap)配對抑制 + duration ramp。
  \item \textbf{Bridge relink} = 雙向中點外推,項目自有機制,非標準 appearance ReID。
  \item \textbf{Semantic relink gate} = appearance + Mahalanobis + IoU joint gate(baseline 關)。
\end{itemize}

\paragraph{提醒.} \(w_{\mathrm{stab}}\)(\Cref{sec:as-stab},成本側 reward)與
\(w_{\mathrm{stab,bid}}\)(\Cref{sec:au-bid},auction bid bias)是\textbf{兩個不同旋鈕},
數值與作用點皆不同,雖都用高度一致性 \(|h_i-h_j|/h_j\)。


\end{document}
